% Acides faibles, bases faibles et couples acide/base — Chimie Tle C — Cours signé OnBuch+
% Programme officiel MINESEC, Terminale C, D, E. Compiler avec : tectonic L08-acides-faibles-bases-faibles-couples-ka.tex
\newcommand{\DOCMATIERE}{Chimie}
\newcommand{\DOCNIVEAU}{Tle C}
\newcommand{\DOCMODULE}{Module 2 · Acides et bases}
\newcommand{\DOCLECON}{08}
\newcommand{\DOCTITRE}{Acides faibles, bases faibles et couples acide/base}
% OnBuch+ — préambule « cours signés » ENRICHI (compilé avec tectonic / XeLaTeX)
% Système visuel « Pop » d'OnBuch+ : fond crème (jamais blanc), bordures encre
% épaisses, ombres « dures » décalées, titres Archivo Black en capitales.
% Version enrichie pour les cours de chimie : chimie (mhchem, chemfig),
% graphes (pgfplots), boîtes pédagogiques supplémentaires (définition,
% propriété, méthode, attention, le savais-tu, exercice résolu, exercice,
% TP, bilan), page de garde refaite et signature OnBuch+ ancrée en bas.
%
% Macros attendues dans main.tex AVANT \input{preamble} :
%   \DOCMATIERE  \DOCNIVEAU  \DOCMODULE  \DOCLECON (numéro)  \DOCTITRE
\documentclass[11pt]{article}

\usepackage{fontspec}
\usepackage[french]{babel}
\usepackage[a4paper,margin=2cm,top=3.4cm,bottom=2.8cm,headheight=1.4cm,footskip=1.3cm]{geometry}
\usepackage{amsmath,amssymb,amsfonts}
\usepackage[table]{xcolor} % [table] : fournit aussi \rowcolors (alternance de lignes)
\usepackage{pagecolor}
\usepackage{tikz}
\usetikzlibrary{arrows.meta,positioning,calc,shapes.geometric,shapes.misc,shapes.arrows,decorations.pathmorphing,decorations.markings,patterns,shadows,mindmap,trees,fit,backgrounds,matrix,intersections}
\usepackage{pgfplots}
\pgfplotsset{compat=1.18}
\usepackage[version=4]{mhchem}
\usepackage{chemfig}
\usepackage{enumitem}
\usepackage{fancyhdr}
% [most] charge toutes les bibliothèques tcolorbox (listings, minted, raster,
% documentation...) et gonfle inutilement la mémoire TeX sur un document long
% (~300 boîtes) : on ne charge que ce qui est réellement utilisé.
\usepackage[skins,breakable]{tcolorbox}
\usepackage{graphicx}
\usepackage{colortbl}
\usepackage{booktabs}
\usepackage{tabularx}
\usepackage{array}
\usepackage{multicol}
\usepackage{siunitx}
% parse-numbers=false : accepte aussi \SI{2,5 \times 10^{-3}}{...}
\sisetup{locale=FR,output-decimal-marker={,},group-minimum-digits=4,parse-numbers=false}
\usepackage{qrcode}
% \degree : commande fréquemment utilisée par les modèles pour un angle ou une
% température en dehors de \SI{}{} (non fournie par les paquets chargés).
\providecommand{\degree}{^{\circ}}
\usepackage{lastpage}
\usepackage{hyperref}
\hypersetup{hidelinks}

% ============================================================
% Palette « Pop » OnBuch+ — mêmes hex que la classe P (plus_ui.dart)
% ============================================================
\definecolor{poporange}{HTML}{F59321}
\definecolor{popdark}{HTML}{E07A0C}
\definecolor{popcream}{HTML}{FFF6E8}
\definecolor{popink}{HTML}{1C1714}
\definecolor{poppurple}{HTML}{7A5AE0}
\definecolor{popgreen}{HTML}{2E9E6A}
\definecolor{poppink}{HTML}{E0457B}
\definecolor{popblue}{HTML}{2F7DE1}
\definecolor{popgold}{HTML}{C98A12}
\definecolor{popmuted}{HTML}{8A7F76}
\definecolor{popline}{HTML}{EADFD2}
\definecolor{popgray}{HTML}{8A7F76}
\colorlet{popgrey}{popgray}
% Alias tolérants (le modèle nomme parfois les couleurs différemment)
\colorlet{popwhite}{white}
\colorlet{popblack}{popink}
\colorlet{popred}{poppink}
\colorlet{popyellow}{popgold}
% Teintes claires pour fonds de tableaux / zones
\colorlet{popblueL}{popblue!10!white}
\colorlet{poporangeL}{poporange!14!white}
\colorlet{popgreenL}{popgreen!12!white}
\colorlet{poppurpleL}{poppurple!12!white}
\colorlet{poppinkL}{poppink!12!white}
\colorlet{popgoldL}{popgold!14!white}

\pagecolor{popcream}

% ============================================================
% Typographie
% ============================================================
\defaultfontfeatures{Ligatures=TeX}
\setmainfont{PlusJakartaSans-Medium.ttf}[Path=../../../fonts/,BoldFont=PlusJakartaSans-Bold.ttf]
% Maths en sans-serif (Fira Math) avec chiffres et lettres droites de Plus
% Jakarta, pour que les formules se fondent dans le texte.
\usepackage[math-style=ISO,bold-style=ISO]{unicode-math}
\setmathfont{FiraMath-Regular.otf}[Path=../../../fonts/]
\setmathfont{PlusJakartaSans-Medium.ttf}[Path=../../../fonts/,range={up/num,up/{Latin,latin},\mathup}]
\usepackage{icomma} % virgule décimale sans espace en mode math
% Caractères absents de la police de texte (sinon carré vide) : rendus par
% leur équivalent mathématique, en texte comme en formule.
\usepackage{newunicodechar}
\newunicodechar{α}{\ensuremath{\alpha}}
\newunicodechar{Α}{\ensuremath{\alpha}}
\newunicodechar{β}{\ensuremath{\beta}}
\newunicodechar{δ}{\ensuremath{\delta}}
\newunicodechar{✓}{\popcheck{popgreen}}
\newfontfamily\popdisplay{ArchivoBlack-Regular.ttf}[Path=../../../fonts/]
\newfontfamily\poplogo{SpaceGrotesk-Bold.ttf}[Path=../../../fonts/]
\newfontfamily\popsemi{PlusJakartaSans-SemiBold.ttf}[Path=../../../fonts/]
\newfontfamily\popxbold{PlusJakartaSans-ExtraBold.ttf}[Path=../../../fonts/]
\newfontfamily\popxboldls{PlusJakartaSans-ExtraBold.ttf}[Path=../../../fonts/,LetterSpace=3.0]

\setlist[enumerate]{leftmargin=1.7em,itemsep=5pt,topsep=4pt}
\setlist[itemize]{leftmargin=1.7em,itemsep=4pt,topsep=4pt,label={\color{poporange}\textbullet}}
\setlength{\parindent}{0pt}
\setlength{\parskip}{5pt}
\renewcommand{\arraystretch}{1.25}

% chemfig : liaisons un peu plus épaisses, encre Pop
\setchemfig{atom sep=2em,bond style={line width=0.9pt,popink},cram width=3pt}

% pgfplots : style Pop par défaut
\pgfplotsset{
  every axis/.append style={axis line style={popink,line width=1pt},tick style={popink},
    grid style={popline,line width=0.5pt},label style={font=\small},tick label style={font=\footnotesize}},
  popcurve/.style={poporange,line width=1.8pt,no markers,smooth},
}

% ============================================================
% Pill / squiggle
% ============================================================
\newcommand{\poppill}[3]{%
  \tikz[baseline=-0.55ex]\node[draw=popink,line width=1.2pt,rounded corners=2.6pt,%
    fill=#2,inner xsep=6.5pt,inner ysep=3pt]{%
    \popxboldls\fontsize{7}{7}\selectfont\color{#3}\MakeUppercase{#1}};%
}
\newcommand{\popsquiggle}[1]{%
  \tikz[baseline=-0.4ex]{
    \draw[#1,line width=2.6pt,line cap=round]
      plot [smooth] coordinates {(0,0.09) (0.34,0.01) (0.68,0.21) (1.02,0.01) (1.36,0.10)};
  }%
}
% Mot-clé surligné (marqueur orange sous le texte)
\newcommand{\cle}[1]{{\popxbold\color{popdark}#1}}

% ============================================================
% En-tête / pied
% ============================================================
\pagestyle{fancy}
\fancyhf{}
\renewcommand{\headrulewidth}{0pt}
\renewcommand{\footrulewidth}{0pt}
\lhead{\raisebox{-0.15ex}{\poplogo\fontsize{13}{13}\selectfont\color{popink}OnBuch\color{poporange}+}}
\rhead{\poppill{\DOCMATIERE\ \textbullet\ \DOCNIVEAU\ \textbullet\ Leçon \DOCLECON}{white}{popink}}
\lfoot{\popxbold\fontsize{7.6}{7.6}\selectfont\color{popmuted}\MakeUppercase{Cours signé OnBuch+}\ \textbullet\ {\popsemi\DOCTITRE}}
\rfoot{\tikz[baseline=-0.6ex]\node[rounded corners=8.5pt,fill=popink,minimum height=17pt,minimum width=17pt,inner xsep=5pt,inner sep=0]{\popxbold\fontsize{8}{8}\selectfont\color{popcream}\thepage\,/\,\pageref*{LastPage}};}

% ============================================================
% Titres
% ============================================================
\newcommand{\chaptitre}[2]{%
  \begin{tcolorbox}[enhanced,colback=poporange,colframe=popink,boxrule=2.6pt,arc=11pt,
      drop shadow=popink,left=15pt,right=15pt,top=12pt,bottom=11pt,before skip=3pt,after skip=18pt]
    \poppill{#2}{popink}{popcream}\par\vspace{9pt}
    {\raggedright\hyphenpenalty=10000\exhyphenpenalty=10000\popdisplay\fontsize{22}{24}\selectfont\color{popcream}\MakeUppercase{#1}\par}
  \end{tcolorbox}
}
% Section numérotée : pastille orange avec le numéro + titre Archivo
\newcounter{popsec}
\newcommand{\coursec}[1]{%
  \stepcounter{popsec}\setcounter{popsub}{0}%
  \par\vspace{16pt}\noindent
  \begin{minipage}{\linewidth}
  \tikz[baseline=-0.7ex]\node[draw=popink,line width=1.6pt,fill=poporange,rounded corners=4pt,minimum size=22pt,inner sep=0]{\popdisplay\fontsize{12}{12}\selectfont\color{popcream}\thepopsec};%
  \hspace{8pt}{\popdisplay\fontsize{15}{17}\selectfont\color{popink}\MakeUppercase{#1}}\par
  \vspace{4pt}\noindent{\color{poporange}\rule{36pt}{3.6pt}}
  \end{minipage}\par\nopagebreak\vspace{8pt}
}
\newcounter{popsub}
\newcommand{\courssub}[1]{%
  \stepcounter{popsub}%
  \par\vspace{9pt}\noindent{\popxbold\fontsize{11.5}{13}\selectfont\color{popink}{\color{poporange}\thepopsec.\thepopsub}\hspace{6pt}#1}\par\nopagebreak\vspace{4pt}
}

% ============================================================
% Boîtes pédagogiques — même recette : fond blanc, bordure épaisse colorée,
% ombre dure encre, badge plein. Titre optionnel [#1].
% ============================================================
\tcbset{popbase/.style={breakable,enhanced,colback=white,boxrule=2.3pt,arc=9pt,
  shadow={3pt}{-3pt}{0pt}{popink},left=14pt,right=14pt,top=11pt,bottom=11pt,before skip=11pt,after skip=15pt,
  fontupper=\popsemi\fontsize{10.6}{14.5}\selectfont\color{popink},
  fontlower=\popsemi\fontsize{10.6}{14.5}\selectfont\color{popink}}}
% Coche dessinée (le glyphe ✓ n'existe pas dans les polices du document)
\newcommand{\popcheck}[1]{\tikz[baseline=-0.5ex]\draw[#1,line width=1.6pt,line cap=round,line join=round](-0.13,0.0)--(-0.04,-0.09)--(0.13,0.1);}
\newcommand{\popboxhead}[3]{\poppill{#1}{#2}{white}\ifx\relax#3\relax\else\hspace{6pt}{\popxbold\fontsize{10.6}{12}\selectfont\color{popink}#3}\fi\par\vspace{7pt}}

\newtcolorbox{aretenir}[1][]{popbase,colframe=popgreen,before upper={\popboxhead{À retenir}{popgreen}{#1}}}
\newtcolorbox{exemplebox}[1][]{popbase,colframe=poporange,before upper={\popboxhead{Exemple}{poporange}{#1}}}
\newtcolorbox{definition}[1][]{popbase,colframe=popblue,before upper={\popboxhead{Définition}{popblue}{#1}}}
\newtcolorbox{propriete}[1][]{popbase,colframe=poppurple,before upper={\popboxhead{Propriété}{poppurple}{#1}}}
\newtcolorbox{methode}[1][]{popbase,colframe=popgold,before upper={\popboxhead{Méthode}{popgold}{#1}}}
\newtcolorbox{attention}[1][]{popbase,colframe=poppink,before upper={\popboxhead{Attention}{poppink}{#1}}}
\newtcolorbox{experience}[1][]{popbase,colframe=popblue,colback=popblueL,before upper={\popboxhead{Expérience}{popblue}{#1}}}
% « Le savais-tu ? » : carte violette pleine (contexte, culture, Cameroun)
\newtcolorbox{savaistu}[1][]{popbase,colback=poppurple,colframe=popink,
  fontupper=\popsemi\fontsize{10.4}{14.2}\selectfont\color{white},
  before upper={\poppill{Le savais-tu\,?}{white}{poppurple}\ifx\relax#1\relax\else\hspace{6pt}{\popxbold\color{white}#1}\fi\par\vspace{7pt}}}
% Exercice résolu : énoncé en haut, \tcblower puis la solution sur fond crème
\newcounter{popexo}\newcounter{popexr}
\newtcolorbox{exoresolu}[1][]{popbase,colframe=poporange,colbacklower=popcream,
  segmentation style={popink,line width=1.2pt,dashed},
  before upper={\stepcounter{popexr}\popboxhead{Exercice résolu \thepopexr}{poporange}{#1}},
  before lower={\poppill{Solution}{popink}{popcream}\par\vspace{6pt}}}
% Exercice (non résolu en place) — la correction est dans la partie Corrigés
\newtcolorbox{exercice}[1][]{popbase,colframe=popink,
  before upper={\stepcounter{popexo}\popboxhead{Exercice \thepopexo}{popink}{#1}}}
% Corrigé
\newtcolorbox{corrige}[1][]{popbase,colframe=popgreen,colback=popgreenL,
  before upper={\popboxhead{Corrigé}{popgreen}{#1}}}
% Objectifs / prérequis / bilan
\newtcolorbox{objectifs}{popbase,colframe=popink,colback=poporangeL,
  before upper={\popboxhead{Objectifs de la leçon}{popink}{}}}
\newtcolorbox{prerequis}{popbase,colframe=popink,colback=popblueL,
  before upper={\popboxhead{Prérequis}{popblue}{}}}
\newtcolorbox{bilan}{popbase,colframe=popink,colback=white,boxrule=2.8pt,
  before upper={\popboxhead{Fiche bilan}{poporange}{L'essentiel à connaître pour l'examen}}}
% Figure encadrée avec légende
\newtcolorbox{popfigure}[1][]{enhanced,breakable=false,colback=white,colframe=popline,boxrule=1.4pt,arc=8pt,
  left=10pt,right=10pt,top=10pt,bottom=8pt,before skip=10pt,after skip=12pt,halign=center,
  lower separated=false,halign lower=center,
  fontlower=\popsemi\fontsize{9}{11.5}\selectfont\color{popmuted},
  #1}
\newcommand{\legende}[1]{\par\vspace{6pt}{\popsemi\fontsize{9}{11.5}\selectfont\color{popmuted}#1}}

% ============================================================
% Signature OnBuch+
% ============================================================
\newcommand{\poplogoinline}[1]{{\poplogo\fontsize{#1}{#1}\selectfont\color{popink}OnBuch\color{poporange}+}}
\newcommand{\signatureonbuch}{%
  \begin{tcolorbox}[enhanced,colback=popink,colframe=popink,boxrule=2.2pt,arc=10pt,
      drop shadow=poporange,left=14pt,right=14pt,top=10pt,bottom=10pt,before skip=0pt,after skip=0pt]
    \tikz[baseline=-0.7ex]\node[circle,fill=popgreen,draw=popcream,line width=1.3pt,minimum size=22pt,inner sep=0]{\popcheck{white}};%
    \hspace{9pt}\begin{minipage}[c]{0.62\linewidth}
      {\popxbold\fontsize{12}{13}\selectfont\color{popcream}Signé\ \textbullet\ {\poplogo OnBuch\color{poporange}+}}\par\vspace{2pt}
      {\raggedright\popsemi\fontsize{8}{10.5}\selectfont\color{popline}\MakeUppercase{Équipe pédagogique OnBuch+}\\\MakeUppercase{Conforme au programme officiel MINESEC}\par}
    \end{minipage}\hfill
    {\poplogo\fontsize{22}{22}\selectfont\color{popcream}OnBuch\color{poporange}+}
  \end{tcolorbox}%
}

% ============================================================
% Page de garde
% #1 titre · #2 matière · #3 niveau · #4 module · #5 n° leçon · #6 durée ·
% #7 description / objectifs (texte ou itemize)
% ============================================================
\newcommand{\pagegarde}[7]{%
  \thispagestyle{empty}
  \newgeometry{a4paper,margin=1.7cm,top=1.6cm,bottom=1.4cm}%
  \begin{tikzpicture}[remember picture,overlay]
    \fill[poporange!18] ([xshift=-3.2cm,yshift=-3.6cm]current page.north east) circle (4.6cm);
    \fill[poppurple!14] ([xshift=2.4cm,yshift=7.6cm]current page.south west) circle (3.8cm);
  \end{tikzpicture}%
  {\poplogo\fontsize{30}{30}\selectfont\color{popink}OnBuch\color{poporange}+}\hfill
  \raisebox{0.6ex}{\poppill{Cours signé \textbullet\ Premium}{popink}{popcream}}\par
  \vspace{1pt}\popsquiggle{poppurple}\par
  \vspace{8pt}
  \begin{tcolorbox}[enhanced,colback=poporange,colframe=popink,boxrule=2.8pt,arc=13pt,
      drop shadow=popink,left=18pt,right=18pt,top=12pt,bottom=13pt,after skip=10pt]
    \poppill{#2 \textbullet\ #3}{popink}{popcream}\hspace{5pt}\poppill{#4}{white}{popink}\par\vspace{8pt}
    {\popdisplay\fontsize{14}{14}\selectfont\color{popink}LEÇON #5}\par\vspace{4pt}
    {\raggedright\hyphenpenalty=10000\exhyphenpenalty=10000\popdisplay\fontsize{25}{27}\selectfont\color{popcream}\MakeUppercase{#1}\par}
    \vspace{8pt}
    {\popxbold\fontsize{9.5}{9.5}\selectfont\color{popink}\MakeUppercase{Durée conseillée : #6}}
  \end{tcolorbox}
  \begin{tcolorbox}[enhanced,colback=white,colframe=popink,boxrule=2.2pt,arc=10pt,
      drop shadow=popink,left=16pt,right=16pt,top=10pt,bottom=10pt,after skip=8pt]
    \poppill{Dans cette leçon}{poppurple}{white}\par\vspace{5pt}
    \popsemi\fontsize{9.3}{12.6}\selectfont\color{popink}#7
  \end{tcolorbox}
  \noindent\begin{minipage}[t]{0.487\linewidth}
    \begin{tcolorbox}[enhanced,colback=popgreen,colframe=popink,boxrule=2.2pt,arc=10pt,
        drop shadow=popink,left=11pt,right=11pt,top=10pt,bottom=10pt,height=3.1cm,valign=center]
      \tcbox[colback=white,colframe=popink,boxrule=1.3pt,arc=5pt,boxsep=0pt,nobeforeafter,
        left=3pt,right=3pt,top=3pt,bottom=3pt,valign=center]{\qrcode[height=1.9cm]{https://play.google.com/store/apps/details?id=com.luvvix.onbuch&hl=fr}}%
      \hspace{8pt}\begin{minipage}[c]{0.5\linewidth}
        {\popdisplay\fontsize{11}{12.5}\selectfont\color{white}\MakeUppercase{Télécharge OnBuch}}\par\vspace{4pt}
        {\raggedright\popsemi\fontsize{8.4}{11}\selectfont\color{white}Gratuit \textbullet\ Google Play\\Tuteur IA, annales, concours, résultats\par}
      \end{minipage}
    \end{tcolorbox}
  \end{minipage}\hfill
  \begin{minipage}[t]{0.487\linewidth}
    \begin{tcolorbox}[enhanced,colback=poppurple,colframe=popink,boxrule=2.2pt,arc=10pt,
        drop shadow=popink,left=11pt,right=11pt,top=10pt,bottom=10pt,height=3.1cm,valign=center]
      \tikz[baseline=-0.7ex]\node[circle,fill=white,draw=popink,line width=1.3pt,minimum size=1.6cm,inner sep=0]{\popdisplay\fontsize{24}{24}\selectfont\color{poppurple}+};%
      \hspace{8pt}\begin{minipage}[c]{0.6\linewidth}
        {\popdisplay\fontsize{11}{12.5}\selectfont\color{white}\MakeUppercase{Passe à OnBuch+}}\par\vspace{4pt}
        {\raggedright\popsemi\fontsize{8.4}{11}\selectfont\color{white}Tous les cours signés, Tuteur IA illimité, fiches PDF — onglet OnBuch+ de l'app\par}
      \end{minipage}
    \end{tcolorbox}
  \end{minipage}
  \vspace{8pt}
  \popcontenu
  \vfill
  \signatureonbuch
  \restoregeometry
  \newpage
}

% Rangée « Ce cours contient » (page de garde)
\newcommand{\poptile}[3]{% #1 couleur #2 gros texte #3 légende
  \begin{tcolorbox}[enhanced,colback=#1,colframe=popink,boxrule=2pt,arc=9pt,shadow={2.5pt}{-2.5pt}{0pt}{popink},
      left=6pt,right=6pt,top=9pt,bottom=9pt,halign=center,valign=center,height=1.75cm,width=\linewidth,nobeforeafter]
    {\popdisplay\fontsize{12.5}{13}\selectfont\color{white}\MakeUppercase{#2}}\par\vspace{3pt}
    {\centering\popsemi\fontsize{7.8}{9.5}\selectfont\color{white}#3\par}
  \end{tcolorbox}}
\newcommand{\popcontenu}{%
  \par\noindent{\popxboldls\fontsize{7.5}{7.5}\selectfont\color{popmuted}CE COURS SIGNÉ CONTIENT}\par\vspace{7pt}
  \noindent\begin{minipage}[t]{0.235\linewidth}\poptile{popblue}{Cours}{complet, illustré, pas à pas}\end{minipage}\hfill
  \begin{minipage}[t]{0.235\linewidth}\poptile{poporange}{Méthodes}{et exercices résolus}\end{minipage}\hfill
  \begin{minipage}[t]{0.235\linewidth}\poptile{poppink}{Exercices}{type examen + corrigés}\end{minipage}\hfill
  \begin{minipage}[t]{0.235\linewidth}\poptile{popgreen}{Fiche bilan}{l'essentiel à retenir}\end{minipage}\par}

% Sommaire
\newtcolorbox{sommaire}{enhanced,colback=white,colframe=popink,boxrule=2.2pt,arc=10pt,
  drop shadow=popink,left=18pt,right=18pt,top=13pt,bottom=13pt,before skip=6pt,after skip=20pt,
  before upper={\poppill{Sommaire}{poppurple}{white}\par\vspace{9pt}\popsemi\fontsize{10.6}{14.5}\selectfont\color{popink}}}

% Signature de fin de cours — poussée tout en bas de la dernière page
\newcommand{\findecours}{%
  \par\vfill\nopagebreak
  \begin{center}\popsquiggle{poporange}\end{center}\vspace{4pt}
  \signatureonbuch
}

%% FIGFIX-BEGIN v4 — correctifs globaux des figures (docs/onbuch-generation/tools/figfix)
\usepackage{adjustbox}\usepackage{etoolbox}
% toute figure TikZ/pgfplots plus large que la ligne est réduite à la largeur disponible
\BeforeBeginEnvironment{tikzpicture}{\begin{adjustbox}{max width=\linewidth}}
\AfterEndEnvironment{tikzpicture}{\end{adjustbox}}
% légendes pgfplots lisibles par-dessus les courbes
\pgfplotsset{every axis/.append style={legend style={fill=white,fill opacity=0.92,text opacity=1,draw=black!15,inner sep=3pt}}}
% étiquettes d'arêtes (syntaxe edge["..."]) avec fond blanc
\tikzset{every edge quotes/.append style={fill=white,inner sep=1.2pt}}
%% FIGFIX-END
\begin{document}

\pagegarde{Acides faibles, bases faibles et couples acide/base}{Chimie}{Tle C}{Module 2 · Acides et bases}{08}{5 h (cours + TP)}{Cette leçon étudie les acides faibles et les bases faibles à travers l'acide éthanoïque, l'éthanoate de sodium et l'ammoniac : ionisation partielle, équilibre chimique, taux d'avancement et coefficient d'ionisation. Elle introduit les couples acide/base, les constantes Ka, Kb et pKa, la relation pH = pKa + log([A-]/[AH]) et les domaines de prédominance, outils indispensables pour classer les couples et prévoir les réactions acido-basiques.
\vspace{4pt}
\begin{multicols}{2}\small
\begin{itemize}
\item À la fin de cette leçon, je sais montrer, à partir du pH et de la concentration, qu'un acide ou une base n'est pas totalement ionisé(e)
\item montrer que la dissolution d'un acide faible dans l'eau conduit à un équilibre chimique et l'étudier à l'aide d'un tableau d'avancement
\item déterminer le taux d'avancement final et le coefficient d'ionisation α d'un acide ou d'une base faible
\item identifier un couple acide/base à partir de l'équation d'un équilibre et citer les couples de l'eau et les ampholytes
\item définir et exploiter les constantes Ka, Kb, pKa et la relation Ka·Kb = Ke
\item établir et utiliser la relation pH = pKa + log([A-]/[AH]) pour tracer les domaines de prédominance
\end{itemize}\end{multicols}}

\begin{sommaire}
\begin{multicols}{2}
\begin{enumerate}
\item Acide faible : l'acide éthanoïque
\item Base faible : éthanoate de sodium et ammoniac
\item Couples acide/base, couples de l'eau et ampholytes
\item Constantes d'acidité Ka, de basicité Kb et pKa
\item Relation pH-pKa et domaines de prédominance
\item Prévision des réactions acido-basiques
\item Travaux pratiques
\item Méthodes et exercices résolus
\item Exercices
\item Corrigés des exercices
\item Fiche bilan
\end{enumerate}
\end{multicols}
\end{sommaire}

\begin{objectifs}
\begin{itemize}
\item À la fin de cette leçon, je sais montrer, à partir du pH et de la concentration, qu'un acide ou une base n'est pas totalement ionisé(e)
\item montrer que la dissolution d'un acide faible dans l'eau conduit à un équilibre chimique et l'étudier à l'aide d'un tableau d'avancement
\item déterminer le taux d'avancement final et le coefficient d'ionisation α d'un acide ou d'une base faible
\item identifier un couple acide/base à partir de l'équation d'un équilibre et citer les couples de l'eau et les ampholytes
\item définir et exploiter les constantes Ka, Kb, pKa et la relation Ka·Kb = Ke
\item établir et utiliser la relation pH = pKa + log([A-]/[AH]) pour tracer les domaines de prédominance
\item classer des couples acide/base selon Ka ou pKa et comparer la force relative des acides et des bases
\item prévoir le sens d'une réaction acido-basique à l'aide de K = 10\^{}(pKa2 - pKa1)
\item déterminer expérimentalement la constante d'acidité d'un couple par mesure de pH
\end{itemize}
\end{objectifs}

\begin{prerequis}
\begin{itemize}
\item Réactions acido-basiques et définitions de Brønsted (Première)
\item pH d'une solution aqueuse, relation pH = -log[H3O+] et produit ionique de l'eau Ke (début de Terminale)
\item Acides forts et bases fortes : ionisation totale, pH = -log C (leçon précédente)
\item Avancement d'une réaction, tableau d'avancement, taux d'avancement final (Première)
\item Constante d'équilibre et quotient de réaction (leçons sur les équilibres chimiques)
\end{itemize}
\end{prerequis}

\newpage

\coursec{Acide faible : l'acide éthanoïque}

Au marché Mokolo, une maman utilise du vinaigre (acide éthanoïque) pour nettoyer le calcaire de sa bouilloire, tandis qu'une autre soigne une brûlure d'estomac avec du bicarbonate. Pourtant, le vinaigre à 8° est bien moins « agressif » qu'un acide fort de même concentration : pourquoi ? De même, l'odeur piquante de l'ammoniac des produits ménagers révèle une base qui ne réagit que partiellement avec l'eau. Comment mesurer cette « faiblesse », la prévoir et l'exploiter ? C'est tout l'objet de cette leçon, au cœur du programme de Terminale et très prisée au Baccalauréat.

Dans cette première partie, nous allons découvrir, à travers l'exemple de l'\cle{acide éthanoïque} \ce{CH3COOH}, ce qu'est un \cle{acide faible}, comment prouver expérimentalement que son ionisation est partielle, et comment quantifier cette ionisation à l'aide du \cle{taux d'avancement final} et du \cle{coefficient d'ionisation}.

\courssub{Rappel : le cas de l'acide fort}

Avant de découvrir l'acide faible, rappelons ce que nous savons déjà d'un \cle{acide fort} comme l'acide chlorhydrique \ce{HCl}.

Lorsqu'on dissout du chlorure d'hydrogène dans l'eau, la réaction avec l'eau est \textbf{totale} : chaque molécule \ce{HCl} cède son proton à une molécule d'eau. On écrit alors avec une flèche simple :

\[ \ce{HCl + H2O -> Cl- + H3O+} \]

Cette transformation est \textbf{totale}, \textbf{rapide} et \textbf{exothermique}. Il ne reste \emph{aucune} molécule \ce{HCl} dans la solution : tout a été transformé en ions \ce{H3O+} et \ce{Cl-}.

\begin{propriete}[pH d'une solution d'acide fort]
Pour une solution d'acide fort de concentration molaire $C$ (avec $C \leq \SI{10^{-1}}{mol.L^{-1}}$ environ), la concentration en ions oxonium est égale à la concentration apportée :
\[ [\ce{H3O+}] = C \qquad \text{donc} \qquad \boxed{\mathrm{pH} = -\log C} \]
\end{propriete}

Par exemple, une solution d'acide chlorhydrique à $C = \SI{1,0e-2}{mol.L^{-1}}$ a un pH égal à $-\log(10^{-2}) = 2,0$. Retiens bien ce résultat : il va nous servir de \textbf{référence} pour démasquer l'acide faible.

\courssub{Expérience : mesure du pH d'une solution d'acide éthanoïque}

\begin{experience}[Mesure du pH d'une solution d'acide éthanoïque]
\textbf{Dispositif et protocole :} on prépare une solution aqueuse d'acide éthanoïque de concentration $C = \SI{1,0e-2}{mol.L^{-1}}$. On étalonne un pH-mètre avec des solutions étalons, puis on plonge l'électrode de verre dans la solution et on lit le pH à \SI{25}{\celsius}.

\textbf{Observation :} le pH-mètre affiche $\mathrm{pH} \approx 3,4$.

\textbf{Interprétation :} si l'acide éthanoïque était un acide fort, on aurait $\mathrm{pH} = -\log C = -\log(10^{-2}) = 2,0$. Or on mesure $3,4 > 2,0$. La solution contient donc \emph{moins} d'ions \ce{H3O+} que prévu :
\[ [\ce{H3O+}] = 10^{-\mathrm{pH}} = 10^{-3,4} \approx \SI{4,0e-4}{mol.L^{-1}} < C = \SI{1,0e-2}{mol.L^{-1}} \]
Toutes les molécules d'acide éthanoïque n'ont donc \textbf{pas} cédé leur proton : l'ionisation est \textbf{partielle}.
\end{experience}

\begin{methode}[Montrer expérimentalement qu'un acide est faible]
\begin{enumerate}
\item Mesurer le pH de la solution de concentration $C$ connue avec un pH-mètre étalonné.
\item Calculer la valeur de référence $-\log C$ (pH qu'aurait un acide fort de même concentration).
\item Comparer : si $\mathrm{pH} > -\log C$, alors $[\ce{H3O+}] < C$ : l'ionisation est partielle, l'acide est \textbf{faible}. Si $\mathrm{pH} = -\log C$, l'acide est fort.
\end{enumerate}
\end{methode}

\begin{definition}[Acide faible]
Un \cle{acide faible} est un acide dont la réaction avec l'eau est \textbf{partielle} (limitée, non totale). Pour une solution d'acide faible de concentration $C$ :
\[ [\ce{H3O+}] < C \qquad \Longleftrightarrow \qquad \mathrm{pH} > -\log C \]
Il reste dans la solution des molécules d'acide non ionisées, en équilibre avec les ions formés.
\end{definition}

\courssub{L'équilibre d'ionisation de l'acide éthanoïque}

Puisque la réaction n'est pas totale, on l'écrit avec une \textbf{double flèche}, symbole d'un \cle{équilibre chimique} :

\[ \ce{CH3COOH + H2O <=> CH3COO- + H3O+} \]

Lisons cette équation avec soin :
\begin{itemize}
\item Dans le \textbf{sens direct} (sens 1), la molécule d'acide éthanoïque \chemfig{CH_3-C(=[1]O)-[7]OH} cède un proton \ce{H+} à une molécule d'eau : il se forme l'\cle{ion éthanoate} \ce{CH3COO-} et l'\cle{ion oxonium} \ce{H3O+}.
\item Dans le \textbf{sens inverse} (sens 2), l'ion éthanoate \ce{CH3COO-} \emph{capture} un proton porté par \ce{H3O+} et régénère l'acide éthanoïque : c'est la preuve que la réaction inverse existe réellement.
\end{itemize}

\begin{attention}[Ne jamais écrire une flèche simple]
L'équation d'ionisation d'un acide faible s'écrit obligatoirement avec la double flèche \ce{<=>}. Écrire \ce{CH3COOH + H2O -> CH3COO- + H3O+} signifierait que la réaction est totale, ce qui est faux : c'est une erreur lourdement sanctionnée au Baccalauréat. À l'équilibre, les quatre espèces \ce{CH3COOH}, \ce{H2O}, \ce{CH3COO-} et \ce{H3O+} coexistent dans la solution.
\end{attention}

Comment être sûr qu'il s'agit bien d'un équilibre et non d'une réaction simplement « arrêtée » ? Deux arguments expérimentaux :
\begin{itemize}
\item \textbf{La réaction inverse existe :} si l'on dissout de l'éthanoate de sodium \ce{CH3COONa} dans l'eau, les ions \ce{CH3COO-} captent des protons de l'eau et la solution devient basique (nous y reviendrons dans la section 2). L'ion éthanoate se comporte donc bien comme une base capable de « remonter » la réaction.
\item \textbf{La dilution modifie le taux d'ionisation :} si l'on dilue la solution, la proportion de molécules ionisées \emph{augmente} (nous le verrons plus bas). Si la réaction était figée, la dilution ne changerait rien à cette proportion. C'est la signature d'un équilibre qui se déplace.
\end{itemize}

\begin{popfigure}
\begin{center}
\begin{tikzpicture}[scale=0.9]
% Bécher acide fort
\draw[thick, popdark] (0,0) -- (0,2.6) (2.6,0) -- (2.6,2.6) (0,0) -- (2.6,0);
\draw[popblue, thick] (0.06,1.7) -- (2.54,1.7);
\fill[popblueL, opacity=0.5] (0.06,0.05) rectangle (2.54,1.68);
% pastilles : toutes ionisées (H3O+)
\foreach \x/\y in {0.4/0.35, 1.0/0.55, 1.7/0.3, 2.2/0.6, 0.6/1.0, 1.3/1.2, 2.0/1.0, 0.9/1.45, 1.8/1.5, 0.35/1.35}{
  \fill[poporange] (\x,\y) circle (0.09);}
\node[popdark, font=\small\bfseries] at (1.3,-0.45) {Acide fort \ce{HCl}};
\node[popdark, font=\scriptsize] at (1.3,-0.85) {$C = \SI{1,0e-2}{mol.L^{-1}}$, pH $= 2,0$};
\node[poporange, font=\scriptsize] at (1.3,2.15) {100 \% ionisé : que des ions \ce{H3O+}};
% Bécher acide faible
\begin{scope}[xshift=6cm]
\draw[thick, popdark] (0,0) -- (0,2.6) (2.6,0) -- (2.6,2.6) (0,0) -- (2.6,0);
\draw[popgreen, thick] (0.06,1.7) -- (2.54,1.7);
\fill[popgreenL, opacity=0.5] (0.06,0.05) rectangle (2.54,1.68);
% molécules non ionisées (ronds verts)
\foreach \x/\y in {0.4/0.35, 1.0/0.55, 1.7/0.3, 2.2/0.6, 0.6/1.0, 1.3/1.2, 2.0/1.0, 0.9/1.45, 1.8/1.5, 0.35/1.35}{
  \fill[popgreen] (\x,\y) circle (0.09);}
% une seule ionisée
\fill[poporange] (1.45,0.75) circle (0.09);
\node[popdark, font=\small\bfseries] at (1.3,-0.45) {Acide faible \ce{CH3COOH}};
\node[popdark, font=\scriptsize] at (1.3,-0.85) {$C = \SI{1,0e-2}{mol.L^{-1}}$, pH $\approx 3,4$};
\node[popgreen!60!black, font=\scriptsize] at (1.3,2.15) {$\approx 4\%$ ionisé seulement};
\end{scope}
% légende pastilles
\fill[poporange] (0.2,-1.5) circle (0.09); \node[right, font=\scriptsize] at (0.4,-1.5) {ion \ce{H3O+} (forme ionisée)};
\fill[popgreen] (4.6,-1.5) circle (0.09); \node[right, font=\scriptsize] at (4.8,-1.5) {molécule \ce{CH3COOH} non ionisée};
\end{tikzpicture}
\end{center}
\legende{Comparaison schématique : dans l'acide fort, toutes les molécules sont ionisées ; dans l'acide éthanoïque, seule une petite fraction l'est.}
\end{popfigure}

\courssub{Étude quantitative : tableau d'avancement}

Pour quantifier cette ionisation partielle, utilisons l'outil que tu maîtrises déjà : le \cle{tableau d'avancement}. Considérons un volume $V$ de solution d'acide éthanoïque de concentration $C$. La quantité initiale d'acide est $n_0 = C \times V$.

\begin{center}
\begin{tabularx}{\linewidth}{|l|X|c|c|c|}
\hline
\rowcolor{popblueL}
\textbf{État} & \textbf{Avancement} & \ce{CH3COOH} & \ce{CH3COO-} & \ce{H3O+} \\
\hline
Initial & $x = 0$ & $CV$ & $0$ & $0$ \\
\hline
En cours & $x$ & $CV - x$ & $x$ & $x$ \\
\hline
Final (équilibre) & $x_f$ & $CV - x_f$ & $x_f$ & $x_f$ \\
\hline
Maximal (si total) & $x_{max}$ & $0$ & $CV$ & $CV$ \\
\hline
\end{tabularx}
\end{center}

L'eau est en large excès (c'est le solvant), elle n'apparaît pas dans le tableau. Deux avancements sont à distinguer soigneusement :

\begin{itemize}
\item L'\cle{avancement maximal} $x_{max}$ : c'est l'avancement qu'on atteindrait \emph{si} la réaction était totale. Ici, l'acide serait entièrement consommé : $CV - x_{max} = 0$, donc $x_{max} = CV$.
\item L'\cle{avancement final} $x_f$ : c'est l'avancement réellement atteint à l'équilibre. Or, d'après le tableau, $n_f(\ce{H3O+}) = x_f$. En mesurant le pH, on accède donc directement à $x_f$ :
\[ \boxed{x_f = [\ce{H3O+}] \times V = 10^{-\mathrm{pH}} \times V} \]
\end{itemize}

\begin{definition}[Taux d'avancement final et coefficient d'ionisation]
Le \cle{taux d'avancement final} $\tau$ compare l'avancement réel à l'avancement maximal :
\[ \tau = \frac{x_f}{x_{max}} = \frac{10^{-\mathrm{pH}} \times V}{C \times V} = \frac{10^{-\mathrm{pH}}}{C} \]
Pour un acide, on l'appelle aussi \cle{coefficient d'ionisation} (ou degré d'ionisation) et on le note $\alpha$ : $\alpha = \tau$. C'est la \textbf{proportion de molécules d'acide qui se sont ionisées}. On a toujours $0 \leq \alpha \leq 1$ (soit $0\,\% \leq \alpha \leq 100\,\%$).
\end{definition}

\begin{exemplebox}[Calcul du coefficient d'ionisation de l'acide éthanoïque]
Reprenons notre solution : $C = \SI{1,0e-2}{mol.L^{-1}}$ et $\mathrm{pH} = 3,4$.
\[ \alpha = \frac{10^{-\mathrm{pH}}}{C} = \frac{10^{-3,4}}{10^{-2}} = 10^{-1,4} \approx 4,0 \times 10^{-2} \]
Soit $\alpha \approx 4\,\%$ : sur 100 molécules d'acide éthanoïque dissoutes, seules 4 environ ont cédé leur proton à l'eau. Les 96 autres restent intactes. Voilà pourquoi le vinaigre est un acide « doux » !
\end{exemplebox}

\begin{attention}[L'erreur classique sur $\tau$]
Beaucoup d'élèves écrivent par mégarde $\tau = \dfrac{C}{10^{-\mathrm{pH}}}$ au lieu de $\tau = \dfrac{10^{-\mathrm{pH}}}{C}$. C'est l'inverse ! Avec notre exemple, on trouverait $\tau = 10^{-2}/10^{-3,4} = 25$, c'est-à-dire $2500\,\%$ : absurde. \textbf{Réflexe de survie :} vérifie toujours que $\tau \leq 1$. Si tu trouves un taux supérieur à 1 (ou négatif), tu t'es trompé quelque part. Autre piège : ne pas confondre $[\ce{H3O+}] = 10^{-\mathrm{pH}}$ avec $10^{+\mathrm{pH}}$.
\end{attention}

\begin{exoresolu}[Acide méthanoïque : fort ou faible ?]
On mesure le pH d'une solution d'acide méthanoïque \ce{HCOOH} (l'acide du venin des fourmis) de concentration $C = \SI{2,0e-3}{mol.L^{-1}}$ : on trouve $\mathrm{pH} = 3,2$.
\begin{enumerate}
\item Montrer que l'acide méthanoïque est un acide faible.
\item Écrire l'équation de sa réaction avec l'eau.
\item Calculer son coefficient d'ionisation $\alpha$.
\end{enumerate}
\tcblower
\textbf{1. Nature de l'acide.} Calculons la valeur de référence : $-\log C = -\log(\num{2,0e-3}) \approx 2,7$. Or le pH mesuré vaut $3,2 > 2,7$, donc $[\ce{H3O+}] = 10^{-3,2} \approx \SI{6,3e-4}{mol.L^{-1}} < C$. L'ionisation est partielle : l'acide méthanoïque est un \textbf{acide faible}.

\textbf{2. Équation de l'équilibre.} L'acide cède partiellement son proton à l'eau :
\[ \ce{HCOOH + H2O <=> HCOO- + H3O+} \]
Il se forme l'ion méthanoate \ce{HCOO-} et l'ion oxonium.

\textbf{3. Coefficient d'ionisation.}
\[ \alpha = \frac{10^{-\mathrm{pH}}}{C} = \frac{10^{-3,2}}{\num{2,0e-3}} = \frac{\num{6,3e-4}}{\num{2,0e-3}} \approx 0,32 \]
Soit $\alpha \approx 32\,\%$. Vérification : $\alpha < 1$, le résultat est cohérent. L'acide méthanoïque est nettement plus ionisé que l'acide éthanoïque ($4\,\%$) à concentration comparable : il est « moins faible ».
\end{exoresolu}

\courssub{Influence de la dilution : la loi d'Ostwald}

Reprenons notre solution d'acide éthanoïque et diluons-la dix fois, puis cent fois. Que devient $\alpha$ ? Les mesures donnent :

\begin{center}
\begin{tabularx}{\linewidth}{|c|X|X|X|}
\hline
\rowcolor{popblueL}
$C$ (\si{mol.L^{-1}}) & $10^{-2}$ & $10^{-3}$ & $10^{-4}$ \\
\hline
pH mesuré & $3,4$ & $3,9$ & $4,5$ \\
\hline
$\alpha = \dfrac{10^{-\mathrm{pH}}}{C}$ & $\approx 4\,\%$ & $\approx 13\,\%$ & $\approx 32\,\%$ \\
\hline
\end{tabularx}
\end{center}

Le constat est net : \textbf{plus la solution est diluée, plus le coefficient d'ionisation augmente}. Ce résultat porte le nom de \cle{loi de dilution d'Ostwald} (Wilhelm Ostwald, prix Nobel de chimie 1909). On l'énonce qualitativement ainsi : \emph{la dilution favorise l'ionisation d'un acide faible}. À dilution infinie, $\alpha$ tend vers 1 : l'acide faible finit par se comporter presque comme un acide fort.

\begin{attention}[Dilution et pH : ne pas tout mélanger]
Quand on dilue un acide faible, $\alpha$ \textbf{augmente}, mais le pH \textbf{augmente aussi} (la solution devient moins acide). En effet, $[\ce{H3O+}] = \alpha \times C$ : $\alpha$ croît, mais $C$ diminue beaucoup plus vite, donc le produit diminue. Diluer rend l'acide \emph{proportionnellement} plus ionisé, mais la solution est \emph{globalement} moins concentrée en \ce{H3O+}.
\end{attention}

\begin{popfigure}
\begin{center}
\begin{tikzpicture}
\begin{axis}[
    width=0.85\linewidth, height=6.5cm,
    xlabel={$C$ (\si{mol.L^{-1}})},
    ylabel={Coefficient d'ionisation $\alpha$ (\%)},
    xmin=0, xmax=0.012, ymin=0, ymax=50,
    xtick={0,0.002,0.004,0.006,0.008,0.01,0.012},
    xticklabels={$0$,{$2 \times 10^{-3}$},{$4 \times 10^{-3}$},{$6 \times 10^{-3}$},{$8 \times 10^{-3}$},{$1 \times 10^{-2}$},{$1,2 \times 10^{-2}$}},
    x tick label style={font=\scriptsize},
    ytick={0,10,20,30,40,50},
    grid=major, grid style={popline, dashed},
    axis line style={popdark},
    label style={popdark},
]
\addplot[popcurve, domain=0.0004:0.012, samples=200] {100*sqrt(1.8e-5/x)};
\addplot[only marks, mark=*, mark size=2pt, poporange] coordinates {(0.01,4.2) (0.001,13) (0.0001,32)};
\node[poporange, font=\scriptsize, anchor=south west] at (axis cs:0.0092,5.5) {points expérimentaux};
\node[popblue, font=\scriptsize, anchor=north west] at (axis cs:0.0006,44) {$\alpha$ croît quand $C$ diminue};
\draw[->, popblue, thick] (axis cs:0.004,30) -- (axis cs:0.0012,20);
\end{axis}
\end{tikzpicture}
\end{center}
\legende{Coefficient d'ionisation $\alpha$ de l'acide éthanoïque en fonction de la concentration $C$ : illustration de la loi d'Ostwald.}
\end{popfigure}

\courssub{Généralisation}

Ce que nous venons d'établir pour l'acide éthanoïque vaut pour \textbf{tous} les acides faibles. On note un acide faible générique \ce{AH} et sa base associée \ce{A-}.

\begin{aretenir}[Acide faible : l'essentiel]
\begin{itemize}
\item Un acide faible \ce{AH} s'ionise \textbf{partiellement} dans l'eau selon l'équilibre :
\[ \ce{AH + H2O <=> A- + H3O+} \]
\item Critère expérimental : $\mathrm{pH} > -\log C$, c'est-à-dire $[\ce{H3O+}] < C$.
\item Avancement final : $x_f = 10^{-\mathrm{pH}} \times V$ ; avancement maximal : $x_{max} = C \times V$.
\item Coefficient d'ionisation : $\alpha = \tau = \dfrac{10^{-\mathrm{pH}}}{C} < 1$.
\item Loi d'Ostwald : la dilution augmente $\alpha$.
\end{itemize}
\end{aretenir}

\begin{savaistu}[Le vinaigre, un acide faible de tous les jours]
Le vinaigre de table est une solution aqueuse d'acide éthanoïque à environ \SI{1}{mol.L^{-1}} (le fameux « 8° » signifie environ \SI{80}{g} d'acide par litre). Son acidité « douce », qui permet de l'utiliser en cuisine et pour détartrer la bouilloire sans danger, vient précisément de son ionisation partielle : à cette concentration, moins de $1\,\%$ des molécules sont ionisées ! Au Cameroun, le vinaigre peut être obtenu par fermentation acétique du vin de palme laissé à l'air : les bactéries du genre \emph{Acetobacter} transforment l'éthanol en acide éthanoïque. C'est aussi pour cela qu'un vin de palme « tourné » pique la langue : il s'est transformé en vinaigre.
\end{savaistu}

\begin{exercice}[À toi de jouer]
Une solution d'acide propanoïque \ce{C2H5COOH} a une concentration $C = \SI{5,0e-3}{mol.L^{-1}}$ et un pH de $3,6$.
\begin{enumerate}
\item Montrer que cet acide est faible.
\item Écrire l'équation de son équilibre avec l'eau.
\item Calculer son coefficient d'ionisation $\alpha$, puis les concentrations $[\ce{H3O+}]$, $[\ce{C2H5COO-}]$ et $[\ce{C2H5COOH}]$ à l'équilibre.
\end{enumerate}
\end{exercice}

\begin{corrige}[Exercice : acide propanoïque]
\textbf{1.} $-\log C = -\log(\num{5,0e-3}) \approx 2,3$. Comme $\mathrm{pH} = 3,6 > 2,3$, on a $[\ce{H3O+}] < C$ : l'ionisation est partielle, l'acide propanoïque est faible.

\textbf{2.} $\ce{C2H5COOH + H2O <=> C2H5COO- + H3O+}$ (double flèche obligatoire).

\textbf{3.} $\alpha = \dfrac{10^{-3,6}}{\num{5,0e-3}} = \dfrac{\num{2,5e-4}}{\num{5,0e-3}} = 5,0 \times 10^{-2}$, soit $\alpha = 5\,\%$. On en déduit : $[\ce{H3O+}] = [\ce{C2H5COO-}] = 10^{-3,6} \approx \SI{2,5e-4}{mol.L^{-1}}$ et $[\ce{C2H5COOH}] = C - [\ce{H3O+}] \approx \SI{4,75e-3}{mol.L^{-1}}$. Vérification : $\alpha < 1$, tout est cohérent.
\end{corrige}

\coursec{Base faible : éthanoate de sodium et ammoniac}

Dans la section précédente, nous avons découvert que l'acide éthanoïque \ce{CH3COOH} ne s'ionise que partiellement dans l'eau : son pH est \emph{supérieur} à la valeur $-\log C$ attendue pour un acide fort. Une question symétrique se pose alors naturellement : qu'en est-il du côté des \emph{bases} ? L'hydroxyde de sodium (la soude) est une base forte, totalement dissociée. Mais toutes les bases se comportent-elles ainsi ? Le savon artisanal fabriqué à partir d'huile de palme, l'ammoniac des produits de nettoyage, ou encore le goût légèrement savonneux d'une solution d'éthanoate de sodium vont nous fournir la réponse : \textbf{non}. Certaines bases ne réagissent que partiellement avec l'eau : ce sont les \cle{bases faibles}.

\courssub{Rappel : le cas de référence de la base forte}

Avant de découvrir ce qu'est une base faible, rappelons le comportement de la base forte par excellence : l'hydroxyde de sodium \ce{NaOH}. Dans l'eau, la dissolution est \emph{totale} :

\[ \ce{NaOH ->[H2O] Na+ + HO-} \]

Chaque mole de soude dissoute fournit exactement une mole d'ions hydroxyde \ce{HO-}. On a donc $[\ce{HO-}] = C$, d'où, à \SI{25}{\celsius} :

\begin{propriete}[Base forte : formule du pH]
Pour une solution de base forte de concentration $C$ (avec $C \gtrsim \SI{e-6}{mol.L^{-1}}$), à \SI{25}{\celsius} :
\[ \mathrm{pH} = 14 + \log C \]
\end{propriete}

Par exemple, pour $C = \SI{1,0e-2}{mol.L^{-1}}$ : $\mathrm{pH} = 14 + \log(10^{-2}) = 12$. Cette valeur servira de \textbf{référence} : toute base qui donnerait, à la même concentration, un pH \emph{inférieur} à 12 ne peut pas être totalement transformée en ions \ce{HO-}.

\courssub{Expérience : la solution d'éthanoate de sodium}

\begin{experience}[Mesure du pH d'une solution d'éthanoate de sodium]
\textbf{Dispositif et protocole.} On dissout de l'éthanoate de sodium \ce{CH3COONa} (solide ionique blanc) dans de l'eau distillée de façon à obtenir une solution de concentration $C = \SI{1,0e-2}{mol.L^{-1}}$. On mesure le pH à \SI{25}{\celsius} avec un pH-mètre étalonné.

\textbf{Observations.} Le pH-mètre affiche :
\[ \mathrm{pH} \approx 8{,}4 \]

\textbf{Interprétation.} Deux constats s'imposent :
\begin{itemize}
\item $\mathrm{pH} = 8{,}4 > 7$ : la solution est bien \textbf{basique}. Il s'est donc formé des ions \ce{HO-} : l'ion éthanoate a bien réagi avec l'eau ;
\item mais $\mathrm{pH} = 8{,}4 < 14 + \log C = 12$ : la solution est \emph{beaucoup moins basique} qu'une base forte de même concentration.
\end{itemize}
Calculons la concentration réelle en ions hydroxyde :
\[ [\ce{HO-}] = \frac{K_e}{[\ce{H3O+}]} = 10^{\mathrm{pH}-14} = 10^{8{,}4-14} = 10^{-5{,}6} \approx \SI{2,5e-6}{mol.L^{-1}} \]
Or $[\ce{HO-}] = \SI{2,5e-6}{mol.L^{-1}} \ll C = \SI{1,0e-2}{mol.L^{-1}}$. La quasi-totalité des ions éthanoate est \emph{restée sous forme} \ce{CH3COO-} : la réaction avec l'eau est \textbf{partielle}.
\end{experience}

\begin{definition}[Base faible]
Une \cle{base faible} est une base dont la réaction avec l'eau est \textbf{partielle} (limitée) : elle conduit à un \emph{équilibre chimique}. Pour une solution de base faible de concentration $C$, à \SI{25}{\celsius} :
\[ [\ce{HO-}] < C \qquad \text{et donc} \qquad \mathrm{pH} < 14 + \log C \]
\end{definition}

\courssub{Équation de l'équilibre : l'ion éthanoate dans l'eau}

L'éthanoate de sodium est un solide ionique : sa dissolution est totale et fournit des ions sodium \ce{Na+} et des ions éthanoate \ce{CH3COO-} :

\[ \ce{CH3COONa ->[H2O] CH3COO- + Na+} \]

L'ion sodium \ce{Na+} est un \textbf{ion spectateur} : il ne réagit pas avec l'eau, il se contente d'assurer l'électroneutralité de la solution. C'est l'ion éthanoate \ce{CH3COO-} qui joue le rôle de base : il \emph{capture} un proton \ce{H+} sur une molécule d'eau, selon l'équilibre :

\[ \ce{CH3COO- + H2O <=> CH3COOH + HO-} \]

Commentons cette équation, car elle est le cœur de la leçon :
\begin{itemize}
\item elle est écrite avec une \textbf{double flèche} \ce{<=>} : la réaction est \emph{limitée} et aboutit à un état d'équilibre, à température ordinaire, sans catalyseur ;
\item l'ion éthanoate \ce{CH3COO-} capte un proton : c'est une \textbf{base} (au sens de Brønsted) ;
\item l'eau \ce{H2O} cède un proton : elle joue ici le rôle d'\textbf{acide} ;
\item il se forme de l'acide éthanoïque \ce{CH3COOH} et des ions hydroxyde \ce{HO-}, responsables du caractère basique de la solution.
\end{itemize}

\begin{attention}[L'ion sodium est spectateur !]
Erreur très fréquente au Baccalauréat : écrire l'équation bilan avec \ce{CH3COONa} ou faire apparaître \ce{Na+} dans l'équilibre. L'ion \ce{Na+} \textbf{ne participe pas} à la réaction : il ne doit \emph{jamais} figurer dans l'équation de l'équilibre acido-basique. La bonne équation est :
\[ \ce{CH3COO- + H2O <=> CH3COOH + HO-} \]
et non \ce{CH3COONa + H2O <=> ...}. De même, ne jamais écrire une flèche simple \ce{->} : la réaction est partielle, la double flèche est obligatoire.
\end{attention}

\courssub{Cas de l'ammoniac \ce{NH3}}

L'ammoniac est la base faible la plus célèbre : c'est le principe actif de nombreux produits de nettoyage ménagers (nettoyants pour vitres, décapants). Dissous dans l'eau, il réagit partiellement selon :

\[ \ce{NH3 + H2O <=> NH4+ + HO-} \]

L'ammoniac \ce{NH3} capte un proton (base) ; l'eau le cède (acide) ; il se forme l'ion ammonium \ce{NH4+} et des ions hydroxyde. Là encore, la réaction est limitée, à température ordinaire, sans catalyseur.

Pour une solution d'ammoniac de concentration $C = \SI{1,0e-2}{mol.L^{-1}}$, le pH mesuré vaut environ $10{,}6$, valeur très inférieure à $14 + \log C = 12$ : l'ammoniac est bien une \textbf{base faible}. Remarquons au passage que $10{,}6 > 8{,}4$ : à même concentration, l'ammoniac est \emph{plus fortement} ionisé que l'ion éthanoate — c'est une base faible « moins faible ». Nous apprendrons à quantifier cette force relative avec le $pK_a$ (et le $pK_b$) dans les sections suivantes.

\begin{popfigure}
\begin{center}
\begin{tikzpicture}[scale=1.0, every node/.style={font=\small}]
% Flacon
\draw[thick, popblue, fill=popblueL] (-0.9,0) -- (-0.9,2.2) -- (-0.35,2.9) -- (-0.35,3.5) -- (0.35,3.5) -- (0.35,2.9) -- (0.9,2.2) -- (0.9,0) -- cycle;
% Niveau liquide
\draw[thick, popblue, fill=popblue!45] (-0.86,0.04) -- (-0.86,1.5) -- (0.86,1.5) -- (0.86,0.04) -- cycle;
% Bouchon
\draw[thick, popdark] (-0.45,3.5) rectangle (0.45,3.85);
% Étiquette
\node[draw=poporange, fill=popcream, rounded corners, align=center, font=\small\bfseries] at (0,0.8) {AMMONIAC\\ \ce{NH3} \SI{1,0e-2}{mol.L^{-1}}\\ pH $\approx 10{,}6$};
% Bulles de NH3 gazeux
\foreach \x/\y/\r in {-0.5/2.0/0.09, 0.1/2.15/0.07, 0.55/1.95/0.08, -0.15/2.45/0.06}{
  \draw[popgreen, fill=popgreenL] (\x,\y) circle (\r);}
\node[font=\scriptsize, popgreen!60!black] at (1.7,2.3) {\ce{NH3} (g)};
\draw[-{Stealth}, popgreen!60!black] (0.6,2.15) -- (1.35,2.3);
% Équation
\node[align=center] at (5.6,1.9) {\Large $\ce{NH3 + H2O <=> NH4+ + HO-}$};
% Double flèche mise en évidence
\draw[poporange, very thick] (4.7,1.62) -- (6.0,1.62);
\draw[poporange, very thick] (4.7,1.62) -- (4.9,1.72);
\draw[poporange, very thick] (4.7,1.62) -- (4.9,1.52);
\node[poporange, font=\small\bfseries, align=center] at (5.35,1.15) {réaction partielle :\\ équilibre chimique};
\node[align=left, font=\small] at (5.6,0.35) {$[\ce{HO-}] \approx \SI{4e-4}{mol.L^{-1}} \ll C$};
\end{tikzpicture}
\end{center}
\legende{Solution aqueuse d'ammoniac : l'équilibre \ce{NH3 + H2O <=> NH4+ + HO-} est limité ; une partie de \ce{NH3} s'échappe même sous forme gazeuse.}
\end{popfigure}

\begin{savaistu}[Pourquoi l'ammoniac sent-il si fort ?]
L'odeur âcre et piquante des produits de nettoyage à l'ammoniac s'explique directement par l'équilibre étudié ici. Dans le flacon, \ce{NH3} dissous coexiste avec \ce{NH4+}. Or \ce{NH3} est un gaz très volatil : une partie s'échappe en permanence de la solution sous forme gazeuse — c'est elle qui irrite vos narines. Ce dégagement \emph{déplace l'équilibre} \ce{NH3 + H2O <=> NH4+ + HO-} dans le sens de la formation de \ce{NH3} (sens inverse), ce qui entretient l'odeur. C'est aussi pour cela qu'il faut reboucher le flacon et ne jamais mélanger l'ammoniac avec l'eau de Javel (dégagement de gaz toxiques \ce{NH2Cl}, \ce{NCl3} !).
\end{savaistu}

\courssub{Étude quantitative : tableau d'avancement, taux d'avancement et coefficient d'ionisation}

Comme pour l'acide éthanoïque, l'outil privilégié pour étudier cet équilibre est le \textbf{tableau d'avancement}. Considérons un volume $V$ de solution d'ammoniac de concentration $C$.

\begin{center}
\begin{tabularx}{\linewidth}{|l|X|X|X|X|}
\hline
\rowcolor{popblueL}
Équation & \ce{NH3} & \ce{+ H2O <=>} & \ce{NH4+} & \ce{+ HO-} \\
\hline
État initial ($x=0$) & $C\,V$ & excès & $0$ & $0$ \\
\hline
État en cours ($x$) & $CV - x$ & excès & $x$ & $x$ \\
\hline
État final ($x_f$) & $CV - x_f$ & excès & $x_f$ & $x_f$ \\
\hline
\end{tabularx}
\end{center}

L'avancement final se déduit de la mesure du pH. En effet, $[\ce{HO-}]_f = \dfrac{K_e}{[\ce{H3O+}]} = 10^{\mathrm{pH}-14}$, donc :

\[ x_f = [\ce{HO-}]_f \times V = 10^{\mathrm{pH}-14} \times V \]

L'avancement maximal (si la réaction était totale) vaudrait $x_{\max} = C\,V$. On définit alors :

\begin{definition}[Taux d'avancement final et coefficient d'ionisation]
Le \cle{taux d'avancement final} $\tau$ est le rapport :
\[ \tau = \frac{x_f}{x_{\max}} = \frac{[\ce{HO-}]_f}{C} \]
Pour une base faible, on l'appelle aussi \cle{coefficient d'ionisation} (ou de protonation) et on le note $\alpha$ :
\[ \alpha = \frac{[\ce{HO-}]}{C} < 1 \]
$\alpha$ représente la \emph{fraction} des molécules de base qui ont effectivement capté un proton. Il s'exprime souvent en pourcentage.
\end{definition}

\begin{methode}[Calculer le coefficient d'ionisation d'une base faible]
\begin{enumerate}
\item Relever le pH mesuré et la concentration $C$ de la solution.
\item Calculer la concentration en ions hydroxyde : $[\ce{HO-}] = \dfrac{K_e}{[\ce{H3O+}]} = 10^{\mathrm{pH}-14}$ (à \SI{25}{\celsius}).
\item En déduire le coefficient d'ionisation : $\alpha = \dfrac{[\ce{HO-}]}{C}$.
\item Conclure : si $\alpha < 1$ (soit $< 100\,\%$), la base est faible ; plus $\alpha$ est proche de 1, plus la base est « forte ».
\end{enumerate}
\end{methode}

\begin{exoresolu}[Coefficient d'ionisation de l'ammoniac]
Une solution d'ammoniac de concentration $C = \SI{1,0e-2}{mol.L^{-1}}$ a un pH égal à $10{,}6$ à \SI{25}{\celsius}. Montrer que l'ammoniac est une base faible et calculer son coefficient d'ionisation $\alpha$.
\tcblower
\textbf{Étape 1 : concentration en ions hydroxyde.}
\[ [\ce{HO-}] = 10^{\mathrm{pH}-14} = 10^{10{,}6-14} = 10^{-3{,}4} \approx \SI{4,0e-4}{mol.L^{-1}} \]
\textbf{Étape 2 : comparaison avec $C$.}
\[ [\ce{HO-}] = \SI{4,0e-4}{mol.L^{-1}} < C = \SI{1,0e-2}{mol.L^{-1}} \]
La réaction \ce{NH3 + H2O <=> NH4+ + HO-} est donc \textbf{partielle} : l'ammoniac est une base faible.
\textbf{Étape 3 : coefficient d'ionisation.}
\[ \alpha = \frac{[\ce{HO-}]}{C} = \frac{4{,}0 \times 10^{-4}}{1{,}0 \times 10^{-2}} = 4{,}0 \times 10^{-2} \approx 4\,\% \]
\textbf{Conclusion.} Seules 4 molécules d'ammoniac sur 100 environ ont capté un proton : 96 \% restent sous forme \ce{NH3}. L'équilibre est très largement déplacé vers les réactifs.
\end{exoresolu}

Faisons de même pour l'éthanoate de sodium de la première expérience : $\alpha = \dfrac{10^{-5{,}6}}{10^{-2}} = 10^{-3{,}6} \approx 0{,}025\,\%$. L'ion éthanoate est une base \emph{extrêmement} faible — cohérent avec le fait que l'acide éthanoïque, son acide conjugué, est relativement « bon » donneur de proton.

\courssub{Influence de la dilution : la loi d'Ostwald}

Que se passe-t-il si l'on dilue la solution d'ammoniac ? L'expérience montre que le coefficient d'ionisation \textbf{augmente} lorsque la concentration $C$ diminue. Par exemple, pour l'ammoniac (valeurs cohérentes avec une constante d'équilibre $K_b$ constante) :

\begin{center}
\begin{tabularx}{\linewidth}{|X|X|X|}
\hline
\rowcolor{popblueL}
$C$ (\si{mol.L^{-1}}) & pH mesuré & $\alpha$ \\
\hline
$1{,}0 \times 10^{-1}$ & $11{,}1$ & $\approx 1{,}3\,\%$ \\
\hline
$1{,}0 \times 10^{-2}$ & $10{,}6$ & $\approx 4{,}0\,\%$ \\
\hline
$1{,}0 \times 10^{-3}$ & $10{,}1$ & $\approx 12{,}6\,\%$ \\
\hline
\end{tabularx}
\end{center}

\begin{propriete}[Loi de dilution d'Ostwald (notion)]
La dilution d'une base faible (ou d'un acide faible) \textbf{augmente son coefficient d'ionisation} $\alpha$ : plus la solution est diluée, plus la proportion de base ionisée est grande. À dilution infinie, $\alpha$ tend vers 1 (ionisation quasi totale), sans que la base devienne pour autant une base forte.
\end{propriete}

Intuitivement : en diluant, on « éloigne » les espèces les unes des autres, ce qui défavorise la recombinaison de \ce{NH4+} et \ce{HO-} et pousse l'équilibre vers la droite. Attention toutefois : $\alpha$ augmente, mais la \emph{concentration} $[\ce{HO-}]$ diminue (le pH se rapproche de 7) : il ne faut pas confondre \emph{proportion} ionisée et \emph{quantité} d'ions présents.

\courssub{Généralisation et comparaison des bases}

\begin{definition}[Généralisation]
Une base faible \ce{B} réagit partiellement avec l'eau selon l'équilibre général :
\[ \ce{B + H2O <=> BH+ + HO-} \]
où \ce{BH+} est l'\textbf{acide conjugué} de la base \ce{B}. La solution est basique ($\mathrm{pH} > 7$) mais vérifie toujours $\mathrm{pH} < 14 + \log C$ à \SI{25}{\celsius}.
\end{definition}

\begin{popfigure}
\begin{center}
\begin{tikzpicture}[scale=0.62]
% Axe pH
\draw[-{Stealth}, thick, popdark] (0,0) -- (15,0);
\foreach \x in {0,1,...,14}{
  \draw[popdark] (\x,-0.12) -- (\x,0.12);
  \node[below, font=\scriptsize, popdark] at (\x,-0.15) {\x};}
\node[below, font=\small, popdark] at (7.5,-0.85) {pH à \SI{25}{\celsius} pour $C = \SI{1,0e-2}{mol.L^{-1}}$};
% Zones
\fill[popblue!15] (0,0.25) rectangle (7,0.75);
\fill[poporange!18] (7,0.25) rectangle (14,0.75);
\node[font=\scriptsize, popblue!70!black] at (3.5,0.5) {milieu acide};
\node[font=\scriptsize, poporange!80!black] at (10.5,0.5) {milieu basique};
% Points
\fill[poppurple] (12,0) circle (0.18);
\node[above, font=\small\bfseries, poppurple, align=center] at (12,0.28) {Base forte (\ce{NaOH})\\ pH $= 12$};
\fill[popgreen!70!black] (10.6,0) circle (0.18);
\node[above, font=\small\bfseries, popgreen!50!black, align=center] at (8.9,1.35) {Ammoniac \ce{NH3}\\ pH $\approx 10{,}6$};
\draw[-{Stealth}, popgreen!50!black] (9.6,1.15) -- (10.5,0.25);
\fill[popgold!80!black] (8.4,0) circle (0.18);
\node[below, font=\small\bfseries, popgold!60!black, align=center] at (6.6,-1.55) {Éthanoate de sodium \ce{CH3COO-}\\ pH $\approx 8{,}4$};
\draw[-{Stealth}, popgold!60!black] (7.4,-1.35) -- (8.3,-0.25);
\end{tikzpicture}
\end{center}
\legende{Comparaison des pH de trois solutions basiques de même concentration $C = \SI{1,0e-2}{mol.L^{-1}}$ : plus la base est faible, plus le pH s'éloigne de la valeur de référence $14 + \log C = 12$.}
\end{popfigure}

\begin{aretenir}[L'essentiel sur les bases faibles]
\begin{itemize}
\item Une \cle{base faible} réagit \textbf{partiellement} avec l'eau : \ce{B + H2O <=> BH+ + HO-} (double flèche obligatoire).
\item Critère expérimental : $\mathrm{pH} < 14 + \log C$, c'est-à-dire $[\ce{HO-}] < C$.
\item Exemples de référence : l'ion éthanoate \ce{CH3COO-} (pH $\approx 8{,}4$) et l'ammoniac \ce{NH3} (pH $\approx 10{,}6$) pour $C = \SI{1,0e-2}{mol.L^{-1}}$.
\item Le coefficient d'ionisation se calcule par $\alpha = \dfrac{[\ce{HO-}]}{C} = \dfrac{10^{\mathrm{pH}-14}}{C} < 1$.
\item La dilution augmente $\alpha$ (loi d'Ostwald) mais fait diminuer le pH vers 7.
\item Les ions « spectateurs » comme \ce{Na+} n'apparaissent jamais dans l'équation de l'équilibre.
\end{itemize}
\end{aretenir}

\begin{exercice}[Savonnier de quartier]
Un artisan savonnier de Bafoussam prépare une lessive de soude diluée, puis la compare à une solution d'ammoniac. La solution d'ammoniac utilisée a pour concentration $C = \SI{5,0e-3}{mol.L^{-1}}$ et un pH mesuré de $10{,}4$ à \SI{25}{\celsius}.
\begin{enumerate}
\item Écrire l'équation de la réaction de l'ammoniac avec l'eau.
\item Calculer $[\ce{H3O+}]$ puis $[\ce{HO-}]$ dans cette solution.
\item Calculer le coefficient d'ionisation $\alpha$ de l'ammoniac. Conclure.
\item Quel serait le pH si l'ammoniac était une base forte de même concentration ?
\end{enumerate}
\end{exercice}

\begin{corrige}[Exercice : Savonnier de quartier]
\textbf{1.} Équilibre de protonation de l'ammoniac :
\[ \ce{NH3 + H2O <=> NH4+ + HO-} \]
\textbf{2.} $[\ce{H3O+}] = 10^{-\mathrm{pH}} = 10^{-10{,}4} \approx \SI{4,0e-11}{mol.L^{-1}}$.
\[ [\ce{HO-}] = \frac{K_e}{[\ce{H3O+}]} = \frac{10^{-14}}{10^{-10{,}4}} = 10^{-3{,}6} \approx \SI{2,5e-4}{mol.L^{-1}} \]
\textbf{3.} $\alpha = \dfrac{[\ce{HO-}]}{C} = \dfrac{2{,}5 \times 10^{-4}}{5{,}0 \times 10^{-3}} = 5{,}0 \times 10^{-2} = 5\,\%$. Comme $\alpha < 1$, l'ammoniac est bien une base faible.
\textbf{4.} Si la base était forte : $\mathrm{pH} = 14 + \log C = 14 + \log(5{,}0 \times 10^{-3}) = 14 - 2{,}3 = 11{,}7$. Le pH réel ($10{,}4$) est bien inférieur, ce qui confirme le caractère faible de la base.
\end{corrige}

Nous savons désormais reconnaître et quantifier une base faible, de manière parfaitement symétrique à l'acide faible de la section précédente. Mais remarquons que l'ion éthanoate \ce{CH3COO-} et l'acide éthanoïque \ce{CH3COOH} se transforment l'un en l'autre par gain ou perte d'un proton : ils forment un \emph{couple}. C'est cette notion fondamentale de \textbf{couple acide/base} que nous allons formaliser dans la section suivante.

\coursec{Couples acide/base, couples de l'eau et ampholytes}

Dans la section précédente, nous avons vu que l'ion \cle{éthanoate} \ce{CH3COO-} se comporte en base faible : il capte un proton de l'eau selon l'équilibre \ce{CH3COO- + H2O <=> CH3COOH + HO-}. Regarde bien cette équation : à gauche, \ce{CH3COO-} capte un proton ; à droite, on retrouve \ce{CH3COOH}, qui n'est autre que l'acide éthanoïque étudié dans la première section. L'acide et la base semblent donc être les deux visages d'une même entité, reliés par un simple proton \ce{H+}. C'est exactement l'idée géniale de la théorie de Brønsted : on ne peut pas parler d'un acide sans parler de sa \emph{base conjuguée}. Cette section formalise cette parenté à travers la notion de \cle{couple acide/base}, notion centrale de tout le chapitre, qui te servira jusqu'au Baccalauréat et bien au-delà.

\courssub{La notion de couple acide/base}

\textbf{L'intuition d'abord.} Imagine une paire de gants : un gant droit et un gant gauche. Ils ne sont identiques qu'à un détail près, et l'un se transforme en l'autre par une opération simple. Un couple acide/base, c'est pareil : deux espèces chimiques qui se transforment l'une en l'autre par la \emph{perte} ou le \emph{gain} d'un seul proton \ce{H+}. L'espèce qui \textbf{cède} le proton est l'acide ; celle qui le \textbf{capte} est la base.

\begin{definition}[Couple acide/base]
Deux espèces chimiques \ce{AH} et \ce{A-} forment un \cle{couple acide/base} si elles s'échangent un proton \ce{H+} selon la demi-équation :
\[ \ce{AH <=> A- + H+} \]
\ce{AH} est l'\cle{acide} du couple, \ce{A-} est sa \cle{base conjuguée}. Le couple se note \textbf{acide/base}, c'est-à-dire \ce{AH}/\ce{A-} : \emph{l'acide est toujours écrit en premier}.
\end{definition}

La flèche double \ce{<=>} est essentielle : elle signifie que la transformation est \emph{réversible}. L'acide \ce{AH} peut céder son proton, et la base \ce{A-} peut le récupérer. C'est cette réversibilité qui explique le caractère limité (équilibre) des réactions des acides et bases faibles vues dans les sections précédentes.

\begin{exemplebox}[Le couple de l'acide éthanoïque]
L'acide éthanoïque \ce{CH3COOH} et l'ion éthanoate \ce{CH3COO-} forment le couple :
\[ \ce{CH3COOH <=> CH3COO- + H+} \qquad \text{couple noté } \ce{CH3COOH}/\ce{CH3COO-} \]
On passe de l'un à l'autre par perte ou gain d'un unique \ce{H+} : la charge passe de $0$ à $-1$, tout le reste de la molécule est inchangé.
\end{exemplebox}

\begin{attention}[La convention d'écriture : acide en premier !]
L'erreur la plus fréquente au Baccalauréat est d'écrire le couple \emph{à l'envers} : \ce{CH3COO-}/\ce{CH3COOH}. C'est \textbf{faux}. La convention internationale impose d'écrire d'abord l'\textbf{acide}, puis la \textbf{base} : \ce{CH3COOH}/\ce{CH3COO-}. Moyen mnémotechnique : \og A comme Acide, A comme Avant \fg. Un couple mal écrit peut te coûter des points même si tu as identifié les bonnes espèces.
\end{attention}

\courssub{Identifier les couples dans un équilibre acido-basique}

Quand on te donne l'équation d'une réaction acido-basique, comment retrouver les deux couples mis en jeu ? La démarche est systématique.

\begin{methode}[Identifier les couples d'un équilibre]
\begin{enumerate}
\item \textbf{Éliminer les ions spectateurs} : barre mentalement les ions qui apparaissent à l'identique des deux côtés de l'équation (comme \ce{Na+} ou \ce{Cl-}) ; ils ne participent pas à l'échange de proton.
\item \textbf{Repérer les paires d'espèces qui ne diffèrent que d'un seul \ce{H+}} : même squelette atomique, une unité de charge et un hydrogène d'écart.
\item \textbf{Dans chaque paire, identifier l'acide} : c'est l'espèce qui possède le proton en plus (la plus protonée, la moins chargée négativement).
\item \textbf{Écrire les couples} sous la forme acide/base.
\end{enumerate}
\end{methode}

\begin{exemplebox}[Application : l'ammoniac dans l'eau]
Considérons l'équilibre étudié à la section précédente :
\[ \ce{NH3 + H2O <=> NH4+ + HO-} \]
\begin{itemize}
\item Il n'y a pas d'ion spectateur.
\item Paire 1 : \ce{NH3} et \ce{NH4+} ne diffèrent que d'un \ce{H+}. L'acide est \ce{NH4+} (il a le proton en plus). Couple : \ce{NH4+}/\ce{NH3}.
\item Paire 2 : \ce{H2O} et \ce{HO-} ne diffèrent que d'un \ce{H+}. L'acide est \ce{H2O}. Couple : \ce{H2O}/\ce{HO-}.
\end{itemize}
Dans cette réaction, \ce{H2O} a \emph{cédé} un proton à \ce{NH3} : l'eau s'est comportée en acide.
\end{exemplebox}

Voici d'autres couples usuels à connaître parfaitement :

\begin{center}
\begin{tabularx}{\linewidth}{|l|X|X|}
\hline
\rowcolor{popblueL} \textbf{Couple acide/base} & \textbf{Demi-équation} & \textbf{Exemple d'espèces} \\
\hline
\ce{NH4+}/\ce{NH3} & \ce{NH4+ <=> NH3 + H+} & ion ammonium / ammoniac \\
\hline
\ce{H2CO3}/\ce{HCO3-} & \ce{H2CO3 <=> HCO3- + H+} & acide carbonique / ion hydrogénocarbonate \\
\hline
\ce{HCO3-}/\ce{CO3^2-} & \ce{HCO3- <=> CO3^2- + H+} & ion hydrogénocarbonate / ion carbonate \\
\hline
\ce{HCOOH}/\ce{HCOO-} & \ce{HCOOH <=> HCOO- + H+} & acide méthanoïque / ion méthanoate \\
\hline
\end{tabularx}
\end{center}

Remarque au passage que l'ion \ce{HCO3-} apparaît dans \emph{deux} couples : tantôt base (face à \ce{H2CO3}), tantôt acide (face à \ce{CO3^2-}). Nous y reviendrons : c'est un ampholyte.

\begin{popfigure}
\begin{center}
\begin{tikzpicture}[>=Stealth, scale=1.0, every node/.style={font=\small}]
% Axe vertical invisible pour alignement
% Colonnes : Acide (gauche) -- Flèche -- Base (droite)
% Couple 1 : Acide éthanoïque / Éthanoate
\node[font=\bfseries, popblue, anchor=east] (a1) at (-1.5,5.5) {\ce{CH3COOH}};
\node[font=\bfseries, poporange, anchor=west] (b1) at (1.5,5.5) {\ce{CH3COO-}};
\draw[->, thick, popblue] (a1) -- (b1) node[midway, above, popdark] {$-\,\ce{H+}$};
\draw[<-, thick, poporange, dashed] (a1) -- (b1) node[midway, below, popdark] {$+\,\ce{H+}$};

% Couple 2 : Ammonium / Ammoniac
\node[font=\bfseries, popblue, anchor=east] (a2) at (-1.5,3.5) {\ce{NH4+}};
\node[font=\bfseries, poporange, anchor=west] (b2) at (1.5,3.5) {\ce{NH3}};
\draw[->, thick, popblue] (a2) -- (b2) node[midway, above, popdark] {$-\,\ce{H+}$};
\draw[<-, thick, poporange, dashed] (a2) -- (b2) node[midway, below, popdark] {$+\,\ce{H+}$};

% Couple 3 : Oxonium / Eau
\node[font=\bfseries, popblue, anchor=east] (a3) at (-1.5,1.5) {\ce{H3O+}};
\node[font=\bfseries, poporange, anchor=west] (b3) at (1.5,1.5) {\ce{H2O}};
\draw[->, thick, popblue] (a3) -- (b3) node[midway, above, popdark] {$-\,\ce{H+}$};
\draw[<-, thick, poporange, dashed] (a3) -- (b3) node[midway, below, popdark] {$+\,\ce{H+}$};

% Couple 4 : Eau / Hydroxyde
\node[font=\bfseries, popblue, anchor=east] (a4) at (-1.5,-0.5) {\ce{H2O}};
\node[font=\bfseries, poporange, anchor=west] (b4) at (1.5,-0.5) {\ce{HO-}};
\draw[->, thick, popblue] (a4) -- (b4) node[midway, above, popdark] {$-\,\ce{H+}$};
\draw[<-, thick, poporange, dashed] (a4) -- (b4) node[midway, below, popdark] {$+\,\ce{H+}$};

% Étiquettes de colonnes
\node[font=\small\bfseries, popblue, anchor=south] at (-1.5,6.3) {ACIDE (cède \ce{H+})};
\node[font=\small\bfseries, poporange, anchor=south] at (1.5,6.3) {BASE CONJUGUÉE (capte \ce{H+})};
% Accolade pour l'eau
\draw[popdark, decorate, decoration={brace, amplitude=5pt, mirror}] (-2.2,1.5) -- (-2.2,-0.5) node[midway, left=5pt, font=\small\itshape, poppurple] {Eau ampholyte};
\end{tikzpicture}
\end{center}
\legende{Quelques couples acide/base. À gauche (bleu), l'acide ; à droite (orange), sa base conjuguée. Le passage de l'acide à la base se fait par perte d'un proton \ce{H+} (flèche pleine) ; le passage inverse par gain d'un proton (flèche pointillée). Note que l'eau apparaît deux fois : tantôt base (couple \ce{H3O+}/\ce{H2O}), tantôt acide (couple \ce{H2O}/\ce{HO-}).}
\end{popfigure}

\courssub{Les couples de l'eau}

L'eau occupe une place tout à fait particulière en chimie acido-basique : elle intervient dans \textbf{deux couples}.

\textbf{Premier couple : l'eau comme base.} Face à un acide, l'eau \emph{capte} un proton pour donner l'ion oxonium \ce{H3O+} :
\[ \ce{H3O+ <=> H2O + H+} \qquad \text{couple } \ce{H3O+}/\ce{H2O} \]
C'est ce couple qui intervient quand on dissout l'acide éthanoïque dans l'eau : \ce{CH3COOH + H2O <=> CH3COO- + H3O+}. Ici, \ce{H2O} joue le rôle de base.

\textbf{Second couple : l'eau comme acide.} Face à une base, l'eau \emph{cède} un proton et devient ion hydroxyde \ce{HO-} :
\[ \ce{H2O <=> HO- + H+} \qquad \text{couple } \ce{H2O}/\ce{HO-} \]
C'est ce couple qui intervient quand on dissout l'éthanoate de sodium ou l'ammoniac dans l'eau : \ce{NH3 + H2O <=> NH4+ + HO-}. Ici, \ce{H2O} joue le rôle d'acide.

\begin{aretenir}[Les deux couples de l'eau]
\begin{itemize}
\item Couple \ce{H3O+}/\ce{H2O} : l'eau est la \textbf{base} (elle capte \ce{H+} pour donner \ce{H3O+}).
\item Couple \ce{H2O}/\ce{HO-} : l'eau est l'\textbf{acide} (elle cède \ce{H+} pour donner \ce{HO-}).
\end{itemize}
L'eau peut donc se comporter en acide \emph{ou} en base selon le partenaire : c'est un \cle{ampholyte}.
\end{aretenir}

\courssub{Les ampholytes (espèces amphotères)}

\begin{definition}[Ampholyte]
Un \cle{ampholyte} (ou espèce \cle{amphotère}) est une espèce chimique qui peut se comporter \textbf{soit en acide, soit en base} selon l'espèce avec laquelle elle réagit. Elle appartient donc à \textbf{deux couples} acide/base : elle est la base de l'un et l'acide de l'autre.
\end{definition}

L'eau est l'ampholyte par excellence. Mais il en existe beaucoup d'autres, en général des ions \emph{hydrogénés} intermédiaires d'un diacide ou d'un triacide :

\begin{center}
\begin{tabularx}{\linewidth}{|l|X|X|}
\hline
\rowcolor{popblueL} \textbf{Ampholyte} & \textbf{Comportement acide (cède \ce{H+})} & \textbf{Comportement basique (capte \ce{H+})} \\
\hline
\ce{H2O} & \ce{H2O <=> HO- + H+} \newline couple \ce{H2O}/\ce{HO-} & \ce{H2O + H+ <=> H3O+} \newline couple \ce{H3O+}/\ce{H2O} \\
\hline
\ce{HCO3-} & \ce{HCO3- <=> CO3^2- + H+} \newline couple \ce{HCO3-}/\ce{CO3^2-} & \ce{HCO3- + H+ <=> H2CO3} \newline couple \ce{H2CO3}/\ce{HCO3-} \\
\hline
\ce{HSO4-} & \ce{HSO4- <=> SO4^2- + H+} \newline couple \ce{HSO4-}/\ce{SO4^2-} & \ce{HSO4- + H+ <=> H2SO4} \newline couple \ce{H2SO4}/\ce{HSO4-} \\
\hline
\ce{H2PO4-} & \ce{H2PO4- <=> HPO4^2- + H+} \newline couple \ce{H2PO4-}/\ce{HPO4^2-} & \ce{H2PO4- + H+ <=> H3PO4} \newline couple \ce{H3PO4}/\ce{H2PO4-} \\
\hline
\end{tabularx}
\end{center}

\begin{exemplebox}[L'ion hydrogénocarbonate, ampholyte du quotidien]
Le bicarbonate de sodium \ce{NaHCO3}, utilisé en pharmacie contre les brûlures d'estomac et en pâtisserie, contient l'ion \ce{HCO3-}.
\begin{itemize}
\item Face à l'acide chlorhydrique de l'estomac, il joue le rôle de \textbf{base} : \ce{HCO3- + H3O+ -> H2CO3 + H2O} (réaction quasi totale, le \ce{H2CO3} se décompose en \ce{CO2} et \ce{H2O} : c'est l'effervescence).
\item Face à une base forte comme la soude, il joue le rôle d'\textbf{acide} : \ce{HCO3- + HO- -> CO3^2- + H2O}.
\end{itemize}
Même espèce, deux comportements : c'est la définition même d'un ampholyte.
\end{exemplebox}

\begin{savaistu}[Le sang, un milieu tamponné grâce aux ampholytes]
Ton sang doit maintenir son pH entre $7{,}35$ et $7{,}45$ : une variation de $0{,}2$ unité peut être mortelle ! Ce tour de force est assuré par des ampholytes. D'une part, le couple \ce{H2CO3}/\ce{HCO3-} : l'ion hydrogénocarbonate neutralise les excès d'acide produits par l'effort musculaire. D'autre part, les \textbf{acides aminés} constituant les protéines (comme l'hémoglobine) possèdent à la fois un groupe acide \ce{-COOH} et un groupe basique \ce{-NH2} : ce sont des ampholytes qui absorbent ou libèrent des protons selon les besoins. La chimie acido-basique de cette leçon est littéralement ce qui te maintient en vie à chaque instant.
\end{savaistu}

\courssub{Toute réaction acido-basique met en jeu deux couples}

Arrivons à la synthèse de cette section, qui est l'un des savoirs les plus importants du chapitre. Un proton \ce{H+} ne circule jamais seul en solution : il est \emph{transféré} d'un acide vers une base. Il faut donc \textbf{deux couples} pour qu'une réaction acido-basique ait lieu : l'acide du couple 1 cède son proton à la base du couple 2.

\begin{propriete}[Équation générale d'une réaction acido-basique]
Une réaction acido-basique est un \cle{transfert de proton} entre l'acide d'un couple et la base d'un autre couple :
\[ \ce{acide1 + base2 <=> base1 + acide2} \]
soit, avec les couples \ce{AH}/\ce{A-} et \ce{BH+}/\ce{B} :
\[ \ce{AH + B <=> A- + BH+} \]
Elle résulte de la somme des deux demi-équations : \ce{AH <=> A- + H+} et \ce{B + H+ <=> BH+}. Le proton échangé n'apparaît pas dans l'équation bilan.
\end{propriete}

\begin{exoresolu}[Identifier les couples et les rôles dans un équilibre]
Énoncé. On dissout du gaz ammoniac \ce{NH3} dans de l'eau. L'équilibre obtenu s'écrit :
\[ \ce{NH3 + H2O <=> NH4+ + HO-} \]
\begin{enumerate}
\item Identifier les deux couples acide/base mis en jeu.
\item Préciser, pour chaque réactif, s'il joue le rôle d'acide ou de base.
\item Écrire les deux demi-équations dont la somme redonne l'équation bilan.
\end{enumerate}
\tcblower
Solution détaillée.
\begin{enumerate}
\item Cherchons les paires d'espèces ne différant que d'un proton \ce{H+} :
\begin{itemize}
\item \ce{NH3} et \ce{NH4+} : même atome d'azote, un \ce{H+} d'écart. L'espèce la plus protonée, \ce{NH4+}, est l'acide. Couple : \ce{NH4+}/\ce{NH3}.
\item \ce{H2O} et \ce{HO-} : un \ce{H+} d'écart. L'espèce la plus protonée, \ce{H2O}, est l'acide. Couple : \ce{H2O}/\ce{HO-}.
\end{itemize}
\item Rôles : \ce{NH3} \emph{capte} un proton (elle devient \ce{NH4+}) : c'est la \textbf{base} du couple \ce{NH4+}/\ce{NH3}. \ce{H2O} \emph{cède} un proton (elle devient \ce{HO-}) : c'est l'\textbf{acide} du couple \ce{H2O}/\ce{HO-}. Ici, l'eau se comporte en acide : elle joue son rôle d'ampholyte.
\item Demi-équations :
\begin{align*}
\ce{H2O &<=> HO- + H+} && \text{(l'eau, acide, cède \ce{H+})}\\
\ce{NH3 + H+ &<=> NH4+} && \text{(l'ammoniac, base, capte \ce{H+})}
\end{align*}
En additionnant membre à membre, le proton \ce{H+} se simplifie et l'on retrouve bien :
\[ \ce{NH3 + H2O <=> NH4+ + HO-} \]
On vérifie la structure générale : $\text{base}_2 + \text{acide}_1 \ce{<=>} \text{acide}_2 + \text{base}_1$.
\end{enumerate}
\end{exoresolu}

\begin{attention}[Pièges classiques à éviter]
\begin{itemize}
\item \textbf{Ne confonds pas} \og base d'un couple \fg{} et \og solution basique \fg{} : dire que \ce{H2O} est la base du couple \ce{H3O+}/\ce{H2O} ne signifie pas que l'eau pure est basique (son pH vaut $7{,}0$ à \SI{25}{\celsius}).
\item \textbf{Un couple n'existe que si les deux espèces ne diffèrent que d'un seul \ce{H+}} : \ce{H2SO4} et \ce{SO4^2-} ne forment \emph{pas} un couple (deux protons d'écart) ; les couples sont \ce{H2SO4}/\ce{HSO4-} et \ce{HSO4-}/\ce{SO4^2-}.
\item \textbf{Vérifie toujours la conservation} : dans une demi-équation \ce{AH <=> A- + H+}, les atomes et les charges doivent être équilibrés des deux côtés.
\end{itemize}
\end{attention}

\begin{exercice}[À toi de jouer]
L'acide méthanoïque \ce{HCOOH} (présent dans le venin des fourmis) réagit avec l'ion éthanoate selon l'équilibre :
\[ \ce{HCOOH + CH3COO- <=> HCOO- + CH3COOH} \]
\begin{enumerate}
\item Identifier les deux couples acide/base mis en jeu.
\item Écrire les deux demi-équations protoniques correspondantes.
\item L'ion \ce{HCO3-} peut-il réagir avec \ce{HCOOH} ? Avec \ce{HO-} ? Écrire les équations et conclure sur la nature de \ce{HCO3-}.
\end{enumerate}
\end{exercice}

\begin{corrige}[Exercice]
\begin{enumerate}
\item Couples : \ce{HCOOH}/\ce{HCOO-} (l'acide \ce{HCOOH} cède son proton) et \ce{CH3COOH}/\ce{CH3COO-} (la base \ce{CH3COO-} capte le proton).
\item Demi-équations : \ce{HCOOH <=> HCOO- + H+} et \ce{CH3COO- + H+ <=> CH3COOH}. Leur somme redonne l'équation bilan.
\item Avec \ce{HCOOH} (acide) : \ce{HCO3- + HCOOH <=> H2CO3 + HCOO-}, \ce{HCO3-} joue le rôle de base (couple \ce{H2CO3}/\ce{HCO3-}). Avec \ce{HO-} (base) : \ce{HCO3- + HO- <=> CO3^2- + H2O}, \ce{HCO3-} joue le rôle d'acide (couple \ce{HCO3-}/\ce{CO3^2-}). Pouvant jouer les deux rôles, \ce{HCO3-} est un \textbf{ampholyte}.
\end{enumerate}
\end{corrige}

\begin{aretenir}[L'essentiel de la section]
\begin{itemize}
\item Un \cle{couple acide/base} relie un acide \ce{AH} et sa base conjuguée \ce{A-} par la demi-équation \ce{AH <=> A- + H+} ; il se note \textbf{acide/base} : \ce{AH}/\ce{A-}.
\item Pour identifier les couples d'un équilibre : éliminer les ions spectateurs, repérer les paires d'espèces ne différant que d'un seul \ce{H+}, l'acide étant l'espèce la plus protonée.
\item L'eau appartient à deux couples : \ce{H3O+}/\ce{H2O} (eau base) et \ce{H2O}/\ce{HO-} (eau acide) ; c'est un \cle{ampholyte}, comme \ce{HCO3-}, \ce{HSO4-} ou \ce{H2PO4-}.
\item Toute réaction acido-basique est un transfert de proton mettant en jeu \textbf{deux couples} : \ce{acide1 + base2 <=> base1 + acide2}.
\end{itemize}
\end{aretenir}

Maintenant que tu sais reconnaître et écrire les couples, une question naturelle se pose : tous les acides faibles se valent-ils ? L'acide éthanoïque est-il \og plus faible \fg{} ou \og plus fort \fg{} que l'ion ammonium ? Pour répondre, il faut un outil quantitatif capable de mesurer la force d'un couple : ce sera la \cle{constante d'acidité} $K_a$ et son compagnon le $pK_a$, objet de la section suivante.

\coursec{Constantes d'acidité $K_a$, de basicité $K_b$ et $pK_a$}

Dans la section précédente, nous avons appris à reconnaître un couple acide/base et à écrire son demi-équilibre formel $\ce{AH <=> A- + H+}$. Mais savoir qu'un couple existe ne suffit pas : il faut pouvoir \emph{comparer} les couples entre eux. L'acide éthanoïque du vinaigre et l'ion ammonium présent dans certains engrais sont tous deux des acides faibles, mais lequel est le « plus faible » ? Pour répondre, les chimistes ont forgé un outil remarquable : un simple nombre, le $pK_a$, qui résume à lui seul la force d'un couple. C'est la carte d'identité de tout couple acide/base, et c'est l'objet de cette section.

\courssub{Constante d'acidité $K_a$ d'un couple acide/base}

\textbf{De l'intuition à la définition.} Plonge un acide faible \ce{AH} dans l'eau. Il se produit la réaction \emph{limitée} (équilibrée) :

\[ \ce{AH + H2O <=> A- + H3O+} \]

Cette réaction est lente ou rapide peu importe ; l'essentiel est qu'elle conduit à un \cle{équilibre chimique} : les quatre espèces \ce{AH}, \ce{A-}, \ce{H3O+} et \ce{H2O} coexistent en permanence, avec des concentrations qui ne varient plus. Or, d'après la loi d'action de masse, tout équilibre chimique est caractérisé par une constante d'équilibre. Appliquée ici, elle donne naissance à la constante d'acidité.

\begin{definition}[Constante d'acidité $K_a$]
Pour le couple acide/base \ce{AH}/\ce{A-}, la \cle{constante d'acidité} $K_a$ est la constante d'équilibre de la réaction de l'acide avec l'eau :
\[ \ce{AH + H2O <=> A- + H3O+} \qquad K_a = \frac{[\ce{A-}]\,[\ce{H3O+}]}{[\ce{AH}]} \]
L'eau, en tant que solvant, n'apparaît pas dans l'expression (sa concentration est constante et incorporée dans $K_a$). Les concentrations sont celles prises \textbf{à l'équilibre}, exprimées en \si{mol.L^{-1}}. $K_a$ est un nombre sans unité qui \textbf{ne dépend que de la température} (et bien sûr du couple considéré).
\end{definition}

\begin{attention}[$K_a$ ne dépend ni de la concentration, ni de la dilution]
Diluer une solution modifie le pH et le taux d'ionisation $\alpha$, mais \textbf{jamais} la valeur de $K_a$ : à température fixe, $K_a$ est une constante du couple, tout comme la masse molaire est une constante de l'espèce. Seule une variation de température modifie $K_a$.
\end{attention}

\textbf{Lecture physique de $K_a$.} Le numérateur contient les produits de l'ionisation (\ce{A-} et \ce{H3O+}), le dénominateur l'acide non ionisé \ce{AH}. Donc :
\begin{itemize}
\item plus $K_a$ est \textbf{grand}, plus l'équilibre est déplacé vers la droite, plus l'acide est \textbf{ionisé}, donc plus il est \textbf{fort} ;
\item plus $K_a$ est \textbf{petit}, plus l'acide reste sous forme \ce{AH}, donc plus il est \textbf{faible}.
\end{itemize}

\courssub{Le $pK_a$ : une échelle plus commode}

Les valeurs de $K_a$ s'étendent sur des ordres de grandeur immenses (de $10^{10}$ pour les acides forts à $10^{-16}$ pour les acides très faibles). Manier des puissances de dix est fastidieux ; on utilise donc une échelle logarithmique, exactement comme on est passé de $[\ce{H3O+}]$ au pH.

\begin{definition}[Le $pK_a$]
Le \cle{$pK_a$} d'un couple acide/base est défini par :
\[ pK_a = -\log K_a \qquad \text{soit de façon équivalente} \qquad K_a = 10^{-pK_a} \]
Le $pK_a$ est un nombre sans unité, généralement compris entre $-10$ et $+20$ environ.
\end{definition}

\begin{aretenir}[La règle d'or de comparaison]
\begin{itemize}
\item Plus $K_a$ est \textbf{grand} (donc $pK_a$ \textbf{petit}), plus l'\textbf{acide} est fort.
\item Plus $K_a$ est \textbf{petit} (donc $pK_a$ \textbf{grand}), plus la \textbf{base conjuguée} est forte.
\end{itemize}
\end{aretenir}

\begin{attention}[Le piège classique du Bac]
Erreur fréquente : croire qu'un grand $pK_a$ signifie un acide fort. \textbf{C'est exactement l'inverse !} À cause du signe « moins » dans $pK_a = -\log K_a$, un \textbf{petit} $pK_a$ correspond à un \textbf{grand} $K_a$, donc à un acide \textbf{fort}. Retiens l'image : le $pK_a$ mesure la « réticence » de l'acide à céder son proton. Un acide fort cède facilement son proton : sa réticence ($pK_a$) est faible.
\end{attention}

\courssub{Constante de basicité $K_b$ et relation fondamentale $K_a \cdot K_b = K_e$}

Une base faible \ce{A-} réagit aussi avec l'eau, selon l'équilibre :

\[ \ce{A- + H2O <=> AH + HO-} \]

\begin{definition}[Constante de basicité $K_b$]
La \cle{constante de basicité} du couple \ce{AH}/\ce{A-} est la constante d'équilibre de la réaction de la base avec l'eau :
\[ K_b = \frac{[\ce{AH}]\,[\ce{HO-}]}{[\ce{A-}]} \]
On définit aussi $pK_b = -\log K_b$. Plus $K_b$ est grand, plus la base est forte.
\end{definition}

Multiplions maintenant les deux constantes du \emph{même} couple :

\[ K_a \times K_b = \frac{[\ce{A-}][\ce{H3O+}]}{[\ce{AH}]} \times \frac{[\ce{AH}][\ce{HO-}]}{[\ce{A-}]} = [\ce{H3O+}][\ce{HO-}] = K_e \]

\begin{propriete}[Relation fondamentale]
Pour tout couple acide/base \ce{AH}/\ce{A-} :
\[ K_a \times K_b = K_e = 10^{-14} \quad (\text{à } \SI{25}{\celsius}) \qquad \text{et} \qquad pK_a + pK_b = pK_e = 14 \]
\end{propriete}

Cette relation est capitale : elle montre que $K_a$ et $K_b$ d'un même couple ne sont \textbf{pas indépendantes}. Connaître le $pK_a$ d'un couple, c'est connaître automatiquement le $pK_b$ de sa base conjuguée.

\begin{savaistu}[Acide fort = base conjuguée négligeable]
La relation $K_a \cdot K_b = K_e$ explique une phrase célèbre : \emph{« plus un acide est fort, plus sa base conjuguée est faible »}. Prenons l'acide chlorhydrique : le couple \ce{HCl}/\ce{Cl-} a un $K_a$ énorme (environ $10^7$). Sa base conjuguée, l'ion chlorure \ce{Cl-}, possède donc un $K_b = K_e/K_a \approx 10^{-21}$ : c'est une base si faible qu'elle est totalement \textbf{inerte} dans l'eau — on dit qu'elle est \emph{spectatrice}. C'est pourquoi l'eau salée de mer ne présente aucune basicité malgré ses milliards de tonnes d'ions chlorure ! À l'inverse, l'ion éthanoate \ce{CH3COO-}, base conjuguée d'un acide faible ($pK_a = 4{,}8$), est une base faible mais \emph{réelle} ($pK_b = 9{,}2$) : c'est elle qui rend légèrement basique une solution d'éthanoate de sodium.
\end{savaistu}

\courssub{Ordres de grandeur et classification des couples}

Deux couples du programme serviront de références tout au long de l'année :
\begin{itemize}
\item couple acide éthanoïque / ion éthanoate \ce{CH3COOH}/\ce{CH3COO-} : $K_a \approx 1{,}7 \times 10^{-5}$, soit $pK_a \approx 4{,}8$ ;
\item couple ion ammonium / ammoniac \ce{NH4+}/\ce{NH3} : $K_a \approx 6{,}3 \times 10^{-10}$, soit $pK_a \approx 9{,}2$.
\end{itemize}

Comparons : $pK_a(\ce{CH3COOH}) = 4{,}8 < pK_a(\ce{NH4+}) = 9{,}2$, donc l'acide éthanoïque est un acide \textbf{plus fort} que l'ion ammonium ; en conséquence, l'ammoniac \ce{NH3} est une base \textbf{plus forte} que l'ion éthanoate.

L'échelle des $pK_a$ permet de classer \emph{tous} les couples :

\begin{propriete}[Classification selon le $pK_a$ (à \SI{25}{\celsius})]
\begin{itemize}
\item $pK_a < 0$ : l'acide est \textbf{fort} (totalement ionisé dans l'eau, plus fort que l'ion oxonium) ; sa base conjuguée est inerte (ex. \ce{HCl}/\ce{Cl-}, $pK_a \approx -7$) ;
\item $0 < pK_a < 14$ : acide \textbf{faible} et base conjuguée \textbf{faible} (ex. \ce{CH3COOH}/\ce{CH3COO-}, \ce{NH4+}/\ce{NH3}) ;
\item $pK_a > 14$ : la base est \textbf{forte} (totalement ionisée dans l'eau, plus forte que l'ion hydroxyde) ; l'acide conjugué est inerte (ex. \ce{NH2-}/\ce{NH3}, $pK_a \approx 35$).
\end{itemize}
Les couples limites du solvent eau sont \ce{H3O+}/\ce{H2O} ($pK_a = 0$) et \ce{H2O}/\ce{HO-} ($pK_a = 14$).
\textbf{Force relative à concentration égale :} l'acide le plus fort donne la solution de \textbf{plus petit pH} ; la base la plus forte donne la solution de \textbf{plus grand pH}.
\end{propriete}

\begin{popfigure}
\begin{tikzpicture}
\begin{axis}[
  width=\linewidth, height=5.2cm,
  axis x line=bottom, axis y line=none,
  xmin=-2, xmax=16, ymin=-1.6, ymax=1.6,
  xlabel={$pK_a$}, xlabel style={font=\small},
  xtick={-2,0,2,4,6,8,10,12,14,16},
  xticklabel style={font=\small},
  ytick=\empty,
  every axis x label/.style={at={(ticklabel* cs:1.02)},anchor=west},
]
% Zones
\fill[popblueL]  (axis cs:-2,-0.35) rectangle (axis cs:0,0.35);
\fill[popgreenL] (axis cs:0,-0.35)  rectangle (axis cs:14,0.35);
\fill[poporangeL](axis cs:14,-0.35) rectangle (axis cs:16,0.35);
\node[font=\scriptsize\bfseries, text=popblue]  at (axis cs:-1,0.75) {acides forts};
\node[font=\scriptsize\bfseries, text=popgreen] at (axis cs:7,0.75) {acides faibles / bases faibles};
\node[font=\scriptsize\bfseries, text=poporange] at (axis cs:15,0.75) {bases fortes};
% Couples
\addplot[only marks, mark=*, mark size=2.2pt, color=popblue] coordinates {(-7,0)};
\node[font=\scriptsize, text=popdark, anchor=south] at (axis cs:-1.2,-0.15) {\ce{HCl}/\ce{Cl-} ($-7$)};
\addplot[only marks, mark=*, mark size=2.2pt, color=popblue] coordinates {(0,0)};
\node[font=\scriptsize, text=popdark, anchor=north west, xshift=1pt] at (axis cs:0.1,-0.15) {\ce{H3O+}/\ce{H2O} ($0$)};
\addplot[only marks, mark=*, mark size=2.2pt, color=poppurple] coordinates {(4.8,0)};
\node[font=\scriptsize, text=popdark, anchor=south] at (axis cs:4.8,0.1) {\ce{CH3COOH}/\ce{CH3COO-} ($4{,}8$)};
\addplot[only marks, mark=*, mark size=2.2pt, color=poppurple] coordinates {(9.2,0)};
\node[font=\scriptsize, text=popdark, anchor=south] at (axis cs:9.2,0.1) {\ce{NH4+}/\ce{NH3} ($9{,}2$)};
\addplot[only marks, mark=*, mark size=2.2pt, color=poporange] coordinates {(14,0)};
\node[font=\scriptsize, text=popdark, anchor=north east, xshift=-1pt] at (axis cs:13.9,-0.15) {\ce{H2O}/\ce{HO-} ($14$)};
% Flèches de force
\draw[-{Stealth[length=2.5mm]}, thick, popblue] (axis cs:1,-1.0) -- (axis cs:-1.5,-1.0) node[midway, below, font=\scriptsize, text=popblue]{acide de plus en plus fort};
\draw[-{Stealth[length=2.5mm]}, thick, poporange] (axis cs:12,-1.0) -- (axis cs:15.5,-1.0) node[midway, below, font=\scriptsize, text=poporange]{base de plus en plus forte};
\end{axis}
\end{tikzpicture}
\legende{Échelle des $pK_a$ à \SI{25}{\celsius} et positionnement de quelques couples usuels du programme.}
\end{popfigure}

Le tableau suivant donne la correspondance entre $K_a$ et $pK_a$ pour les couples rencontrés au lycée :

\begin{center}
\begin{tabularx}{\linewidth}{|l|X|c|c|}
\hline
\rowcolor{popblueL}
\textbf{Couple acide/base} & \textbf{Acide (forme)} & $\boldsymbol{K_a}$ & $\boldsymbol{pK_a}$ \\
\hline
\ce{HCl}/\ce{Cl-} & acide chlorhydrique (fort) & $\approx 10^{7}$ & $\approx -7$ \\
\hline
\ce{H3O+}/\ce{H2O} & ion oxonium & $1$ & $0$ \\
\hline
\ce{CH3COOH}/\ce{CH3COO-} & acide éthanoïque (vinaigre) & $1{,}7\times 10^{-5}$ & $4{,}8$ \\
\hline
\ce{NH4+}/\ce{NH3} & ion ammonium & $6{,}3\times 10^{-10}$ & $9{,}2$ \\
\hline
\ce{H2O}/\ce{HO-} & eau (acide très faible) & $10^{-14}$ & $14$ \\
\hline
\end{tabularx}
\end{center}

\begin{attention}[Bien lire un tableau de couples]
Dans un tableau de couples classés par $pK_a$ croissant, les acides les plus forts sont \textbf{en haut} (petit $pK_a$) et les bases les plus fortes \textbf{en bas} (grand $pK_a$). Ne confonds pas la force d'un acide avec celle de sa base conjuguée : elles varient \emph{en sens inverse}.
\end{attention}

\courssub{Détermination expérimentale de la constante d'acidité $K_a$}

Comment le chimiste obtient-il la valeur $pK_a = 4{,}8$ de l'acide éthanoïque ? Par une simple mesure de pH ! L'idée est élégante : dans une solution d'acide faible seul, chaque ion \ce{H3O+} formé s'accompagne d'un ion \ce{A-} (stœchiométrie 1:1 de l'équilibre d'ionisation). La mesure du pH donne donc directement deux concentrations sur trois.

\begin{methode}[Déterminer expérimentalement le $K_a$ d'un couple]
\begin{enumerate}
\item Préparer une solution de l'acide \ce{AH} de concentration apportée $C$ connue (ex. $C = \SI{1,0e-2}{mol.L^{-1}}$).
\item Mesurer le pH de la solution au pH-mètre préalablement étalonné.
\item En déduire $[\ce{H3O+}] = 10^{-pH}$.
\item Exploiter la stœchiométrie de $\ce{AH + H2O <=> A- + H3O+}$ : $[\ce{A-}] = [\ce{H3O+}]$ (les ions \ce{HO-} de l'eau sont négligeables en milieu acide).
\item Par conservation : $[\ce{AH}] = C - [\ce{H3O+}]$.
\item Appliquer la formule : $K_a = \dfrac{[\ce{H3O+}]^2}{C - [\ce{H3O+}]}$ puis $pK_a = -\log K_a$.
\item Vérifier la cohérence : $[\ce{H3O+}]$ doit être nettement inférieure à $C$ (acide faible) ; sinon l'approximation est à revoir.
\end{enumerate}
\end{methode}

\begin{exemplebox}[Le $pK_a$ de l'acide éthanoïque mesuré au laboratoire]
Une solution d'acide éthanoïque de concentration $C = \SI{1,0e-2}{mol.L^{-1}}$ a un pH mesuré de $3{,}4$ à \SI{25}{\celsius}. Alors :
\begin{align*}
[\ce{H3O+}] &= 10^{-3{,}4} = 4{,}0\times 10^{-4}\ \si{mol.L^{-1}} = [\ce{CH3COO-}] \\
[\ce{CH3COOH}] &= C - [\ce{H3O+}] = 10^{-2} - 4{,}0\times 10^{-4} = 9{,}6\times 10^{-3}\ \si{mol.L^{-1}} \\
K_a &= \frac{(10^{-3{,}4})^2}{10^{-2} - 10^{-3{,}4}} = \frac{(4{,}0\times 10^{-4})^2}{9{,}6\times 10^{-3}} \approx 1{,}7\times 10^{-5}
\end{align*}
d'où $pK_a = -\log(1{,}7\times 10^{-5}) \approx 4{,}8$. On retrouve la valeur des tables ! Au passage, le coefficient d'ionisation vaut $\alpha = [\ce{H3O+}]/C = 4{,}0\times 10^{-4}/10^{-2} = 0{,}04$, soit seulement \SI{4}{\percent} : l'acide éthanoïque est bien un acide faible.
\end{exemplebox}

\begin{exoresolu}[Du pH au $pK_a$ : le couple de l'acide méthanoïque]
L'acide méthanoïque \ce{HCOOH} (acide des fourmis) donne, à la concentration $C = \SI{2,0e-2}{mol.L^{-1}}$, une solution de $\text{pH} = 2{,}7$ à \SI{25}{\celsius}.\\
1. Écrire l'équation de la réaction de cet acide avec l'eau.\\
2. Calculer $K_a$ et $pK_a$ du couple \ce{HCOOH}/\ce{HCOO-}.\\
3. En déduire $K_b$ et $pK_b$ de l'ion méthanoate. Comparer la force de \ce{HCOOH} à celle de \ce{CH3COOH} ($pK_a = 4{,}8$).
\tcblower
\textbf{1.} Équilibre d'ionisation (réaction limitée, à \SI{25}{\celsius}) :
\[ \ce{HCOOH + H2O <=> HCOO- + H3O+} \]
\textbf{2.} On calcule successivement :
\begin{align*}
[\ce{H3O+}] &= 10^{-2{,}7} = 2{,}0\times 10^{-3}\ \si{mol.L^{-1}} = [\ce{HCOO-}] \\
[\ce{HCOOH}] &= C - [\ce{H3O+}] = 2{,}0\times 10^{-2} - 2{,}0\times 10^{-3} = 1{,}8\times 10^{-2}\ \si{mol.L^{-1}} \\
K_a &= \frac{(2{,}0\times 10^{-3})^2}{1{,}8\times 10^{-2}} \approx 2{,}2\times 10^{-4}
\qquad pK_a = -\log(2{,}2\times 10^{-4}) \approx 3{,}7
\end{align*}
\textbf{3.} D'après la relation fondamentale :
\[ K_b = \frac{K_e}{K_a} = \frac{10^{-14}}{2{,}2\times 10^{-4}} \approx 4{,}5\times 10^{-11} \qquad pK_b = 14 - 3{,}7 = 10{,}3 \]
Comparaison : $pK_a(\ce{HCOOH}) = 3{,}7 < pK_a(\ce{CH3COOH}) = 4{,}8$, donc l'acide méthanoïque est \textbf{plus fort} que l'acide éthanoïque (il s'ionise davantage : ici $\alpha = \SI{10}{\percent}$ contre \SI{4}{\percent} pour l'éthanoïque à même concentration). Réciproquement, l'ion méthanoate est une base plus faible que l'ion éthanoate.
\end{exoresolu}

\begin{exercice}[$pK_a$ et dilution]
On mesure le pH de deux solutions d'acide benzoïque \ce{C6H5COOH} à \SI{25}{\celsius} : solution (1) $C_1 = \SI{1,0e-1}{mol.L^{-1}}$, $\text{pH}_1 = 2{,}95$ ; solution (2) $C_2 = \SI{1,0e-3}{mol.L^{-1}}$, $\text{pH}_2 = 3{,}9$.\\
1. Calculer $K_a$ et $pK_a$ à partir de chaque solution. Que constates-tu ?\\
2. Calculer le coefficient d'ionisation $\alpha$ dans chaque cas. Quelle loi illustre ce résultat ?
\end{exercice}

\begin{corrige}[Exercice 1]
\textbf{1.} Solution (1) : $[\ce{H3O+}] = 10^{-2{,}95} = 1{,}12\times 10^{-3}$ ; $K_a = \dfrac{(1{,}12\times 10^{-3})^2}{10^{-1} - 1{,}12\times 10^{-3}} \approx 1{,}27\times 10^{-5}$.\\
Solution (2) : $[\ce{H3O+}] = 10^{-3{,}9} = 1{,}26\times 10^{-4}$ ; $K_a = \dfrac{(1{,}26\times 10^{-4})^2}{10^{-3} - 1{,}26\times 10^{-4}} \approx 1{,}8\times 10^{-5}$.\\
On constate que $K_a$ est (quasi) identique dans les deux cas, soit $pK_a \approx 4{,}8$--$4{,}9$ : $K_a$ ne dépend pas de la concentration, c'est bien une \textbf{constante} du couple à température donnée.\\
\textbf{2.} $\alpha_1 = \dfrac{1{,}12\times 10^{-3}}{10^{-1}} \approx 1{,}1\,\%$ ; $\alpha_2 = \dfrac{1{,}26\times 10^{-4}}{10^{-3}} \approx 12{,}6\,\%$. La dilution augmente l'ionisation : c'est la \textbf{loi de dilution d'Ostwald}.
\end{corrige}

\begin{aretenir}[L'essentiel de la section]
\begin{itemize}
\item $K_a = \dfrac{[\ce{A-}][\ce{H3O+}]}{[\ce{AH}]}$ : constante d'acidité du couple \ce{AH}/\ce{A-}, sans unité, ne dépendant que de la température.
\item $pK_a = -\log K_a$ ; petit $pK_a$ = acide fort ; grand $pK_a$ = base conjuguée forte.
\item $K_a \times K_b = K_e = 10^{-14}$ et $pK_a + pK_b = 14$ à \SI{25}{\celsius}.
\item Classification : $pK_a < 0$ acide fort ; $0 < pK_a < 14$ acide faible ; $pK_a > 14$ base forte.
\item Détermination expérimentale : mesurer le pH d'une solution de concentration $C$ connue, puis $K_a = \dfrac{(10^{-pH})^2}{C - 10^{-pH}}$.
\end{itemize}
\end{aretenir}

Maintenant que chaque couple possède son « étiquette » $pK_a$, une question naturelle se pose : dans une solution dont on fixe le pH, quelle forme du couple domine, l'acide \ce{AH} ou la base \ce{A-} ? C'est la relation $pH = pK_a + \log\frac{[\ce{A-}]}{[\ce{AH}]}$ qui nous le dira — et elle expliquera au passage le changement de couleur des indicateurs colorés. Rendez-vous à la section suivante.

\coursec{Relation pH-pKa et domaines de prédominance}

Dans la section précédente, nous avons appris à quantifier la force d'un couple acide/base grâce à sa constante d'acidité $K_a$ et à son $pK_a$. Mais une question pratique demeure : lorsque l'on dissout de l'acide éthanoïque dans de l'eau, ou lorsque le pH d'une solution varie au cours d'un dosage, \emph{quelle forme du couple est majoritaire ?} Est-ce l'acide \ce{CH3COOH} ou sa base conjuguée \ce{CH3COO-} ? La réponse à cette question est capitale : c'est elle qui explique le changement de couleur des indicateurs colorés, le pouvoir tampon du sang, et même la conservation des aliments. Toute cette section repose sur une seule relation, simple et puissante, que nous allons d'abord établir rigoureusement.

\courssub{Établissement de la relation pH--pKa}

Considérons un couple acide/base quelconque \ce{AH}/\ce{A-}. L'équilibre d'ionisation de l'acide dans l'eau s'écrit :

\[ \ce{AH + H2O <=> A- + H3O+} \]

Cet équilibre est \cle{rapide}, \cle{partiel} et se réalise à température ambiante sans catalyseur. Sa constante d'acidité s'écrit :

\[ K_a = \frac{[\ce{A-}] \times [\ce{H3O+}]}{[\ce{AH}]} \]

Prenons le logarithme décimal (en changeant le signe) de cette expression. Rappelons que $\log(a \times b) = \log a + \log b$ et que $-\log[\ce{H3O+}] = \mathrm{pH}$ :

\begin{align*}
K_a &= \frac{[\ce{A-}]\,[\ce{H3O+}]}{[\ce{AH}]}\\
-\log K_a &= -\log[\ce{H3O+}] - \log\frac{[\ce{A-}]}{[\ce{AH}]}\\
pK_a &= \mathrm{pH} - \log\frac{[\ce{A-}]}{[\ce{AH}]}
\end{align*}

En isolant le pH, on obtient la relation fondamentale de cette section, dite \emph{relation de Henderson-Hasselbalch} :

\begin{definition}[Relation pH--pKa (Henderson-Hasselbalch)]
Pour tout couple acide/base \ce{AH}/\ce{A-} présent en solution aqueuse, le pH de la solution est relié au $pK_a$ du couple par :
\[ \boxed{\; \mathrm{pH} = pK_a + \log\frac{[\ce{A-}]}{[\ce{AH}]} \;} \]
où $[\ce{A-}]$ est la concentration de la forme \textbf{basique} (au numérateur) et $[\ce{AH}]$ celle de la forme \textbf{acide} (au dénominateur). Cette relation est valable dès que l'équilibre $\ce{AH + H2O <=> A- + H3O+}$ est établi.
\end{definition}

\begin{attention}[Le piège du rapport inversé]
L'erreur la plus fréquente au Baccalauréat est d'écrire $\log\dfrac{[\ce{AH}]}{[\ce{A-}]}$ au lieu de $\log\dfrac{[\ce{A-}]}{[\ce{AH}]}$. Retiens ce moyen mnémotechnique : \textbf{la Base est en haut} (au numérateur), comme dans « B » de \emph{Bon élève}. Vérification rapide : si la base prédomine, le rapport est supérieur à 1, son logarithme est positif, donc $\mathrm{pH} > pK_a$ : c'est cohérent, un milieu riche en base a un pH élevé. Si tu obtiens une contradiction de ce type, c'est que tu as inversé le rapport.
\end{attention}

\courssub{Comparaison du pH au pKa : les trois cas}

Analysons la relation de Henderson-Hasselbalch dans trois situations.

\textbf{Premier cas : $\mathrm{pH} = pK_a$.} Alors $\log\dfrac{[\ce{A-}]}{[\ce{AH}]} = 0$, donc $\dfrac{[\ce{A-}]}{[\ce{AH}]} = 1$ : les deux formes conjuguées sont en \cle{quantités égales} : $[\ce{AH}] = [\ce{A-}]$. C'est la situation d'une solution tampon idéale, à mi-chemin du dosage d'un acide faible par une base forte (la fameuse demi-équivalence).

\textbf{Deuxième cas : $\mathrm{pH} < pK_a$.} Le logarithme est négatif, donc le rapport est inférieur à 1 : la forme \cle{acide} \ce{AH} prédomine. Plus précisément, si $\mathrm{pH} < pK_a - 1$, alors $\dfrac{[\ce{A-}]}{[\ce{AH}]} < 10^{-1}$, c'est-à-dire $[\ce{AH}] > 10\,[\ce{A-}]$ : la forme acide est \emph{ultra-majoritaire} (plus de 90 \% de l'espèce dissoute).

\textbf{Troisième cas : $\mathrm{pH} > pK_a$.} Le logarithme est positif, donc le rapport est supérieur à 1 : la forme \cle{basique} \ce{A-} prédomine. Si $\mathrm{pH} > pK_a + 1$, alors $[\ce{A-}] > 10\,[\ce{AH}]$ : la forme basique est ultra-majoritaire.

\begin{methode}[Quelle forme prédomine ?]
Pour déterminer l'espèce prédominante d'un couple \ce{AH}/\ce{A-} dans une solution de pH connu :
\begin{enumerate}
\item Relever le $pK_a$ du couple (donné ou calculé à partir de $K_a$ : $pK_a = -\log K_a$).
\item Comparer le pH de la solution au $pK_a$.
\item Conclure : si $\mathrm{pH} < pK_a$, la forme \textbf{acide} \ce{AH} prédomine ; si $\mathrm{pH} > pK_a$, la forme \textbf{basique} \ce{A-} prédomine ; si $\mathrm{pH} = pK_a$, les deux formes sont en concentrations égales.
\item Pour préciser les proportions, calculer le rapport $\dfrac{[\ce{A-}]}{[\ce{AH}]} = 10^{\mathrm{pH} - pK_a}$.
\end{enumerate}
\end{methode}

\begin{exemplebox}[L'acide éthanoïque à pH = 6,8]
Le couple de l'acide éthanoïque est \ce{CH3COOH}/\ce{CH3COO-}, de $pK_a = 4,8$. Une solution contenant ce couple est portée à $\mathrm{pH} = 6,8$ (par exemple en ajoutant de la soude). Quelle forme prédomine ?

On calcule : $\dfrac{[\ce{CH3COO-}]}{[\ce{CH3COOH}]} = 10^{\mathrm{pH} - pK_a} = 10^{6,8 - 4,8} = 10^{2} = 100$.

Il y a donc \textbf{100 fois plus} d'ions éthanoate que de molécules d'acide éthanoïque : l'ion \ce{CH3COO-} prédomine largement (environ 99 \% de l'espèce dissoute). Comme $\mathrm{pH} = pK_a + 2 > pK_a + 1$, on est bien dans le domaine où la base est ultra-majoritaire. C'est exactement ce qui se passe lorsqu'on ajoute de la soude au vinaigre : l'acide est progressivement « converti » en ion éthanoate.
\end{exemplebox}

\courssub{Diagramme de prédominance}

Le résultat précédent se visualise sur un \cle{diagramme de prédominance} : un axe de pH partagé en deux domaines séparés par la valeur du $pK_a$. À gauche (milieu acide), la forme acide règne ; à droite (milieu basique), c'est la forme basique.

\begin{popfigure}
\begin{center}
\begin{tikzpicture}[x=0.85cm]
% Axe de pH
\draw[->, thick, popdark] (0,0) -- (11,0) node[below left=2pt and -30pt] {\small pH};
\foreach \x/\lab in {0/0, 2/2, 4/4, 4.8/4{,}8, 6/6, 8/8, 10/10}{
  \draw[popdark] (\x,0.08) -- (\x,-0.08);
  \node[below, popdark, font=\small] at (\x,-0.1) {\lab};
}
% Frontière pKa
\draw[very thick, popink] (4.8,-0.55) -- (4.8,1.15);
\node[above, popink, font=\small\bfseries] at (4.8,1.15) {$pK_a = 4{,}8$};
% Domaine acide
\fill[popblue!20] (0,0.35) rectangle (4.8,0.95);
\node[popblue, font=\small\bfseries] at (2.4,0.65) {\ce{CH3COOH} prédomine};
\node[popblue, font=\footnotesize] at (2.4,0.15) {$[\ce{CH3COOH}] > [\ce{CH3COO-}]$};
% Domaine basique
\fill[poporange!25] (4.8,0.35) rectangle (10.6,0.95);
\node[poporange!80!black, font=\small\bfseries] at (7.7,0.65) {\ce{CH3COO-} prédomine};
\node[poporange!80!black, font=\footnotesize] at (7.7,0.15) {$[\ce{CH3COO-}] > [\ce{CH3COOH}]$};
% Zone ultra-majoritaire
\draw[<->, thick, popblue] (0,-0.75) -- (3.8,-0.75) node[midway, below, font=\footnotesize] {AH ultra-majoritaire ($\mathrm{pH} < pK_a - 1$)};
\draw[<->, thick, poporange!80!black] (5.8,-0.75) -- (10.6,-0.75) node[midway, below, font=\footnotesize] {\ce{A-} ultra-majoritaire ($\mathrm{pH} > pK_a + 1$)};
\end{tikzpicture}
\end{center}
\legende{Diagramme de prédominance du couple \ce{CH3COOH}/\ce{CH3COO-} ($pK_a = 4{,}8$). La frontière verticale correspond à $\mathrm{pH} = pK_a$, où les deux formes sont en concentrations égales.}
\end{popfigure}

\courssub{Diagramme de distribution}

Le diagramme de prédominance dit \emph{qui} prédomine, mais pas \emph{dans quelle proportion}. Pour cela, on trace le \cle{diagramme de distribution} : les pourcentages de chaque forme en fonction du pH. En notant $C = [\ce{AH}] + [\ce{A-}]$ la concentration totale, un calcul à partir de la relation de Henderson-Hasselbalch donne :

\[ \%(\ce{AH}) = \frac{100}{1 + 10^{\mathrm{pH} - pK_a}} \qquad \%(\ce{A-}) = \frac{100 \times 10^{\mathrm{pH} - pK_a}}{1 + 10^{\mathrm{pH} - pK_a}} \]

\begin{popfigure}
\begin{center}
\begin{tikzpicture}
\begin{axis}[
  width=0.92\linewidth, height=6.2cm,
  xlabel={pH}, ylabel={Pourcentage (\%)},
  xmin=0, xmax=10, ymin=0, ymax=105,
  xtick={0,2,4,4.8,6,8,10}, ytick={0,25,50,75,100},
  grid=both, grid style={popline!40},
  legend style={at={(0.5,-0.28)}, anchor=north, legend columns=2, font=\small},
  axis line style={popdark}, label style={popdark}, tick label style={popdark, font=\small}
]
\addplot[popcurve, popblue, domain=0:10, samples=200] {100/(1+10^(x-4.8))};
\addlegendentry{\% \ce{CH3COOH} (forme acide)}
\addplot[popcurve, poporange, domain=0:10, samples=200] {100*10^(x-4.8)/(1+10^(x-4.8))};
\addlegendentry{\% \ce{CH3COO-} (forme basique)}
\draw[dashed, popink, thick] (axis cs:4.8,0) -- (axis cs:4.8,105);
\node[popink, font=\small, rotate=90, anchor=south] at (axis cs:4.55,55) {$\mathrm{pH} = pK_a = 4{,}8$};
\addplot[only marks, mark=*, popink, mark size=2pt] coordinates {(4.8,50)};
\end{axis}
\end{tikzpicture}
\end{center}
\legende{Diagramme de distribution du couple \ce{CH3COOH}/\ce{CH3COO-}. Les deux courbes se croisent à $\mathrm{pH} = pK_a$, où chaque forme représente exactement 50 \% de l'espèce dissoute.}
\end{popfigure}

Les deux courbes, symétriques, se croisent obligatoirement au point $(pK_a \,;\, 50\,\%)$. On y lit directement, par exemple, qu'à $\mathrm{pH} = 2,8$ l'acide représente environ 99 \% de l'espèce dissoute, tandis qu'à $\mathrm{pH} = 6,8$ c'est l'ion éthanoate qui atteint 99 \%.

\courssub{Application : les indicateurs colorés}

\begin{definition}[Indicateur coloré acido-basique]
Un \cle{indicateur coloré} acido-basique est un couple acide/base \ce{HIn}/\ce{In-} dont la forme acide \ce{HIn} et la forme basique \ce{In-} ont des \textbf{couleurs différentes}. L'équilibre en solution s'écrit :
\[ \ce{HIn + H2O <=> In- + H3O+} \qquad \mathrm{pH} = pK_{a,\mathrm{In}} + \log\frac{[\ce{In-}]}{[\ce{HIn}]} \]
(équilibre rapide, partiel, à température ambiante). La couleur observée dépend donc du pH de la solution.
\end{definition}

D'après la relation pH--pKa appliquée au couple de l'indicateur :
\begin{itemize}
\item si $\mathrm{pH} < pK_{a,\mathrm{In}} - 1$, la forme acide \ce{HIn} est ultra-majoritaire : on observe la \emph{teinte acide} ;
\item si $\mathrm{pH} > pK_{a,\mathrm{In}} + 1$, la forme basique \ce{In-} est ultra-majoritaire : on observe la \emph{teinte basique} ;
\item si $pK_{a,\mathrm{In}} - 1 < \mathrm{pH} < pK_{a,\mathrm{In}} + 1$, les deux formes coexistent en proportions notables : on observe une teinte intermédiaire, dite \emph{teinte sensible}. Cet intervalle est la \cle{zone de virage}.
\end{itemize}

\begin{aretenir}[Zone de virage]
La \textbf{zone de virage} d'un indicateur coloré est l'intervalle de pH, centré sur son $pK_a$, dans lequel il change progressivement de couleur :
\[ pK_{a,\mathrm{In}} - 1 < \mathrm{pH} < pK_{a,\mathrm{In}} + 1 \]
Exemples usuels : hélianthine ($pK_a \approx 3,7$ ; rouge en milieu acide, jaune-orange en milieu basique), bleu de bromothymol ou BBT ($pK_a \approx 7$ ; jaune / bleu, vert à la teinte sensible), phénolphtaléine ($pK_a \approx 9$ ; incolore / rose fuchsia).
\end{aretenir}

\begin{popfigure}
\begin{center}
\begin{tikzpicture}[x=0.72cm, y=0.8cm]
% Echelle de pH
\draw[->, thick, popdark] (0,0) -- (15,0);
\foreach \p in {0,1,...,14}{
  \draw[popdark] (\p,0.06) -- (\p,-0.06);
  \node[below, popdark, font=\footnotesize] at (\p,-0.08) {\p};
}
\node[below, popdark, font=\small] at (15.3,-0.1) {pH};
% Hélianthine : virage 2,7 - 4,7 (pKa 3,7)
\fill[red!70!black] (0,0.5) rectangle (2.7,1.0);
\fill[poporange!50!red] (2.7,0.5) rectangle (4.7,1.0);
\fill[orange!80!yellow] (4.7,0.5) rectangle (14.5,1.0);
\node[left, font=\small\bfseries, popdark] at (0,0.75) {Hélianthine};
\node[font=\footnotesize, white] at (1.35,0.75) {rouge};
\node[font=\footnotesize, popdark] at (3.7,0.75) {orange};
\node[font=\footnotesize, popdark] at (9,0.75) {jaune-orange};
% BBT : virage 6 - 8 (pKa 7)
\fill[yellow!85!orange] (0,1.4) rectangle (6,1.9);
\fill[popgreen!50!popblue] (6,1.4) rectangle (8,1.9);
\fill[popblue!80!black] (8,1.4) rectangle (14.5,1.9);
\node[left, font=\small\bfseries, popdark] at (0,1.65) {BBT};
\node[font=\footnotesize, popdark] at (3,1.65) {jaune};
\node[font=\footnotesize, white] at (7,1.65) {vert};
\node[font=\footnotesize, white] at (11,1.65) {bleu};
% Phénolphtaléine : virage 8 - 10 (pKa 9)
\fill[popcream] (0,2.3) rectangle (8,2.8);
\fill[poppink!50!popcream] (8,2.3) rectangle (10,2.8);
\fill[poppink!85!black] (10,2.3) rectangle (14.5,2.8);
\node[left, font=\small\bfseries, popdark] at (0,2.55) {Phénolphtaléine};
\node[font=\footnotesize, popmuted] at (4,2.55) {incolore};
\node[font=\footnotesize, popdark] at (9,2.55) {rose pâle};
\node[font=\footnotesize, white] at (12,2.55) {rose fuchsia};
% Marqueurs pKa
\foreach \p/\y in {3.7/0.75, 7/1.65, 9/2.55}{
  \draw[dashed, popink, thick] (\p,\y-0.25) -- (\p,\y+0.28);
}
\end{tikzpicture}
\end{center}
\legende{Bandeau des indicateurs colorés usuels. Les zones colorées intermédiaires correspondent aux zones de virage ($pK_a \pm 1$) ; les pointillés marquent le $pK_a$ de chaque indicateur.}
\end{popfigure}

\begin{propriete}[Choix d'un indicateur pour un dosage]
Pour repérer l'équivalence d'un dosage acido-basique par changement de couleur, on choisit un indicateur dont la \textbf{zone de virage encadre le pH à l'équivalence} $\mathrm{pH}_E$ :
\[ pK_{a,\mathrm{In}} - 1 < \mathrm{pH}_E < pK_{a,\mathrm{In}} + 1 \]
Exemples : dosage d'un acide fort par une base forte ($\mathrm{pH}_E = 7$) : le BBT convient parfaitement ; dosage d'un acide faible par la soude ($\mathrm{pH}_E \approx 8,7$) : on choisit la phénolphtaléine ; dosage d'une base faible par un acide fort ($\mathrm{pH}_E \approx 5$) : l'hélianthine est adaptée.
\end{propriete}

\begin{savaistu}[Le chou rouge, indicateur du village]
Le jus de chou rouge contient des \emph{anthocyanes}, pigments naturels qui changent de couleur selon le pH : rouge en milieu acide (jus de citron, vinaigre), violet vers la neutralité, puis vert voire jaune en milieu basique (eau de lessive de cendres, solution de « potasse » artisanale). On peut donc réaliser en classe, avec du matériel du quotidien, un véritable indicateur coloré naturel pour comparer l'acidité du jus de bissap, du vin de palme frais ou d'une eau de source. Le jus de bissap (fleurs d'hibiscus), naturellement rouge acide, vire au bleu-vert en milieu basique : la nature offre ses propres indicateurs !
\end{savaistu}

\begin{exoresolu}[Prédominance et couleur d'un indicateur]
\textbf{Énoncé.} On verse quelques gouttes de bleu de bromothymol (BBT, $pK_{a,\mathrm{In}} = 7,0$) dans une solution aqueuse S de $\mathrm{pH} = 5,0$.
\begin{enumerate}
\item Écrire l'équation de l'équilibre acido-basique du BBT dans l'eau.
\item Calculer le rapport $\dfrac{[\ce{In-}]}{[\ce{HIn}]}$ dans la solution S.
\item Quelle est la couleur observée ? Justifier.
\end{enumerate}
\tcblower
\textbf{Solution détaillée.}
\begin{enumerate}
\item Le BBT est un couple acide/base \ce{HIn}/\ce{In-}. Son équilibre avec l'eau, rapide et partiel, s'écrit :
\[ \ce{HIn + H2O <=> In- + H3O+} \]
\item La relation de Henderson-Hasselbalch donne :
\[ \frac{[\ce{In-}]}{[\ce{HIn}]} = 10^{\mathrm{pH} - pK_{a,\mathrm{In}}} = 10^{5,0 - 7,0} = 10^{-2} = 0,01 \]
Il y a donc 100 fois plus de forme acide \ce{HIn} que de forme basique \ce{In-}.
\item Comme $\mathrm{pH} = 5,0 < pK_{a,\mathrm{In}} - 1 = 6,0$, la forme acide est ultra-majoritaire : la solution prend la \textbf{teinte acide du BBT, c'est-à-dire le jaune}. Le point de vigilance : ne pas inverser le rapport dans le logarithme ; la base \ce{In-} est toujours au numérateur.
\end{enumerate}
\end{exoresolu}

\begin{exercice}[À toi de jouer]
Le couple de l'acide hypochloreux, principe actif de l'eau de Javel, est \ce{HClO}/\ce{ClO-} de $pK_a = 7,5$.
\begin{enumerate}
\item Tracer le diagramme de prédominance de ce couple.
\item L'eau de Javel commerciale a un pH voisin de 11. Quelle forme du couple y prédomine ? Calculer le rapport des concentrations.
\item On acidifie de l'eau de Javel jusqu'à $\mathrm{pH} = 4$. Quelle forme devient majoritaire ? Pourquoi cette opération est-elle dangereuse en pratique ?
\end{enumerate}
\end{exercice}

\begin{corrige}[Exercice 1]
\begin{enumerate}
\item Axe de pH coupé en $\mathrm{pH} = 7,5$ : à gauche ($\mathrm{pH} < 7,5$), \ce{HClO} prédomine ; à droite ($\mathrm{pH} > 7,5$), \ce{ClO-} prédomine ; à $\mathrm{pH} = 7,5$, $[\ce{HClO}] = [\ce{ClO-}]$.
\item À $\mathrm{pH} = 11$ : $\dfrac{[\ce{ClO-}]}{[\ce{HClO}]} = 10^{11 - 7,5} = 10^{3,5} \approx 3{,}2 \times 10^{3}$. L'ion hypochlorite \ce{ClO-} est ultra-majoritaire ($\mathrm{pH} > pK_a + 1$).
\item À $\mathrm{pH} = 4 < pK_a - 1 = 6,5$, c'est \ce{HClO} qui devient ultra-majoritaire. En milieu acide, l'ion hypochlorite \ce{ClO-} réagit avec les ions chlorure \ce{Cl-} (présents dans l'eau de Javel commerciale) pour former du dichlore \ce{Cl2}, gaz très toxique : \ce{ClO- + Cl- + 2H+ -> Cl2 + H2O}. C'est pourquoi il ne faut \textbf{jamais} mélanger l'eau de Javel avec un produit acide (détartrant, vinaigre, acide chlorhydrique).
\end{enumerate}
\end{corrige}

\begin{aretenir}[L'essentiel de la section]
\begin{itemize}
\item Relation fondamentale : $\mathrm{pH} = pK_a + \log\dfrac{[\ce{A-}]}{[\ce{AH}]}$, la \textbf{base au numérateur}.
\item $\mathrm{pH} < pK_a$ : la forme acide prédomine ; $\mathrm{pH} > pK_a$ : la forme basique prédomine ; $\mathrm{pH} = pK_a$ : égalité des concentrations.
\item Le diagramme de prédominance partage l'axe des pH en deux domaines séparés par le $pK_a$ ; le diagramme de distribution donne les pourcentages, avec croisement des courbes à 50 \% pour $\mathrm{pH} = pK_a$.
\item Un indicateur coloré est un couple \ce{HIn}/\ce{In-} aux couleurs différentes ; sa zone de virage est $pK_a \pm 1$ ; pour un dosage, elle doit contenir le pH à l'équivalence.
\end{itemize}
\end{aretenir}

\coursec{Prévision des réactions acido-basiques}

Dans la section précédente, nous avons appris à lire les domaines de prédominance : à un pH donné, on sait quelle forme (acide ou base) d'un couple domine. Il nous reste une question essentielle, celle que le correcteur du Baccalauréat pose presque chaque année : \emph{si l'on mélange un acide d'un couple et la base d'un autre couple, que se passe-t-il ?} La réaction a-t-elle lieu ? Est-elle totale ou limitée ? Peut-on l'utiliser pour un dosage ? Toute cette section est consacrée à un outil remarquablement simple qui répond à ces questions : la constante de réaction $K$, calculée à partir des seuls p$K_a$ des couples en présence.

\courssub{Réaction entre deux couples acide/base : écriture générale}

Considérons deux couples acide/base :
\begin{itemize}
\item le couple 1 : \ce{A1H}/\ce{A1-}, de constante d'acidité $K_{a1}$ ;
\item le couple 2 : \ce{A2H}/\ce{A2-}, de constante d'acidité $K_{a2}$.
\end{itemize}

Mettons en présence l'\cle{acide du couple 1}, \ce{A1H}, et la \cle{base du couple 2}, notée \ce{B2} (c'est-à-dire \ce{A2-}). L'acide \ce{A1H} cède un proton \ce{H+} ; la base \ce{B2} le capte. La réaction qui se produit est :

\[ \ce{A1H + B2 <=> A1- + B2H+} \]

c'est-à-dire, avec les notations des couples :

\[ \ce{A1H + A2- <=> A1- + A2H} \]

\begin{attention}[Bien identifier les partenaires]
La réaction acido-basique met toujours en jeu \textbf{l'acide d'un couple} et \textbf{la base d'un autre couple}. On ne fait jamais réagir deux acides entre eux, ni deux bases entre elles : il faut un donneur de proton (\ce{A1H}) et un accepteur de proton (\ce{A2-}). Avant d'écrire quoi que ce soit, demande-toi : \og qui donne le proton ? qui le reçoit ? \fg
\end{attention}

\begin{exemplebox}[Le vinaigre et l'ammoniac]
Le vinaigre contient de l'acide éthanoïque \ce{CH3COOH} (couple \ce{CH3COOH}/\ce{CH3COO-}, p$K_{a1} = 4{,}8$). L'ammoniac \ce{NH3} est la base du couple \ce{NH4+}/\ce{NH3} (p$K_{a2} = 9{,}2$). Le mélange donne :

\[ \ce{CH3COOH + NH3 <=> CH3COO- + NH4+} \]

L'acide éthanoïque cède son proton à l'ammoniac. Reste à savoir : cette réaction va-t-elle jusqu'au bout ? C'est exactement l'objet de ce qui suit.
\end{exemplebox}

\courssub{Constante de la réaction : $K = 10^{pK_{a2} - pK_{a1}}$}

\courssub{Établissement de la formule}

Écrivons les deux équilibres d'acidité mis en jeu :

\[ \ce{A1H + H2O <=> A1- + H3O+} \qquad K_{a1} = \dfrac{[\ce{A1-}][\ce{H3O+}]}{[\ce{A1H}]} \]

\[ \ce{A2H + H2O <=> A2- + H3O+} \qquad K_{a2} = \dfrac{[\ce{A2-}][\ce{H3O+}]}{[\ce{A2H}]} \]

La réaction étudiée, \ce{A1H + A2- <=> A1- + A2H}, a pour constante d'équilibre :

\[ K = \dfrac{[\ce{A1-}][\ce{A2H}]}{[\ce{A1H}][\ce{A2-}]} \]

Multiplions le numérateur et le dénominateur par $[\ce{H3O+}]$ :

\[ K = \dfrac{[\ce{A1-}][\ce{H3O+}]}{[\ce{A1H}]} \times \dfrac{[\ce{A2H}]}{[\ce{A2-}][\ce{H3O+}]} = \dfrac{K_{a1}}{K_{a2}} \]

En passant aux p$K_a$ (rappel : $K_a = 10^{-pK_a}$) :

\[ K = \dfrac{10^{-pK_{a1}}}{10^{-pK_{a2}}} = 10^{pK_{a2} - pK_{a1}} \]

\begin{definition}[Constante d'une réaction acido-basique]
Pour la réaction entre l'acide \ce{A1H} du couple 1 (p$K_{a1}$) et la base \ce{A2-} du couple 2 (p$K_{a2}$) :
\[ \ce{A1H + A2- <=> A1- + A2H} \]
la constante d'équilibre vaut :
\[ \boxed{K = \dfrac{K_{a1}}{K_{a2}} = 10^{\,pK_{a2} - pK_{a1}}} \]
où \textbf{p$K_{a1}$ est le p$K_a$ du couple de l'acide qui réagit} (le donneur de proton) et \textbf{p$K_{a2}$ celui du couple de l'acide qui se forme} (le produit).
\end{definition}

\begin{attention}[Le piège classique : inverser les p$K_a$]
L'erreur la plus fréquente au Bac est d'écrire $K = 10^{pK_{a1} - pK_{a2}}$. Retiens ce moyen mnémotechnique : \og \textbf{formé moins réagi} \fg. Le p$K_a$ du couple de l'acide \textbf{formé} (produit) apparaît en premier dans l'exposant, moins le p$K_a$ du couple de l'acide qui \textbf{a réagi} (réactif). Une inversion change $10^{4{,}4}$ en $10^{-4{,}4}$ : la conclusion est alors exactement à l'inverse de la réalité !
\end{attention}

\courssub{Critère quantitatif : totale, limitée ou impossible ?}

La valeur de $K$ renseigne directement sur l'état final du système :

\begin{aretenir}[Critère de prévision]
Pour la réaction \ce{A1H + A2- <=> A1- + A2H} :
\begin{itemize}
\item si $K > 10^{4}$, c'est-à-dire $pK_{a2} - pK_{a1} > 4$ : la réaction est \textbf{quasi totale} (le réactif limitant est pratiquement entièrement consommé) ;
\item si $K < 10^{-4}$, c'est-à-dire $pK_{a2} - pK_{a1} < -4$ : la réaction \textbf{n'a pratiquement pas lieu} ;
\item si $10^{-4} < K < 10^{4}$ : on obtient un \textbf{équilibre} où les quatre espèces coexistent en proportions notables.
\end{itemize}
\end{aretenir}

L'idée intuitive est limpide : plus l'acide qui réagit est fort (petit p$K_{a1}$) et plus la base qui réagit est forte (grand p$K_{a2}$ de son couple), plus le transfert de proton est favorable, et plus $K$ est grande.

\begin{methode}[Prévoir une réaction acido-basique en quatre étapes]
\begin{enumerate}
\item \textbf{Identifier les deux couples} présents et repérer qui est l'acide (donneur) et qui est la base (accepteur) parmi les espèces mises en présence.
\item \textbf{Écrire l'équation bilan} \ce{A1H + A2- <=> A1- + A2H} en vérifiant l'équilibre en atomes et en charges.
\item \textbf{Relever les p$K_a$} dans le tableau de couples et calculer $K = 10^{pK_{a2} - pK_{a1}}$ (p$K_a$ de l'acide formé moins p$K_a$ de l'acide réactif).
\item \textbf{Conclure} en comparant $K$ à $10^{4}$ et $10^{-4}$ : quasi totale, équilibre, ou pas de réaction.
\end{enumerate}
\end{methode}

\begin{exoresolu}[Réaction de l'acide éthanoïque sur l'ammoniac]
On mélange une solution d'acide éthanoïque \ce{CH3COOH} (couple \ce{CH3COOH}/\ce{CH3COO-}, p$K_{a1} = 4{,}8$) avec une solution d'ammoniac \ce{NH3} (couple \ce{NH4+}/\ce{NH3}, p$K_{a2} = 9{,}2$). Écrire l'équation de la réaction, calculer sa constante $K$ et conclure.
\tcblower
\textbf{Étape 1 — Couples et partenaires.} \ce{CH3COOH} est l'acide du couple 1 ; \ce{NH3} est la base du couple 2. C'est donc \ce{CH3COOH} qui cède son proton.

\textbf{Étape 2 — Équation bilan.}
\[ \ce{CH3COOH + NH3 <=> CH3COO- + NH4+} \]
Vérification : même nombre d'atomes de chaque élément des deux côtés, charge totale nulle à gauche et $(-1) + (+1) = 0$ à droite. L'équation est correcte.

\textbf{Étape 3 — Calcul de $K$.} L'acide qui \emph{réagit} est \ce{CH3COOH} (p$K_{a1} = 4{,}8$) ; l'acide \emph{formé} est \ce{NH4+} (p$K_{a2} = 9{,}2$) :
\[ K = 10^{pK_{a2} - pK_{a1}} = 10^{9{,}2 - 4{,}8} = 10^{4{,}4} \approx 2{,}5 \times 10^{4} \]

\textbf{Étape 4 — Conclusion.} $K \approx 2{,}5 \times 10^{4} > 10^{4}$ : la réaction est \textbf{quasi totale}. Le mélange d'acide éthanoïque et d'ammoniac conduit pratiquement à une solution d'éthanoate d'ammonium \ce{CH3COO- + NH4+}. C'est d'ailleurs une réaction utilisée en laboratoire pour préparer ce sel.
\end{exoresolu}

\courssub{La règle du gamma : lecture graphique sur l'échelle des p$K_a$}

Plaçons les couples sur une \cle{échelle verticale de p$K_a$}, croissante de haut en bas. En haut se trouvent les acides les plus forts (petits p$K_a$) ; en bas, les bases les plus fortes (grands p$K_a$ de leurs couples). La règle du gamma ($\gamma$) s'énonce ainsi :

\begin{propriete}[Règle du gamma]
Sur une échelle verticale de p$K_a$ croissants vers le bas, la réaction naturelle (spontanée) se produit entre \textbf{l'acide le plus fort} (situé le plus haut) et \textbf{la base la plus forte} (située le plus bas), pour donner l'acide le plus faible et la base la plus faible. La flèche qui relie l'acide réactif à la base réactive dessine la lettre grecque $\gamma$.
\end{propriete}

\begin{popfigure}
\begin{center}
\begin{tikzpicture}[scale=1.0]
% Axe vertical pKa
\draw[->, thick, popdark] (0,-0.3) -- (0,10.6) node[above, font=\small\bfseries]{p$K_a$};
\node[rotate=90, font=\small, popmuted] at (-0.9,5.2) {p$K_a$ croissants vers le bas};
% Graduations
\foreach \y/\v in {0.5/0, 2.0/2, 3.5/4, 5.0/6, 6.5/8, 8.0/10, 9.5/12}
  \draw[popdark] (-0.12,\y) -- (0.12,\y) node[left=2pt, font=\small]{\v};
% Couple 1 : CH3COOH/CH3COO- pKa 4,8 -> y ~ 2.9
\draw[popblue, very thick] (0.6,2.9) -- (3.4,2.9);
\node[font=\small\bfseries, popblue, anchor=west] at (3.5,3.25) {\ce{CH3COOH}};
\node[font=\small\bfseries, popblue, anchor=west] at (3.5,2.55) {\ce{CH3COO-}};
\node[font=\footnotesize, popmuted, anchor=east] at (0.5,2.9) {4,8};
% Couple 2 : NH4+/NH3 pKa 9,2 -> y ~ 7.3
\draw[poporange, very thick] (0.6,7.3) -- (3.4,7.3);
\node[font=\small\bfseries, poporange, anchor=west] at (3.5,7.65) {\ce{NH4+}};
\node[font=\small\bfseries, poporange, anchor=west] at (3.5,6.95) {\ce{NH3}};
\node[font=\footnotesize, popmuted, anchor=east] at (0.5,7.3) {9,2};
% Ronds reactifs
\fill[popblue] (2.0,2.9) circle (3pt) node[above=4pt, font=\footnotesize\bfseries, popblue]{acide qui réagit};
\fill[poporange] (2.0,7.3) circle (3pt) node[below=4pt, font=\footnotesize\bfseries, poporange]{base qui réagit};
% Fleche gamma
\draw[->, very thick, popgreen, rounded corners=6pt] (2.0,2.9) -- (4.9,2.9) -- (4.9,7.3) -- (2.15,7.3);
\node[font=\small\bfseries, popgreen, anchor=west] at (5.05,5.1) {$\gamma$ : réaction};
\node[font=\small\bfseries, popgreen, anchor=west] at (5.05,4.6) {spontanée};
% Labels haut/bas
\node[font=\footnotesize\itshape, popmuted, anchor=west] at (6.6,0.5) {acides forts};
\node[font=\footnotesize\itshape, popmuted, anchor=west] at (6.6,9.5) {bases fortes};
\end{tikzpicture}
\end{center}
\legende{Règle du gamma : sur l'échelle verticale des p$K_a$, l'acide éthanoïque (couple de p$K_a = 4{,}8$) réagit spontanément avec l'ammoniac, base du couple de p$K_a = 9{,}2$. La flèche verte dessine un $\gamma$.}
\end{popfigure}

La règle du gamma permet de prévoir \emph{sans calcul} le sens d'une réaction dès que l'on dispose d'un tableau de couples classés par p$K_a$ croissant : il suffit de repérer l'acide placé le plus haut et la base placée le plus bas parmi les espèces mises en présence.

\begin{exemplebox}[Lecture d'un tableau de couples]
On donne les couples classés par p$K_a$ croissant : \ce{HCl}/\ce{Cl-} ($-\!7$) ; \ce{H3O+}/\ce{H2O} ($0$) ; \ce{CH3COOH}/\ce{CH3COO-} ($4{,}8$) ; \ce{NH4+}/\ce{NH3} ($9{,}2$) ; \ce{H2O}/\ce{HO-} ($14$). Si l'on mélange \ce{HCl} et \ce{NH3}, la règle du gamma désigne immédiatement \ce{HCl} (acide le plus haut) et \ce{NH3} (base la plus basse) : la réaction \ce{HCl + NH3 -> Cl- + NH4+} est spontanée, avec $K = 10^{9{,}2-(-7)} = 10^{16{,}2}$, immense. En revanche, mélanger \ce{CH3COO-} et \ce{NH4+} (les produits) ne donne quasiment rien : $K = 10^{4{,}8-9{,}2} = 10^{-4{,}4} < 10^{-4}$.
\end{exemplebox}

\courssub{Cas particulier fondamental : acide fort et base faible}

Un \cle{acide fort} est un acide dont le couple a un p$K_a$ négatif (p$K_a < 0$) : il est totalement ionisé dans l'eau et n'existe en solution que sous la forme \ce{H3O+}. Le couple réellement en jeu est donc \ce{H3O+}/\ce{H2O}, de p$K_a = 0$.

Pour la réaction d'un acide fort sur une base faible \ce{B} (couple \ce{BH+}/\ce{B} de p$K_a$ quelconque, positif) :

\[ \ce{H3O+ + B -> H2O + BH+} \qquad K = 10^{pK_a - 0} = 10^{pK_a} \]

Comme le p$K_a$ du couple de la base faible est toujours supérieur à 0 (et le plus souvent supérieur à 4), on a $K \gg 10^{4}$ :

\begin{aretenir}[Acide fort + base faible]
La réaction entre un acide fort (présent sous forme \ce{H3O+}) et une base faible est \textbf{toujours quasi totale}. De même, la réaction entre une base forte (\ce{HO-}) et un acide faible \ce{AH}, \ce{AH + HO- -> A- + H2O}, de constante $K = 10^{14 - pK_a}$, est toujours quasi totale. C'est le fondement de tous les dosages acido-basiques.
\end{aretenir}

\courssub{Lien avec les dosages acido-basiques}

\begin{propriete}[Conditions d'une réaction de dosage]
Une réaction utilisée pour un \cle{dosage} (titrage) doit être :
\begin{enumerate}
\item \textbf{totale} : $K > 10^{4}$, garanti par un écart de p$K_a$ supérieur à 4 ;
\item \textbf{rapide} : l'équilibre est atteint instantanément (c'est le cas des réactions acido-basiques, qui sont des transferts de protons) ;
\item \textbf{unique} : une seule réaction se produit entre le titrant et le titré.
\end{enumerate}
Le critère $K = 10^{pK_{a2} - pK_{a1}} > 10^{4}$ est donc l'outil théorique qui valide (ou invalide) le choix d'un titrant.
\end{propriete}

\begin{popfigure}
\begin{center}
\begin{tikzpicture}[
  node distance=0.55cm,
  box/.style={draw=popblue, thick, fill=popblueL, rounded corners=4pt, text width=4.6cm, align=center, font=\small, inner sep=5pt},
  dec/.style={draw=poporange, thick, fill=poporangeL, diamond, aspect=2.2, align=center, font=\small, inner sep=2pt},
  end/.style={draw=popgreen, thick, fill=popgreenL, rounded corners=4pt, text width=4.2cm, align=center, font=\small\bfseries, inner sep=5pt},
  arr/.style={->, thick, popdark}
]
\node[box] (cpl) {Identifier les deux couples et écrire \ce{A1H + A2- <=> A1- + A2H}};
\node[box, below=of cpl] (calc) {Calculer $K = 10^{pK_{a2} - pK_{a1}}$};
\node[dec, below=of calc] (test) {$K > 10^{4}$ ?};
\node[end, below left=0.7cm and -1.2cm of test] (oui) {Réaction quasi totale : dosage possible};
\node[end, below right=0.7cm and -1.2cm of test] (non) {Équilibre limité ou réaction nulle : dosage impossible};
\draw[arr] (cpl) -- (calc);
\draw[arr] (calc) -- (test);
\draw[arr] (test) -| node[above left, font=\small\bfseries, popgreen]{oui} (oui);
\draw[arr] (test) -| node[above right, font=\small\bfseries, poppink]{non} (non);
\end{tikzpicture}
\end{center}
\legende{Organigramme de décision : écrire les couples, calculer $K$, comparer à $10^{4}$, conclure.}
\end{popfigure}

\begin{exemplebox}[Validation du titrage du vinaigre par la soude]
Le vinaigre est une solution d'acide éthanoïque (\ce{CH3COOH}/\ce{CH3COO-}, p$K_a = 4{,}8$). On le dose par la soude, solution d'hydroxyde de sodium contenant les ions \ce{HO-} (couple \ce{H2O}/\ce{HO-}, p$K_a = 14$). La réaction de dosage est :
\[ \ce{CH3COOH + HO- -> CH3COO- + H2O} \]
Sa constante vaut :
\[ K = 10^{14 - 4{,}8} = 10^{9{,}2} \approx 1{,}6 \times 10^{9} \gg 10^{4} \]
La réaction est quasi totale, instantanée et unique : c'est une excellente réaction de dosage, et le saut de pH à l'équivalence est net.
\end{exemplebox}

\begin{savaistu}[Pourquoi le saut de pH du titrage du vinaigre est-il si net ?]
C'est précisément le critère $K = 10^{pK_{a2} - pK_{a1}}$ qui garantit le succès du titrage du vinaigre par la soude, grand classique du Baccalauréat C, D et E au Cameroun. Parce que $K \approx 10^{9}$, chaque goutte de soude versée avant l'équivalence est intégralement consommée : le pH évolue de façon contrôlée, puis \emph{bascule} brutalement à l'équivalence, lorsque le moindre excès de \ce{HO-} n'a plus d'acide à neutraliser. Si la réaction n'était pas totale, la courbe pH $= f(V)$ s'arrondirait autour de l'équivalence et le saut deviendrait indétectable : impossible alors de repérer le point équivalent, ni avec un indicateur coloré, ni avec un pH-mètre. Les savonniers et vinaigriers artisans appliquent le même principe sans le savoir : pour vérifier l'acidité d'une production, il faut un réactif qui réagit \emph{complètement} avec l'acide.
\end{savaistu}

\begin{exercice}[À toi de jouer]
On donne les couples : \ce{HCOOH}/\ce{HCOO-} (p$K_a = 3{,}8$) et \ce{C6H5COOH}/\ce{C6H5COO-} (p$K_a = 4{,}2$).
\begin{enumerate}
\item On mélange une solution d'acide méthanoïque \ce{HCOOH} et une solution de benzoate de sodium \ce{C6H5COO-}. Écrire l'équation de la réaction et calculer sa constante $K$. Conclure.
\item On mélange maintenant de l'acide benzoïque \ce{C6H5COOH} avec des ions méthanoate \ce{HCOO-}. Que se passe-t-il ?
\item Placer ces deux couples sur une échelle verticale de p$K_a$ et vérifier les conclusions par la règle du gamma.
\end{enumerate}
\end{exercice}

\begin{corrige}[Exercice 1]
\textbf{1.} L'acide qui réagit est \ce{HCOOH} (p$K_{a1} = 3{,}8$) ; l'acide formé est \ce{C6H5COOH} (p$K_{a2} = 4{,}2$) :
\[ \ce{HCOOH + C6H5COO- <=> HCOO- + C6H5COOH} \qquad K = 10^{4{,}2 - 3{,}8} = 10^{0{,}4} \approx 2{,}5 \]
$10^{-4} < K < 10^{4}$ : la réaction conduit à un \textbf{équilibre} où les quatre espèces coexistent ; elle n'est ni totale, ni nulle.

\textbf{2.} Pour la réaction inverse, $K' = 10^{3{,}8 - 4{,}2} = 10^{-0{,}4} \approx 0{,}4$ : c'est la même situation d'équilibre vue de l'autre côté ($K' = 1/K$). Aucune des deux réactions n'est totale.

\textbf{3.} Sur l'échelle verticale, \ce{HCOOH} (p$K_a = 3{,}8$) est légèrement au-dessus de \ce{C6H5COOH} (p$K_a = 4{,}2$). La règle du gamma confirme que \ce{HCOOH} réagit avec \ce{C6H5COO-}, mais l'écart de p$K_a$ (seulement $0{,}4$) est trop faible pour que la réaction soit totale.
\end{corrige}

\begin{aretenir}[L'essentiel de la section]
\begin{itemize}
\item La réaction entre l'acide d'un couple 1 et la base d'un couple 2 s'écrit \ce{A1H + A2- <=> A1- + A2H}.
\item Sa constante vaut $K = \dfrac{K_{a1}}{K_{a2}} = 10^{pK_{a2} - pK_{a1}}$ : p$K_a$ de l'acide \textbf{formé} moins p$K_a$ de l'acide \textbf{réactif}.
\item $K > 10^{4}$ : réaction quasi totale ; $K < 10^{-4}$ : pas de réaction ; entre les deux : équilibre.
\item La règle du gamma, sur une échelle verticale de p$K_a$, désigne l'acide le plus fort et la base la plus forte comme partenaires de la réaction spontanée.
\item Acide fort + base faible (et base forte + acide faible) : réaction toujours quasi totale.
\item Une réaction de dosage doit être totale ($K > 10^{4}$), rapide et unique : c'est ce critère qui valide tout titrage acido-basique.
\end{itemize}
\end{aretenir}

\newpage

\coursec{Travaux pratiques}

\courssub{Objectifs du TP}

À l'issue de cette séance de travaux pratiques, tu dois être capable de :

\begin{itemize}
\item préparer par pesée et par dilution des solutions de concentration connue ;
\item utiliser correctement un pH-mètre (étalonnage à deux points, rinçage, lecture) ;
\item montrer, à partir d'une mesure de pH et de la concentration apportée, que l'ion éthanoate est une \cle{base faible} (réaction non totale avec l'eau) ;
\item calculer le \cle{coefficient d'ionisation} $\alpha$ (ou degré d'hydrolyse) d'une base faible ;
\item déterminer expérimentalement la \cle{constante d'acidité} $K_a$ et le $pK_a$ du couple \ce{CH3COOH}/\ce{CH3COO-} par deux méthodes indépendantes (mesure du pH de l'acide seul, puis pH d'un mélange équimolaire acide/sel) ;
\item comparer les valeurs expérimentales à la valeur tabulée $pK_a = 4,8$ (à \SI{25}{\celsius}) et discuter les sources d'erreur.
\end{itemize}

\begin{aretenir}[Idée directrice du TP]
Pour une base forte totalement dissociée de concentration $C$, on aurait $pH = 14 + \log C$ (à \SI{25}{\celsius}). Si le pH mesuré est \emph{inférieur} à cette valeur, la base n'est que partiellement ionisée : c'est une \cle{base faible}. Ensuite, la mesure du pH d'une solution d'acide éthanoïque permet de remonter à $K_a$, car à l'équilibre $[\ce{H3O+}] = [\ce{CH3COO-}]$ et $[\ce{CH3COOH}] = C - [\ce{H3O+}]$. Le mélange équimolaire acide/base conjuguée donne directement $pH = pK_a$.
\end{aretenir}

\courssub{Matériel et produits}

\begin{center}
\begin{tabularx}{\linewidth}{|l|X|}
\hline
\rowcolor{popblue!40} \textbf{Matériel} & \textbf{Produits} \\
\hline
\rowcolor{popblueL} pH-mètre étalonné (ou papier pH de précision) & Acide éthanoïque, solution à \SI{0,10}{mol.L^{-1}} \\
\hline
Balance de précision (au centigramme \SI{0,01}{g}) & Éthanoate de sodium solide \ce{CH3COONa} ($M = \SI{82,0}{g.mol^{-1}}$) \\
\hline
\rowcolor{popblueL} Fioles jaugées de \SI{100}{mL} (2) et \SI{50}{mL} (1) & Eau distillée (ou eau bouillie puis refroidie, sans \ce{CO2}) \\
\hline
Pipettes jaugées de \SI{10}{mL} et \SI{20}{mL} (avec poire aspirante) & Solutions tampons pH = 4,00 ; pH = 7,00 ; pH = 9,00 (étalonnage) \\
\hline
\rowcolor{popblueL} Béchers de \SI{100}{mL} (4), agitateur magnétique + barreau & \\
\hline
Spatule, entonnoir, pissette, lunettes, blouse, gants & \\
\hline
\end{tabularx}
\end{center}

\begin{methode}[Préparation rigoureuse de la solution d'éthanoate de sodium \SI{1,0e-2}{mol.L^{-1}}]
Peser directement \SI{0,082}{g} sur une balance au centigramme est imprécis (erreur relative \textgreater 10\%). On procède en deux temps :
\begin{enumerate}
\item \textbf{Solution mère \SI{0,10}{mol.L^{-1}} :} Peser \textbf{\SI{0,820}{g}} d'éthanoate de sodium (erreur \textless 1\%), dissoudre dans une fiole jaugée de \SI{100}{mL}, compléter au trait de jauge.
\item \textbf{Dilution au dixième :} Prélever \SI{10,0}{mL} de cette solution mère à l'aide d'une pipette jaugée, verser dans une fiole jaugée de \SI{100}{mL} (propre, rincée à l'eau distillée), compléter au trait de jauge. On obtient \SI{100}{mL} de solution \SI{1,0e-2}{mol.L^{-1}}.
\end{enumerate}
\end{methode}

\begin{savaistu}[Et si un produit manque ?]
Dans beaucoup de lycées camerounais, le pH-mètre est rare ou en panne. Deux alternatives honnêtes :
\begin{itemize}
\item le \cle{papier pH de précision} (bandes 3,8--5,5 et 8--10) permet des mesures à $\pm 0,3$ unité : suffisant pour montrer que la base est faible, moins précis pour $K_a$ ;
\item l'éthanoate de sodium peut être \emph{préparé sur place} en neutralisant exactement \SI{10}{mL} d'acide éthanoïque \SI{0,10}{mol.L^{-1}} par la soude \SI{0,10}{mol.L^{-1}} en présence de phénolphtaléine : le point de virage correspond à une solution d'éthanoate de sodium. Le vinaigre de table (acide éthanoïque dilué à environ \SI{5}{\percent}) peut aussi remplacer l'acide pur après dosage préalable.
\end{itemize}
\end{savaistu}

\courssub{Sécurité}

\begin{attention}[Gestes qui protègent]
\begin{itemize}
\item L'\textbf{acide éthanoïque concentré} est \emph{corrosif} et \emph{inflammable} (pictogrammes : corrosion, flamme). Les solutions diluées utilisées ici sont moins dangereuses mais provoquent des irritations : port obligatoire des \textbf{lunettes} et de la \textbf{blouse}, gants pour la manipulation de la solution mère.
\item Ne jamais pipeter à la bouche : utiliser une \textbf{poire aspirante} ou un propipette.
\item En cas de contact avec la peau ou les yeux : rincer abondamment à l'eau claire pendant plusieurs minutes et prévenir l'enseignant.
\item L'éthanoate de sodium solide est peu dangereux mais poussiéreux : éviter d'inhaler, se laver les mains en fin de séance.
\item La sonde du pH-mètre est fragile et coûteuse : ne jamais la laisser sécher à l'air, la conserver dans sa solution de conservation (\ce{KCl} saturé), ne jamais frotter le bulbe en verre.
\end{itemize}
\end{attention}

\courssub{Schéma du dispositif}

\begin{popfigure}
\begin{center}
\begin{tikzpicture}[scale=1.0, >=Stealth]
% Bécher
\draw[thick, popdark] (0,0) -- (0,3.2) (2.6,0) -- (2.6,3.2);
\draw[thick, popdark] (0,0) .. controls (0.2,-0.25) and (2.4,-0.25) .. (2.6,0);
% Liquide
\fill[popblue!25] (0.06,0.05) rectangle (2.54,2.0);
\draw[popblue, thick] (0.06,2.0) -- (2.54,2.0);
% Sonde pH-mètre
\draw[thick, popdark, fill=popcream] (1.05,4.6) rectangle (1.55,1.2);
\draw[thick, popdark, fill=popblue!50] (1.12,1.9) rectangle (1.48,1.2);
\draw[popmuted] (1.3,4.6) -- (1.3,5.3);
\node[popmuted, font=\small] at (1.3,5.6) {câble};
% Boîtier pH-mètre
\draw[thick, popdark, fill=popgold!30, rounded corners] (4.0,4.2) rectangle (6.6,5.6);
\node[font=\small\bfseries, popdark] at (5.3,5.25) {pH-mètre};
\draw[fill=white, draw=popdark] (4.3,4.4) rectangle (6.3,5.0);
\node[font=\small\ttfamily, popink] at (5.3,4.7) {pH = 8,87};
\draw[popmuted] (1.3,5.3) -- (1.3,5.9) -- (4.0,5.9) -- (4.0,5.3);
% Agitateur
\draw[thick, poppurple] (2.1,0.35) ellipse (0.35 and 0.12);
\node[font=\scriptsize, poppurple] at (3.4,0.35) {barreau magnétique};
% Légendes
\node[font=\small, popdark, align=center] at (1.3,-0.8) {solution à étudier\\ (bécher de \SI{100}{mL})};
\node[font=\small, popdark, align=left] at (5.3,3.0) {sonde de verre\\ plongeant dans\\ la solution};
\draw[->, popmuted] (4.6,3.2) -- (1.7,2.3);
\end{tikzpicture}
\end{center}
\legende{Dispositif de mesure du pH : le bulbe de la sonde doit être entièrement immergé, sans toucher le barreau magnétique ; agitation modérée pendant la mesure, puis lecture après stabilisation (dérive \textless 0,01 pH/min).}
\end{popfigure}

\courssub{Protocole}

\textbf{Étalonnage préalable du pH-mètre (deux points)}\par
Rincer la sonde à l'eau distillée, l'essuyer délicatement par tamponnement avec du papier absorbant (ne pas frotter). Étalonner l'appareil avec les solutions tampons \textbf{pH = 7,00} puis \textbf{pH = 9,00} (ou 10,00) car les mesures attendues sont basiques (partie 1) et acides (partie 2). Rincer la sonde à l'eau distillée entre chaque tampon et avant la première mesure. Noter la température de la salle $T$ (\SI{}{\celsius}).

\textbf{Partie 1 : Mise en évidence d'une base faible (ion éthanoate)}\par
\begin{enumerate}
\item Préparer \SI{100}{mL} de solution d'éthanoate de sodium de concentration $C_b = \SI{1,0e-2}{mol.L^{-1}}$ selon la \textbf{Méthode} ci-dessus (solution mère + dilution).
\item Verser environ \SI{40}{mL} dans un bécher propre, immerger la sonde, agiter modérément, attendre la stabilisation. Noter $pH_1$ et la température.
\item Calculer le pH théorique d'une base forte de même concentration : $pH_{\text{th}} = 14 + \log C_b = 12,00$ (à \SI{25}{\celsius}). Comparer.
\end{enumerate}

\textbf{Partie 2 : Détermination du $pK_a$ à partir de l'acide seul}\par
\begin{enumerate}
\item Prélever \SI{10,0}{mL} de la solution d'acide éthanoïque à \SI{0,10}{mol.L^{-1}} (pipette jaugée) et diluer au dixième dans une fiole jaugée de \SI{100}{mL} : on obtient une solution de concentration $C_a = \SI{1,0e-2}{mol.L^{-1}}$. Homogénéiser.
\item Rincer la sonde, mesurer le pH de cette solution. Noter $pH_2$.
\end{enumerate}

\textbf{Partie 3 (Vérification à la demi-équivalence) :}\par
\begin{enumerate}
\item Mélanger \SI{20,0}{mL} de la solution d'acide éthanoïque $C_a = \SI{1,0e-2}{mol.L^{-1}}$ et \SI{20,0}{mL} de la solution d'éthanoate de sodium $C_b = \SI{1,0e-2}{mol.L^{-1}}$ dans un bécher.
\item Rincer la sonde, mesurer le pH du mélange. Noter $pH_3$.
\end{enumerate}

\courssub{Observations et résultats attendus}

\begin{center}
\begin{tabularx}{\linewidth}{|l|X|X|}
\hline
\rowcolor{popgreen!40} \textbf{Mesure} & \textbf{Valeur typique attendue (à \SI{25}{\celsius})} & \textbf{Interprétation immédiate} \\
\hline
\rowcolor{popgreenL} $pH_1$ (éthanoate de sodium, \SI{1,0e-2}{mol.L^{-1}}) & $8,9 \pm 0,1$ & $pH_1 < 12$ : base faible (hydrolyse partielle) \\
\hline
$pH_2$ (acide éthanoïque, \SI{1,0e-2}{mol.L^{-1}}) & $3,4 \pm 0,1$ & $pH_2 > 2$ : acide faible (ionisation partielle) \\
\hline
\rowcolor{popgreenL} $pH_3$ (mélange équimolaire $V_a = V_b$) & $4,8 \pm 0,2$ & $pH_3 \approx pK_a$ (relation de Henderson-Hasselbalch) \\
\hline
Température de la salle $T$ & \SI{25}{\celsius} (noter la valeur réelle) & $K_e$, $K_a$, $pK_a$ dépendent de $T$ \\
\hline
\end{tabularx}
\end{center}

\courssub{Exploitation}

\begin{exoresolu}[Exploitation guidée des mesures]
\textbf{Questions.}
\begin{enumerate}
\item Pour la solution d'éthanoate de sodium, calculer la concentration en ions hydroxide réellement présents à partir de $pH_1 = 8,9$. Conclure sur la force de la base.
\item Calculer le coefficient d'ionisation $\alpha$ (degré d'hydrolyse) de l'ion éthanoate.
\item Écrire l'équation de l'équilibre d'ionisation de l'acide éthanoïque dans l'eau et dresser le tableau d'avancement correspondant.
\item À partir de $pH_2 = 3,4$, calculer $[\ce{H3O+}]$, puis $K_a$ et $pK_a$ du couple \ce{CH3COOH}/\ce{CH3COO-}.
\item Justifier que dans le mélange de la partie 3, on a $[\ce{CH3COOH}] = [\ce{CH3COO-}]$, et en déduire une seconde valeur du $pK_a$.
\item Comparer les deux valeurs expérimentales à la valeur tabulée $pK_a = 4,8$ (à \SI{25}{\celsius}) et citer trois sources d'erreur principales.
\end{enumerate}
\tcblower
\textbf{Solution détaillée.}
\begin{enumerate}
\item \textbf{Concentration en ions hydroxyde :}
\[ [\ce{H3O+}] = 10^{-pH_1} = 10^{-8,9} = \SI{1,26e-9}{mol.L^{-1}} \]
À \SI{25}{\celsius}, $K_e = [\ce{H3O+}][\ce{HO-}] = 10^{-14}$.
\[ [\ce{HO-}] = \frac{K_e}{[\ce{H3O+}]} = \frac{10^{-14}}{1,26 \times 10^{-9}} = \SI{7,94e-6}{mol.L^{-1}} \]
\textbf{Conclusion :} $[\ce{HO-}] = \SI{7,9e-6}{mol.L^{-1}} \ll C_b = \SI{1,0e-2}{mol.L^{-1}}$. La réaction de l'ion éthanoate avec l'eau ($\ce{CH3COO- + H2O <=> CH3COOH + HO-}$) est \emph{très limitée} : l'ion éthanoate est une \cle{base faible}.
\item \textbf{Coefficient d'ionisation $\alpha$ :}
\[ \alpha = \frac{[\ce{HO-}]_{\text{eq}}}{C_b} = \frac{7,94 \times 10^{-6}}{1,0 \times 10^{-2}} = \SI{7,9e-4}{} \]
Soit $\alpha \approx \SI{0,079}{\percent}$. Moins de \SI{0,1}{\percent} des ions éthanoate ont capté un proton.
\item \textbf{Équilibre et tableau d'avancement (acide éthanoïque) :}
\[ \ce{CH3COOH_{(aq)} + H2O_{(l)} <=> CH3COO-_{(aq)} + H3O+_{(aq)}} \]
\begin{center}
\begin{tabular}{lcccc}
\hline
 & \ce{CH3COOH} & \ce{CH3COO-} & \ce{H3O+} \\
\hline
État initial (mol.L$^{-1}$) & $C_a$ & $0$ & $\approx 0$ (eau pure) \\
État d'équilibre (mol.L$^{-1}$) & $C_a - x_f$ & $x_f$ & $x_f$ \\
\hline
\end{tabular}
\end{center}
\item \textbf{Calcul de $K_a$ et $pK_a$ (Méthode 1) :}
$pH_2 = 3,4 \Rightarrow [\ce{H3O+}] = 10^{-3,4} = \SI{3,98e-4}{mol.L^{-1}} = x_f$.
D'après le tableau : $[\ce{CH3COO-}] = x_f = \SI{3,98e-4}{mol.L^{-1}}$.
$[\ce{CH3COOH}] = C_a - x_f = \SI{1,0e-2}{} - \SI{3,98e-4}{} = \SI{9,60e-3}{mol.L^{-1}}$.
\[ K_a = \frac{[\ce{CH3COO-}][\ce{H3O+}]}{[\ce{CH3COOH}]} = \frac{(\SI{3,98e-4}{})^2}{\SI{9,60e-3}{}} = \SI{1,65e-5}{} \]
\[ pK_a = -\log(K_a) = -\log(\SI{1,65e-5}{}) = 4,78 \approx 4,8 \]
\item \textbf{Méthode 2 (Mélange équimolaire / Demi-équivalence) :}
Les deux solutions mères ont même concentration $C = \SI{1,0e-2}{mol.L^{-1}}$ et sont mélangées à volumes égaux ($V_a = V_b = \SI{20,0}{mL}$).
Quantités de matière : $n(\ce{CH3COOH}) = n(\ce{CH3COO-}) = C \times V$.
Dans le mélange (volume total $V_{\text{tot}} = \SI{40,0}{mL}$) :
\[ [\ce{CH3COOH}] = \frac{n(\ce{CH3COOH})}{V_{\text{tot}}} = \frac{C \times V}{2V} = \frac{C}{2} \]
\[ [\ce{CH3COO-}] = \frac{n(\ce{CH3COO-})}{V_{\text{tot}}} = \frac{C}{2} \]
Donc $[\ce{CH3COOH}] = [\ce{CH3COO-}]$.
La relation de Henderson-Hasselbalch : $pH = pK_a + \log\dfrac{[\ce{CH3COO-}]}{[\ce{CH3COOH}]}$ donne :
\[ pH_3 = pK_a + \log(1) = pK_a \]
Mesure $pH_3 = 4,8 \Rightarrow pK_a \approx 4,8 \quad (K_a \approx \SI{1,6e-5}{})$.
\item \textbf{Comparaison et sources d'erreur :}
Les deux méthodes donnent $pK_a \approx 4,8$, en excellent accord avec la valeur tabulée ($pK_a = 4,76$ à \SI{25}{\celsius}, souvent arrondi à 4,8 au Bac).
\textbf{Sources d'erreur principales :}
\begin{enumerate}
\item \textbf{Température :} $K_e$ et $K_a$ varient avec $T$. Si $T \neq \SI{25}{\celsius}$, la valeur théorique de référence change.
\item \textbf{Étalonnage pH-mètre :} Mauvais choix des tampons (ex: pH 4 et 7 pour mesurer pH 9), sonde sale, dérive de l'appareil.
\item \textbf{\ce{CO2} dissous :} L'eau distillée non bouillie contient du \ce{CO2} qui forme de l'acide carbonique, acidifiant les solutions basiques (surévalue $[\ce{H3O+}]$, sous-évalue $[\ce{HO-}]$ et $pH_1$).
\item \textbf{Précision des volumes/masses :} Erreur sur la pesée (surtout si masse unique \SI{0,082}{g}), erreur sur les volumes pipetés/jaugés.
\end{enumerate}
\end{enumerate}
\end{exoresolu}

\begin{attention}[Pièges fréquents dans les calculs]
\begin{itemize}
\item Ne pas écrire $[\ce{CH3COOH}] = C_a$ à l'équilibre : il faut retrancher $[\ce{H3O+}]$, même si la correction est petite ici ($x_f/C_a \approx 4\%$).
\item Attention aux puissances de dix : $pH = 3,4$ donne $[\ce{H3O+}] = 10^{-3,4}$ et non $3,4 \times 10^{-1}$.
\item Le coefficient $\alpha$ est sans unité ; le donner aussi en pourcentage pour l'interpréter (\SI{0,08}{\percent}).
\item Ne pas confondre $K_a$ (constante d'acidité de l'acide) et $K_b$ (constante de basicité de la base conjuguée) : $K_a \times K_b = K_e$.
\item Pour le mélange équimolaire, l'égalité des concentrations \emph{dans le mélange final} vient de l'égalité des quantités de matière initiales et du même volume total, pas de l'égalité des concentrations mères seules.
\end{itemize}
\end{attention}

\begin{aretenir}[Liens avec le cours : Couples de l'eau et Ampholytes]
\begin{itemize}
\item L'eau réalise deux couples acide/base : \ce{H3O+}/\ce{H2O} (acide fort / base très faible) et \ce{H2O}/\ce{HO-} (acide très faible / base forte).
\item L'ion éthanoate \ce{CH3COO-} est la base conjuguée de l'acide éthanoïque \ce{CH3COOH}. Le couple est \ce{CH3COOH}/\ce{CH3COO-}.
\item L'eau est un \cle{ampholyte} : elle peut agir comme acide (donne \ce{H+} à \ce{CH3COO-}) ou comme base (reçoit \ce{H+} de \ce{CH3COOH}).
\item \textbf{Loi d'Ostwald (notion) :} Si l'on dilue la solution d'éthanoate de sodium, $\alpha$ augmente (le degré d'hydrolyse croît quand la concentration diminue), mais $[\ce{HO-}] = \alpha C$ diminue.
\end{itemize}
\end{aretenir}

\courssub{Conclusion}

Ce TP illustre la démarche complète d'identification d'un acide ou d'une base faible : la comparaison du pH mesuré au pH qu'aurait une espèce forte de même concentration montre que l'ion éthanoate ne réagit que très partiellement avec l'eau ($\alpha \approx \SI{0,08}{\percent}$), ce qui prouve l'existence d'un \cle{équilibre chimique}. L'exploitation quantitative des mesures conduit, par deux voies indépendantes, à $pK_a \approx 4,8$ pour le couple \ce{CH3COOH}/\ce{CH3COO-}, en excellent accord avec la valeur tabulée. La méthode du mélange équimolaire ($pH = pK_a$ à la demi-équivalence) sera réinvestie dans l'étude des dosages acido-basiques et des solutions tampons : c'est le principe même du vinaigre, de nombreux médicaments et des tampons biologiques qui stabilisent le pH du sang (système carbonates/bicarbonates).

\newpage

\coursec{Méthodes et exercices résolus}

\coursec{Leçon 08 : Acides faibles, bases faibles et couples acide/base}

\noindent Dans cette leçon, nous allons comprendre pourquoi certains acides (comme celui du vinaigre ou du vin de palme) et certaines bases (comme l'ammoniac ou les carbonates de la cendre) ne se dissocient pas totalement dans l'eau. Nous apprendrons à quantifier cette \cle{ionisation partielle}, à identifier les \cle{couples acide/base} mis en jeu, à utiliser leurs constantes ($K_a$, $K_b$, $\mathrm{p}K_a$) pour classer leurs forces, à prédire le pH d'une solution et à prévoir le sens d'une réaction acido-basique. Ces outils sont indispensables pour comprendre le dosage du vinaigre, la fabrication artisanale du savon, l'action tampon du sang ou encore le choix d'un indicateur coloré.

\begin{savaistu}[Le vinaigre et le bili-bili : acides faibles du quotidien]
Au Cameroun, le vinaigre de palme (issu de la fermentation acétique du vin de palme) et le bili-bili (bière de mil ou de maïs) doivent leur acidité à la présence d'acide éthanoïque et d'acide lactique, tous deux \textbf{acides faibles}. Contrairement à l'acide chlorhydrique (acide fort) utilisé dans l'industrie (SONARA, CIMENCAM), ces acides naturels ne libèrent qu'une faible fraction de leurs protons dans l'eau. C'est cette ionisation partielle qui donne au vinaigre son pH typique (autour de 2,5--3,0) et au bili-bili sa saveur acidulée sans être corrosif. Comprendre l'équilibre de ces acides faibles permet de contrôler la qualité de ces produits traditionnels.
\end{savaistu}

\courssub{Acide faible : l'acide éthanoïque}

\begin{definition}[Acide faible]
Un \cle{acide faible} est un acide dont la réaction avec l'eau (ionisation) est \textbf{limitée} : elle ne va pas jusqu'à son terme et conduit à un \cle{équilibre chimique}. L'acide n'est que \textbf{partiellement ionisé} ; il coexiste dans la solution avec sa base conjuguée et les ions oxonium.
\end{definition}

\begin{propriete}[Équation d'ionisation de l'acide éthanoïque]
L'acide éthanoïque (acide acétique), $\ce{CH3COOH}$, réagit avec l'eau selon l'équilibre :
\[ \ce{CH3COOH_{(aq)} + H2O_{(l)} <=> CH3COO-_{(aq)} + H3O+_{(aq)}} \]
La flèche d'équilibre \ce{<=>} traduit le caractère réversible et incomplet de la transformation.
\end{propriete}

\begin{methode}[Le réflexe « comparaison pH / $-\log C$ » pour un acide]
Pour savoir si un acide est fort ou faible à partir d'une mesure de pH :
\begin{enumerate}
\item Calculer la concentration molaire $C$ de la solution acide (attention aux dilutions).
\item Calculer la valeur théorique $-\log C$ (pH attendu si l'acide était fort, c'est-à-dire totalement ionisé : $[\ce{H3O+}] = C$).
\item Comparer au pH \emph{mesuré} expérimentalement :
\begin{itemize}
\item si $\mathrm{pH} = -\log C$ (à l'incertitude près) $\Rightarrow$ acide \cle{fort} (ionisation totale) ;
\item si $\mathrm{pH} > -\log C$ $\Rightarrow$ acide \cle{faible} (ionisation partielle, $[\ce{H3O+}] < C$).
\end{itemize}
\item Conclure : « Le pH mesuré est supérieur à $-\log C$, la concentration en ions $\ce{H3O+}$ est inférieure à $C$, donc l'ionisation est partielle : c'est un acide faible. »
\end{enumerate}
\textbf{Rappel :} à \SI{25}{\celsius}, $\mathrm{pH} = -\log[\ce{H3O+}]$ avec $[\ce{H3O+}]$ en \SI{}{mol.L^{-1}}.
\end{methode}

\begin{exoresolu}[Le vinaigre du marché est-il un acide fort ?]
Un vinaigre commercial (solution aqueuse d'acide éthanoïque $\ce{CH3COOH}$) est dilué 10 fois. La solution diluée $S$ a une concentration $C = \SI{8,0e-2}{mol.L^{-1}}$. Le pH mesuré à \SI{25}{\celsius} vaut $2,9$.
\begin{enumerate}
\item Montrer que l'acide éthanoïque est un acide faible.
\item Calculer la concentration des ions $\ce{H3O+}$ et en déduire celle des ions $\ce{CH3COO-}$.
\end{enumerate}
\tcblower
\textbf{1. Nature de l'acide.}

Si l'acide éthanoïque était fort, l'ionisation serait totale :
\[ \ce{CH3COOH + H2O -> CH3COO- + H3O+} \quad\Rightarrow\quad [\ce{H3O+}] = C \]
Le pH théorique serait :
\[ \mathrm{pH}_{\text{th}} = -\log C = -\log(8,0\times 10^{-2}) = 1,10. \]
Or le pH mesuré est $2,90$. Comme $2,90 > 1,10$, la solution contient \textbf{moins} d'ions $\ce{H3O+}$ que prévu pour une ionisation totale.
\par L'acide éthanoïque est donc un \cle{acide faible} ; sa dissolution conduit à l'équilibre :
\[ \ce{CH3COOH + H2O <=> CH3COO- + H3O+}. \]

\textbf{2. Concentrations ioniques.}

D'après la définition du pH :
\[ [\ce{H3O+}] = 10^{-\mathrm{pH}} = 10^{-2,90} = \SI{1,26e-3}{mol.L^{-1}} \approx \SI{1,3e-3}{mol.L^{-1}}. \]
D'après la stœchiométrie de l'équilibre (1 mol $\ce{CH3COOH}$ ionisée donne 1 mol $\ce{H3O+}$ et 1 mol $\ce{CH3COO-}$), et en négligeant la contribution de l'autoprotolyse de l'eau ($[\ce{H3O+}]_{\text{eau}} \approx \SI{1e-7}{mol.L^{-1}} \ll \SI{1,3e-3}{mol.L^{-1}}$) :
\[ [\ce{CH3COO-}] \approx [\ce{H3O+}] = \SI{1,3e-3}{mol.L^{-1}}. \]

\textbf{Vérification :} $[\ce{H3O+}] = \SI{1,3e-3}{mol.L^{-1}} \ll C = \SI{8,0e-2}{mol.L^{-1}}$. Seule une faible fraction ($\approx 1,6\,\%$) de l'acide est ionisée.

\textbf{Conclusion :} le pH du vinaigre dilué ($2,9$) est bien supérieur à $-\log C$ ($1,1$) : l'acide éthanoïque est un acide faible, partiellement ionisé dans l'eau.
\end{exoresolu}

\courssub{Base faible : éthanoate de sodium et ammoniac}

\begin{definition}[Base faible]
Une \cle{base faible} est une base dont la réaction avec l'eau (prise de proton sur $\ce{H2O}$) est \textbf{limitée} et conduit à un équilibre chimique. Elle n'est que \textbf{partiellement protonée} ; elle coexiste avec son acide conjugué et les ions hydroxyde.
\end{definition}

\begin{propriete}[Équations d'ionisation de bases faibles courantes]
\begin{itemize}
\item \textbf{Ammoniac} ($\ce{NH3}$, gaz dissous) : $\ce{NH3_{(aq)} + H2O_{(l)} <=> NH4+_{(aq)} + HO-_{(aq)}}$
\item \textbf{Ion éthanoate} (issu de l'éthanoate de sodium $\ce{CH3COONa}$, sel fort totalement dissocié) :
$\ce{CH3COO-_{(aq)} + H2O_{(l)} <=> CH3COOH_{(aq)} + HO-_{(aq)}}$
\end{itemize}
Dans les deux cas, la production d'ions $\ce{HO-}$ est partielle : $[\ce{HO-}] < C$ (concentration initiale de la base).
\end{propriete}

\begin{methode}[Le réflexe « comparaison pH / $14 + \log C$ » pour une base]
À \SI{25}{\celsius} ($K_e = 10^{-14}$), pour une solution de base de concentration $C$ :
\begin{enumerate}
\item Calculer $14 + \log C$ (pH théorique si la base était forte : $[\ce{HO-}] = C \Rightarrow \mathrm{pOH} = -\log C \Rightarrow \mathrm{pH} = 14 + \log C$).
\item Comparer au pH mesuré :
\begin{itemize}
\item si $\mathrm{pH} = 14 + \log C$ $\Rightarrow$ base \cle{forte} ;
\item si $\mathrm{pH} < 14 + \log C$ $\Rightarrow$ base \cle{faible} ($[\ce{HO-}] < C$).
\end{itemize}
\end{enumerate}
\end{methode}

\begin{exoresolu}[L'éthanoate de sodium est-il une base forte ?]
On dissout de l'éthanoate de sodium $\ce{CH3COONa}$ pour obtenir une solution de concentration $C = \SI{1,0e-1}{mol.L^{-1}}$. Le pH mesuré à \SI{25}{\celsius} est $8,9$.
\begin{enumerate}
\item Rappeler l'équation de la réaction de l'ion éthanoate avec l'eau.
\item Montrer qu'il s'agit d'une base faible.
\end{enumerate}
\tcblower
\textbf{1.} L'éthanoate de sodium est un sel fort : $\ce{CH3COONa -> Na+ + CH3COO-}$. L'ion $\ce{CH3COO-}$ réagit avec l'eau :
\[ \ce{CH3COO- + H2O <=> CH3COOH + HO-}. \]

\textbf{2.} Si l'ion éthanoate était une base forte, on aurait $[\ce{HO-}] = C = \SI{0,10}{mol.L^{-1}}$.
Le pH théorique serait : $\mathrm{pH}_{\text{th}} = 14 + \log(0,10) = 14 - 1 = 13$.
Or le pH mesuré est $8,9 < 13$. La concentration en ions $\ce{HO-}$ est bien inférieure à $C$ :
\[ [\ce{HO-}] = 10^{\mathrm{pH}-14} = 10^{-5,1} = \SI{7,9e-6}{mol.L^{-1}} \ll C. \]
L'ion éthanoate est une \cle{base faible} (sa réaction avec l'eau est un équilibre).
\end{exoresolu}

\courssub{Étudier l'équilibre d'ionisation par l'avancement : taux d'avancement final et coefficient d'ionisation}

\begin{definition}[Coefficient d'ionisation $\alpha$ (ou taux d'avancement final $\tau$)]
Pour un acide faible $\ce{AH}$ (ou une base faible $\ce{B}$) de concentration initiale $C$, le \cle{coefficient d'ionisation} $\alpha$ est le rapport entre la quantité d'acide (ou base) ionisée à l'équilibre et la quantité initiale introduite :
\[ \alpha = \frac{x_f}{x_{\text{max}}} = \frac{[\ce{H3O+}]_{\text{eq}}}{C} \quad\text{(acide)}\qquad\text{ou}\qquad \alpha = \frac{[\ce{HO-}]_{\text{eq}}}{C} \quad\text{(base)} \]
avec $0 < \alpha < 1$ (souvent exprimé en $\%$). $\alpha \ll 1$ caractérise un acide/base faible.
\end{definition}

\begin{methode}[Tableau d'avancement et calcul de $\alpha$]
Pour l'équilibre $\ce{AH + H2O <=> A- + H3O+}$ (concentration initiale $C$, volume $V$) :
\begin{enumerate}
\item Dresser le tableau d'avancement en quantités de matière ($n$ en mol) :
\begin{center}
\begin{tabularx}{\linewidth}{|l|X|X|X|}
\hline
\rowcolor{popblueL} État & \ce{AH} & \ce{A-} & \ce{H3O+} \\
\hline
Initial & $C V$ & $0$ & $0$ (négligeable) \\
\hline
Final & $CV - x_f$ & $x_f$ & $x_f$ \\
\hline
\end{tabularx}
\end{center}
\item Avancement maximal (réaction supposée totale) : $x_{\text{max}} = C V$.
\item Avancement final $x_f$ déduit du pH : $x_f = [\ce{H3O+}]_{\text{eq}} \times V = 10^{-\mathrm{pH}} \times V$.
\item Coefficient d'ionisation :
\[ \alpha = \frac{x_f}{x_{\text{max}}} = \frac{[\ce{H3O+}]_{\text{eq}}}{C} = \frac{10^{-\mathrm{pH}}}{C}. \]
\item Pour une base $\ce{B}$ : $\alpha = \dfrac{[\ce{HO-}]_{\text{eq}}}{C} = \dfrac{10^{\mathrm{pH}-14}}{C}$ (à \SI{25}{\celsius}).
\end{enumerate}
\end{methode}

\begin{propriete}[Loi de dilution d'Ostwald (notion)]
Pour un acide (ou une base) faible, \textbf{plus la solution est diluée (plus $C$ diminue), plus le coefficient d'ionisation $\alpha$ augmente}. La dilution déplace l'équilibre dans le sens de l'ionisation (sens qui augmente le nombre d'espèces en solution).
\end{propriete}

\begin{exoresolu}[Coefficient d'ionisation de l'ammoniac et effet de la dilution]
On dispose d'une solution $S_1$ d'ammoniac $\ce{NH3}$ de concentration $C_1 = \SI{1,0e-1}{mol.L^{-1}}$ ; son pH vaut $11,1$ à \SI{25}{\celsius}. On dilue $S_1$ dix fois pour obtenir $S_2$ ($C_2 = \SI{1,0e-2}{mol.L^{-1}}$) de pH $10,6$.
\begin{enumerate}
\item Écrire l'équation de la réaction de l'ammoniac avec l'eau.
\item Calculer $\alpha_1$ dans $S_1$ et $\alpha_2$ dans $S_2$.
\item Conclure sur l'effet de la dilution.
\end{enumerate}
\tcblower
\textbf{1. Équation de la réaction.}
\[ \ce{NH3_{(aq)} + H2O_{(l)} <=> NH4+_{(aq)} + HO-_{(aq)}}. \]

\textbf{2. Coefficients d'ionisation.}
Pour un volume $V$ de solution, le tableau d'avancement donne $[\ce{HO-}] = x_f/V$ et $x_{\text{max}} = CV$, d'où $\alpha = [\ce{HO-}]/C$ avec $[\ce{HO-}] = K_e/[\ce{H3O+}] = 10^{\mathrm{pH}-14}$.

\textbf{Solution $S_1$ ($C_1 = \SI{0,10}{mol.L^{-1}}$, $\mathrm{pH}_1 = 11,1$) :}
\[ [\ce{HO-}]_1 = 10^{11,1-14} = 10^{-2,9} = \SI{1,26e-3}{mol.L^{-1}} \]
\[ \alpha_1 = \frac{1,26\times10^{-3}}{1,0\times10^{-1}} = 1,26\times10^{-2} \approx \textbf{1,3\,\%}. \]

\textbf{Solution $S_2$ ($C_2 = \SI{0,010}{mol.L^{-1}}$, $\mathrm{pH}_2 = 10,6$) :}
\[ [\ce{HO-}]_2 = 10^{10,6-14} = 10^{-3,4} = \SI{3,98e-4}{mol.L^{-1}} \]
\[ \alpha_2 = \frac{3,98\times10^{-4}}{1,0\times10^{-2}} = 3,98\times10^{-2} \approx \textbf{4,0\,\%}. \]

\textbf{3. Conclusion.}
Dans les deux cas $\alpha \ll 1$ : l'ammoniac est une \cle{base faible}. De plus $\alpha_2 (4,0\,\%) > \alpha_1 (1,3\,\%)$ : la dilution \textbf{augmente} le coefficient d'ionisation. C'est la \cle{loi de dilution d'Ostwald}.
\end{exoresolu}

\begin{experience}[Détermination expérimentale de $K_a$ par pH-métrie (point à mi-équivalence)]
\textbf{Objectif :} Mesurer la constante d'acidité $K_a$ d'un acide faible (ex: acide éthanoïque) par dosage pH-métrique avec une base forte ($\ce{NaOH}$).
\par\textbf{Principe :} Au cours du dosage $\ce{AH + HO- -> A- + H2O}$, le pH évolue. Au \textbf{point à mi-équivalence} (volume versé $V = V_E/2$), on a $n_{\ce{AH}} = n_{\ce{A-}}$, donc $[\ce{AH}] = [\ce{A-}]$. La relation de Henderson-Hasselbalch donne $\mathrm{pH} = \mathrm{p}K_a + \log(1) = \mathrm{p}K_a$.
\par\textbf{Matériel :} pH-mètre étalonné (tampons pH 4, 7, 10), burette, erlenmeyer, agitateur magnétique, solution d'acide faible ($C_A, V_A$), solution de soude ($C_B$).
\par\textbf{Protocole :}
\begin{enumerate}
\item Verser $V_A$ de l'acide dans l'erlenmeyer, plonger l'électrode, noter $\mathrm{pH}_0$.
\item Ajouter la soude par petites fractions (ex: \SI{0,5}{mL}), noter le pH après chaque addition jusqu'après l'équivalence.
\item Tracer la courbe $\mathrm{pH} = f(V_B)$. Repérer le volume d'équivalence $V_E$ (point d'inflexion, ou tangentes).
\item Lire le pH au volume $V_E/2$ : cette valeur est $\mathrm{p}K_a$ de l'acide. En déduire $K_a = 10^{-\mathrm{p}K_a}$.
\end{enumerate}
\textbf{Exploitation :} Pour l'acide éthanoïque, on trouve $\mathrm{pH}_{1/2} \approx 4,8 \Rightarrow \mathrm{p}K_a = 4,8 \Rightarrow K_a = 1,6\times10^{-5}$.
\par\textbf{Précaution :} L'électrode doit être rincée à l'eau distillée entre chaque mesure ; l'agitation doit être constante sans créer de vortex.
\end{experience}

\courssub{Couples acide/base, couples de l'eau et ampholytes}

\begin{definition}[Couple acide/base (Brønsted)]
Un \cle{couple acide/base} est un duo d'espèces chimiques qui ne diffèrent que par la présence ou l'absence d'un \textbf{seul proton $\ce{H+}$}. On l'écrit \ce{AH}/\ce{A-} où $\ce{AH}$ est l'\textbf{acide} (donneur de $\ce{H+}$) et $\ce{A-}$ la \textbf{base conjuguée} (accepteur de $\ce{H+}$).
\[ \ce{AH <=> A- + H+} \]
\end{definition}

\begin{definition}[Ampholyte]
Une espèce chimique capable de se comporter \textbf{tantôt en acide, tantôt en base} selon le partenaire réactionnel est un \cle{ampholyte}. L'eau est l'ampholyte de référence.
\end{definition}

\begin{propriete}[Les deux couples de l'eau]
L'eau peut jouer les deux rôles :
\begin{itemize}
\item Comme \textbf{base} (accepte $\ce{H+}$) : couple \ce{H3O+}/\ce{H2O} (acide $\ce{H3O+}$, base $\ce{H2O}$).
\item Comme \textbf{acide} (donne $\ce{H+}$) : couple \ce{H2O}/\ce{HO-} (acide $\ce{H2O}$, base $\ce{HO-}$).
\end{itemize}
Le produit ionique de l'eau $K_e = [\ce{H3O+}][\ce{HO-}] = 10^{-14}$ à \SI{25}{\celsius} lie ces deux couples.
\end{propriete}

\begin{methode}[Identifier les couples dans un équilibre acido-basique]
\begin{enumerate}
\item Repérer l'espèce qui \textbf{cède} un $\ce{H+}$ : c'est l'acide du couple 1 ; l'espèce formée est sa base conjuguée.
\item Repérer l'espèce qui \textbf{capte} ce $\ce{H+}$ : c'est la base du couple 2 ; l'espèce formée est son acide conjugué.
\item Écrire les deux couples sous la forme \ce{Acide}/\ce{Base}.
\item Vérifier : les deux formes d'un même couple ne diffèrent que d'un $\ce{H+}$.
\item Ne pas oublier les couples de l'eau si $\ce{H2O}$, $\ce{H3O+}$ ou $\ce{HO-}$ interviennent.
\end{enumerate}
\end{methode}

\begin{exoresolu}[Couples en jeu dans trois équilibres]
Identifier les deux couples acide/base et le rôle de chaque réactif pour :
\[ \text{(1)}\ \ce{CH3COOH + H2O <=> CH3COO- + H3O+} \]
\[ \text{(2)}\ \ce{NH3 + H2O <=> NH4+ + HO-} \]
\[ \text{(3)}\ \ce{HCO3- + CH3COOH <=> H2CO3 + CH3COO-} \]
\tcblower
\textbf{(1)} $\ce{CH3COOH}$ perd $\ce{H+}$ $\rightarrow$ acide du couple \ce{CH3COOH}/\ce{CH3COO-}. $\ce{H2O}$ gagne $\ce{H+}$ $\rightarrow$ base du couple \ce{H3O+}/\ce{H2O}. \textit{Rôle de l'eau : base.}

\textbf{(2)} $\ce{NH3}$ gagne $\ce{H+}$ $\rightarrow$ base du couple \ce{NH4+}/\ce{NH3}. $\ce{H2O}$ perd $\ce{H+}$ $\rightarrow$ acide du couple \ce{H2O}/\ce{HO-}. \textit{Rôle de l'eau : acide.}

\textbf{(3)} $\ce{CH3COOH}$ cède $\ce{H+}$ $\rightarrow$ couple \ce{CH3COOH}/\ce{CH3COO-} (acide). $\ce{HCO3-}$ capte $\ce{H+}$ et devient $\ce{H2CO3}$ $\rightarrow$ base du couple \ce{H2CO3}/\ce{HCO3-}.

\textbf{Synthèse :} L'eau est base dans (1) et acide dans (2) : c'est un \cle{ampholyte}. L'ion hydrogénocarbonate $\ce{HCO3-}$ est base dans (3) mais peut aussi céder $\ce{H+}$ (couple \ce{HCO3-}/\ce{CO3^2-}) : c'est aussi un ampholyte, à l'origine du pouvoir tampon du sang et du jus de bissap.
\end{exoresolu}

\begin{savaistu}[Ampholytes et vie quotidienne au Cameroun]
Le pouvoir tampon du sang (maintien du pH vers 7,4) repose sur le couple \ce{H2CO3}/\ce{HCO3-} (dioxyde de carbone dissous / hydrogénocarbonate). Le jus de bissap (\textit{Hibiscus sabdariffa}) contient des anthocyanes (indicateurs naturels) et des acides organiques (citrique, malique) dont les formes acides/bases coexistent, stabilisant sa couleur et son acidité. Dans la savonnerie artisanale (huile de palme + lessive de cendre), la soude (base forte) réagit totalement, mais les résidus de carbonate ($\ce{CO3^2-}$, base faible, couple \ce{HCO3-}/\ce{CO3^2-}) modulent le pH final du savon.
\end{savaistu}

\courssub{Constantes d'acidité $K_a$, de basicité $K_b$ et $\mathrm{p}K_a$}

\begin{definition}[Constante d'acidité $K_a$ et $\mathrm{p}K_a$]
Pour un couple acide/base $\ce{AH}/\ce{A-}$ en solution aqueuse diluée (activités $\approx$ concentrations), la \cle{constante d'acidité} est :
\[ K_a = \frac{[\ce{A-}][\ce{H3O+}]}{[\ce{AH}]} \quad (\text{sans unité, concentrations en \SI{}{mol.L^{-1}}}). \]
La \cle{constante d'acidité logarithmique} est $\mathrm{p}K_a = -\log K_a$.
\end{definition}

\begin{definition}[Constante de basicité $K_b$ et $\mathrm{p}K_b$]
Pour une base $\ce{B}$ (couple $\ce{BH+}/\ce{B}$), la \cle{constante de basicité} est :
\[ K_b = \frac{[\ce{BH+}][\ce{HO-}]}{[\ce{B}]}, \quad \mathrm{p}K_b = -\log K_b. \]
\end{definition}

\begin{propriete}[Relation fondamentale $K_a \times K_b = K_e$]
Pour un couple conjugué $\ce{AH}/\ce{A-}$ à \SI{25}{\celsius} :
\[ K_a(\ce{AH}/\ce{A-}) \times K_b(\ce{A-}/\ce{AH}) = K_e = 10^{-14} \quad\Leftrightarrow\quad \mathrm{p}K_a + \mathrm{p}K_b = 14. \]
\textbf{Conséquence :} Plus l'acide est fort ($K_a$ grand, $\mathrm{p}K_a$ petit), plus sa base conjuguée est faible ($K_b$ petit, $\mathrm{p}K_b$ grand), et inversement.
\end{propriete}

\begin{aretenir}[Classement des couples sur l'axe des $\mathrm{p}K_a$]
On place les couples sur un axe de $\mathrm{p}K_a$ \textbf{croissant} (de 0 à 14 en milieu aqueux) :
\begin{itemize}
\item \textbf{Acides forts} : $\mathrm{p}K_a < 0$ (ex: \ce{H3O+}/\ce{H2O}, $\mathrm{p}K_a = -1,7$). Leurs bases conjuguées sont \emph{inertes} (ne réagissent pas avec l'eau).
\item \textbf{Acides faibles} : $0 < \mathrm{p}K_a < 14$ (ex: \ce{CH3COOH}/\ce{CH3COO-}, $\mathrm{p}K_a = 4,8$ ; \ce{NH4+}/\ce{NH3}, $\mathrm{p}K_a = 9,2$).
\item \textbf{Bases fortes} : $\mathrm{p}K_a > 14$ (ex: \ce{H2O}/\ce{HO-}, $\mathrm{p}K_a = 15,7$). Leurs acides conjugués sont inertes.
\end{itemize}
\textbf{Règle :} L'acide le plus fort est celui de \textbf{plus petit $\mathrm{p}K_a$} (en bas de l'axe). La base la plus forte est la base conjuguée de l'acide le plus faible (en haut de l'axe).
\end{aretenir}

\begin{exoresolu}[Classement de couples et constante de basicité]
Données à \SI{25}{\celsius} :
\begin{center}
\begin{tabularx}{\linewidth}{|l|X|X|}
\hline
\rowcolor{popblueL} Couple & $K_a$ & $\mathrm{p}K_a$ \\
\hline
\ce{HCOOH}/\ce{HCOO-} (acide méthanoïque) & $1,8\times10^{-4}$ & $3,7$ \\
\hline
\ce{CH3COOH}/\ce{CH3COO-} (acide éthanoïque) & $1,8\times10^{-5}$ & $4,8$ \\
\hline
\ce{NH4+}/\ce{NH3} (ion ammonium) & $6,3\times10^{-10}$ & $9,2$ \\
\hline
\end{tabularx}
\end{center}
\begin{enumerate}
\item Classer les trois acides par force croissante.
\item Classer les trois bases conjuguées par force croissante.
\item Calculer $K_b$ de l'ammoniac.
\end{enumerate}
\tcblower
\textbf{1. Force des acides (ordre croissant = $\mathrm{p}K_a$ croissant) :}
\[ \ce{NH4+} (\mathrm{p}K_a=9,2) < \ce{CH3COOH} (\mathrm{p}K_a=4,8) < \ce{HCOOH} (\mathrm{p}K_a=3,7). \]
L'acide méthanoïque est le plus fort des trois.

\textbf{2. Force des bases conjuguées (ordre inverse des acides) :}
\[ \ce{HCOO-} < \ce{CH3COO-} < \ce{NH3}. \]
L'ammoniac est la plus forte base ; l'ion éthanoate est plus basique que l'ion méthanoate.

\textbf{3. $K_b$ de l'ammoniac (base du couple \ce{NH4+}/\ce{NH3}) :}
\[ K_b = \frac{K_e}{K_a(\ce{NH4+}/\ce{NH3})} = \frac{10^{-14}}{6,3\times10^{-10}} = \SI{1,59e-5}{} \approx \SI{1,6e-5}{}. \]
Vérification : $\mathrm{p}K_b = 14 - 9,2 = 4,8 \Rightarrow K_b = 10^{-4,8} = 1,58\times10^{-5}$.
\end{exoresolu}

\courssub{Relation $\mathrm{pH}$--$\mathrm{p}K_a$ et domaines de prédominance}

\begin{propriete}[Relation de Henderson-Hasselbalch]
Pour un couple $\ce{AH}/\ce{A-}$ en solution, à l'équilibre :
\[ K_a = \frac{[\ce{A-}][\ce{H3O+}]}{[\ce{AH}]} \;\Rightarrow\; [\ce{H3O+}] = K_a\,\frac{[\ce{AH}]}{[\ce{A-}]}. \]
En prenant le $-\log$ des deux membres :
\[ \boxed{\mathrm{pH} = \mathrm{p}K_a + \log\frac{[\ce{A-}]}{[\ce{AH}]}} \quad\text{(base au numérateur, acide au dénominateur)}. \]
\end{propriete}

\begin{propriete}[Domaines de prédominance des formes acide et basique]
D'après la relation $\mathrm{pH} = \mathrm{p}K_a + \log\frac{[\text{base}]}{[\text{acide}]}$ :
\begin{itemize}
\item $\mathrm{pH} < \mathrm{p}K_a - 1$ $\Rightarrow$ $\log\frac{[\text{base}]}{[\text{acide}]} < -1$ $\Rightarrow$ $\frac{[\text{base}]}{[\text{acide}]} < 0,1$ : la forme \textbf{acide $\ce{AH}$ prédomine} ($[\ce{AH}] > 10[\ce{A-}]$).
\item $\mathrm{pH} > \mathrm{p}K_a + 1$ $\Rightarrow$ $\frac{[\text{base}]}{[\text{acide}]} > 10$ : la forme \textbf{basique $\ce{A-}$ prédomine} ($[\ce{A-}] > 10[\ce{AH}]$).
\item $\mathrm{pH} = \mathrm{p}K_a$ $\Rightarrow$ $[\ce{AH}] = [\ce{A-}]$ : \textbf{zone tampon} (capacité tampon maximale), zone de virage d'un indicateur coloré.
\end{itemize}
\end{propriete}

\begin{methode}[Choix d'un indicateur coloré pour un dosage]
Un indicateur coloré est un couple acide/base $\ce{InH}/\ce{In-}$ dont les deux formes ont des couleurs différentes. Sa zone de virage est $[\mathrm{p}K_a - 1\,;\, \mathrm{p}K_a + 1]$.
\begin{enumerate}
\item Déterminer (ou connaître) le pH à l'équivalence du dosage.
\item Choisir l'indicateur dont la \textbf{zone de virage contient ce pH d'équivalence}.
\end{enumerate}
\end{methode}

\begin{exoresolu}[Domaines de prédominance et choix d'un indicateur]
Couple \ce{CH3COOH}/\ce{CH3COO-} : $\mathrm{p}K_a = 4,8$.
\begin{enumerate}
\item Établir la relation $\mathrm{pH}$--$\mathrm{p}K_a$.
\item À $\mathrm{pH} = 5,8$, calculer $[\ce{CH3COO-}]/[\ce{CH3COOH}]$ et donner la forme prédominante.
\item Dosage d'acide éthanoïque par la soude : $\mathrm{pH}_{\text{équiv}} = 8,7$. Choisir entre hélianthine ($3,1$--$4,4$) et phénolphtaléine ($8,2$--$10,0$).
\end{enumerate}
\tcblower
\textbf{1.}
\[ K_a = \frac{[\ce{CH3COO-}][\ce{H3O+}]}{[\ce{CH3COOH}]} \Rightarrow \mathrm{pH} = \mathrm{p}K_a + \log\frac{[\ce{CH3COO-}]}{[\ce{CH3COOH}]}. \]

\textbf{2.}
\[ \log\frac{[\ce{CH3COO-}]}{[\ce{CH3COOH}]} = 5,8 - 4,8 = 1,0 \;\Rightarrow\; \frac{[\ce{CH3COO-}]}{[\ce{CH3COOH}]} = 10. \]
Il y a 10 fois plus de base que d'acide. Comme $\mathrm{pH} = \mathrm{p}K_a + 1$, on est à la limite du domaine de prédominance de la forme basique : \ce{CH3COO-} prédomine.

\textbf{3.} L'équivalence se situe à $\mathrm{pH} = 8,7$.
\begin{itemize}
\item Hélianthine : zone $3,1$--$4,4$. $8,7 \notin [3,1;4,4]$ $\rightarrow$ virage bien avant l'équivalence (erreur).
\item Phénolphtaléine : zone $8,2$--$10,0$. $8,7 \in [8,2;10,0]$ $\rightarrow$ \textbf{choix correct} (virage incolore $\to$ rose au voisinage de l'équivalence).
\end{itemize}

\begin{popfigure}
\begin{tikzpicture}[scale=1.1]
\draw[->,thick] (0,0) -- (10.5,0) node[right]{pH};
\foreach \x/\lab in {0/0, 3.38/4,8, 6.77/9,6, 9.86/14}
  \draw (\x,0.08) -- (\x,-0.08) node[below=2pt]{\lab};
% pKa = 4.8 -> x = 4.8 * 9.86/14 = 3.38
\fill[poppurple] (3.38,0) circle (2pt) node[above=3pt]{$\mathrm{p}K_a = 4,8$};
% Transition zone pKa-1 to pKa+1 -> pH 3.8 to 5.8 -> x = 2.68 to 4.08
\draw[poppurple,thick] (2.68,0) -- (4.08,0);
\fill[poppurple] (2.68,0) circle (1.5pt); \fill[poppurple] (4.08,0) circle (1.5pt);
\node[above] at (1.5,0.25){\small \ce{CH3COOH} prédomine};
\node[above] at (7.5,0.25){\small \ce{CH3COO-} prédomine};
\draw[poppurple,<->] (2.68,0.15) -- (4.08,0.15) node[midway,above,font=\small]{Zone de virage};
\end{tikzpicture}
\legende{Diagramme de prédominance du couple \ce{CH3COOH}/\ce{CH3COO-} ($\mathrm{p}K_a = 4,8$).}
\end{popfigure}
\end{exoresolu}

\courssub{Prévision des réactions acido-basiques}

\begin{methode}[Règle du gamma et constante de réaction $K$]
Pour prévoir le sens et l'étendue d'une réaction entre un acide et une base :
\begin{enumerate}
\item Identifier les deux couples acide/base mis en jeu : \ce{AH1}/\ce{A1-} et \ce{AH2}/\ce{A2-}.
\item Les placer sur un axe de $\mathrm{p}K_a$ croissant (acide à gauche, base à droite de chaque couple).
\item La réaction \textbf{spontanée (naturelle)} est celle entre l'\textbf{acide du couple de plus petit $\mathrm{p}K_a$} (le plus fort) et la \textbf{base du couple de plus grand $\mathrm{p}K_a$} (la plus forte) : \emph{la flèche « gamma » va de l'acide fort vers la base forte}.
\item Écrire l'équation : $\ce{AH1 + A2- <=> A1- + AH2}$.
\item Calculer la constante d'équilibre :
\[ \boxed{K = \frac{K_{a1}}{K_{a2}} = 10^{\mathrm{p}K_{a2} - \mathrm{p}K_{a1}}} \]
où $K_{a1}$ est le $K_a$ de l'acide \textbf{réactif} (celui qui donne $\ce{H+}$) et $K_{a2}$ celui de l'acide \textbf{formé} (acide conjugué de la base réactive).
\item Interpréter $K$ :
\begin{itemize}
\item $K \ge 10^4$ : réaction \textbf{quasi totale} (flèche \ce{->} souvent utilisée) ;
\item $K \le 10^{-4}$ : réaction \textbf{quasi nulle} ;
\item $10^{-4} < K < 10^4$ : \textbf{équilibre} significatif.
\end{itemize}
\end{enumerate}
\textbf{Astuce :} La réaction naturelle a toujours $K > 1$ (exposant $\mathrm{p}K_{a2}-\mathrm{p}K_{a1} > 0$). Si tu trouves $K < 1$, tu as écrit la réaction inverse.
\end{methode}

\begin{exoresolu}[Réaction entre l'acide éthanoïque et l'ammoniac]
Données : $\mathrm{p}K_a(\ce{CH3COOH}/\ce{CH3COO-}) = 4,8$ ; $\mathrm{p}K_a(\ce{NH4+}/\ce{NH3}) = 9,2$.
\begin{enumerate}
\item Écrire l'équation de la réaction qui se produit.
\item Calculer $K$ et conclure.
\item La réaction inverse (\ce{NH4+} + \ce{CH3COO-}) est-elle importante ?
\end{enumerate}
\tcblower
\textbf{1. Couples et sens naturel.}
Couple 1 : \ce{CH3COOH}/\ce{CH3COO-} ($\mathrm{p}K_{a1}=4,8$). Couple 2 : \ce{NH4+}/\ce{NH3} ($\mathrm{p}K_{a2}=9,2$).
L'acide le plus fort est \ce{CH3COOH} (plus petit $\mathrm{p}K_a$). La base la plus forte est \ce{NH3} (base du couple de plus grand $\mathrm{p}K_a$).
Réaction naturelle :
\[ \ce{CH3COOH + NH3 <=> CH3COO- + NH4+}. \]

\textbf{2. Constante d'équilibre.}
L'acide réactif (donneur de $\ce{H+}$) est \ce{CH3COOH} ($K_{a1}$). L'acide formé est \ce{NH4+} ($K_{a2}$).
\[ K = 10^{\mathrm{p}K_{a2} - \mathrm{p}K_{a1}} = 10^{9,2 - 4,8} = 10^{4,4} \approx 2,5\times10^{4}. \]
Comme $K \ge 10^4$, la réaction est \textbf{quasi totale}. Le mélange acide éthanoïque + ammoniac réagit pratiquement complètement pour former de l'éthanoate d'ammonium $\ce{CH3COONH4}$.

\textbf{3. Réaction inverse.}
\[ K' = \frac{1}{K} = 10^{-4,4} \approx 4,0\times10^{-5} \ll 10^{-4}. \]
La réaction entre \ce{NH4+} et \ce{CH3COO-} est \textbf{quasi nulle} : ces ions coexistent sans réagir notablement.
\end{exoresolu}

\courssub{Entraînement : exercices types Bac}

\begin{exercice}[Acide lactique dans le bili-bili]
L'acide lactique $\ce{CH3CH(OH)COOH}$ (noté $\ce{AH}$) est un acide faible présent dans le bili-bili. On prépare une solution $S$ de concentration $C = \SI{5,0e-2}{mol.L^{-1}}$. Le pH mesuré à \SI{25}{\celsius} est $2,52$.
\begin{enumerate}
\item Montrer que l'acide lactique est un acide faible.
\item Calculer le coefficient d'ionisation $\alpha$ de l'acide lactique dans $S$.
\item On dilue $S$ 10 fois. Prédire l'évolution de $\alpha$ et justifier.
\item Donner la valeur approchée du $\mathrm{p}K_a$ de l'acide lactique sachant qu'à l'équivalence d'un dosage par $\ce{NaOH}$, le pH vaut $8,2$ et qu'au point à mi-équivalence le pH vaut $3,8$.
\end{enumerate}
\end{exercice}

\begin{corrige}[Exercice : Acide lactique dans le bili-bili]
\textbf{1.} $-\log C = -\log(5,0\times10^{-2}) = 1,30$. pH mesuré $= 2,52 > 1,30$. L'acide lactique est un \textbf{acide faible} (ionisation partielle).
\par\textbf{2.} $[\ce{H3O+}] = 10^{-2,52} = \SI{3,02e-3}{mol.L^{-1}}$. $\alpha = \dfrac{[\ce{H3O+}]}{C} = \dfrac{3,02\times10^{-3}}{5,0\times10^{-2}} = 0,0604 \approx \textbf{6,0\,\%}$.
\par\textbf{3.} Par dilution ($C$ diminue), $\alpha$ \textbf{augmente} (loi d'Ostwald). L'équilibre se décale dans le sens de l'ionisation (augmentation du nombre d'espèces).
\par\textbf{4.} Au point à mi-équivalence, $[\ce{AH}] = [\ce{A-}]$, donc $\mathrm{pH} = \mathrm{p}K_a$. $\mathrm{p}K_a \approx \textbf{3,8}$ ($K_a \approx 1,6\times10^{-4}$). (Le pH à l'équivalence $8,2$ confirme que la base conjuguée $\ce{A-}$ est basique, mais ne donne pas directement $\mathrm{p}K_a$).
\end{corrige}

\begin{exercice}[Réaction entre acide éthanoïque et ion hydrogénocarbonate]
On mélange une solution d'acide éthanoïque $\ce{CH3COOH}$ et une solution d'hydrogénocarbonate de sodium $\ce{NaHCO3}$, apportant la même quantité de matière $n_0$ de chaque réactif. Données : $\mathrm{p}K_a(\ce{H2CO3}/\ce{HCO3-}) = 6,3$ ; $\mathrm{p}K_a(\ce{CH3COOH}/\ce{CH3COO-}) = 4,8$.
\begin{enumerate}
\item Identifier les deux couples acide/base mis en jeu.
\item Écrire l'équation de la réaction entre l'acide le plus fort et la base la plus forte.
\item Calculer la constante $K$ de cette réaction. Peut-on la considérer comme totale ?
\item Montrer que le taux d'avancement final vaut $\tau_f = \dfrac{\sqrt{K}}{1+\sqrt{K}}$ et le calculer.
\item En déduire le pH du mélange à l'équilibre.
\end{enumerate}
\end{exercice}

\begin{corrige}[Exercice : Acide éthanoïque et hydrogénocarbonate]
\textbf{1.} Couples : \ce{CH3COOH}/\ce{CH3COO-} ($\mathrm{p}K_{a1}=4,8$) et \ce{H2CO3}/\ce{HCO3-} ($\mathrm{p}K_{a2}=6,3$).
\par\textbf{2.} L'acide le plus fort est \ce{CH3COOH} (plus petit $\mathrm{p}K_a$) ; la base la plus forte est \ce{HCO3-} (plus grand $\mathrm{p}K_a$). Ce sont justement les deux réactifs mélangés :
\[ \ce{CH3COOH + HCO3- <=> CH3COO- + H2CO3} \]
\par\textbf{3.} $K = \dfrac{K_{a1}}{K_{a2}} = 10^{\mathrm{p}K_{a2} - \mathrm{p}K_{a1}} = 10^{6,3 - 4,8} = 10^{1,5} \approx 31,6$.
Comme $10^{-4} < K < 10^{4}$, la réaction n'est \textbf{pas totale} : elle conduit à un équilibre, mais elle est favorisée dans le sens direct ($K>1$).
\par\textbf{4.} Tableau d'avancement (quantités en mol) : à l'équilibre, $n(\ce{CH3COOH}) = n(\ce{HCO3-}) = n_0 - x_f$ et $n(\ce{CH3COO-}) = n(\ce{H2CO3}) = x_f$. Le volume se simplifie :
\[ K = \frac{x_f^2}{(n_0 - x_f)^2} \iff \frac{x_f}{n_0 - x_f} = \sqrt{K} \iff x_f = \frac{\sqrt{K}}{1+\sqrt{K}}\,n_0. \]
D'où $\tau_f = \dfrac{x_f}{n_0} = \dfrac{\sqrt{K}}{1+\sqrt{K}} = \dfrac{5,62}{6,62} \approx 0,85$ : environ $85\,\%$ des réactifs ont réagi.
\par\textbf{5.} Avec le couple \ce{CH3COOH}/\ce{CH3COO-} :
\[ \mathrm{pH} = \mathrm{p}K_{a1} + \log\frac{[\ce{CH3COO-}]}{[\ce{CH3COOH}]} = 4,8 + \log\frac{x_f}{n_0 - x_f} = 4,8 + \log\sqrt{K} = 4,8 + 0,75 \approx 5,55. \]
\textbf{Vérification} avec l'autre couple : $\mathrm{pH} = 6,3 + \log\dfrac{[\ce{HCO3-}]}{[\ce{H2CO3}]} = 6,3 - \log\sqrt{K} = 6,3 - 0,75 = 5,55$. Les deux couples donnent bien le même pH, ce qui est le cas pour toute solution à l'équilibre. On remarque que $\mathrm{pH} = \dfrac{\mathrm{p}K_{a1} + \mathrm{p}K_{a2}}{2}$.
\end{corrige}

\begin{exercice}[Dosage d'une solution d'acide éthanoïque]
On dose \SI{20,0}{mL} d'une solution d'acide éthanoïque de concentration inconnue $C_A$ par une solution de soude $\ce{NaOH}$ de concentration $C_B = \SI{0,100}{mol.L^{-1}}$. On suit le pH. Au point à mi-équivalence ($V_B = \SI{7,5}{mL}$), $\mathrm{pH} = 4,8$. À l'équivalence ($V_E = \SI{15,0}{mL}$), $\mathrm{pH} = 8,7$.
\begin{enumerate}
\item Déterminer $\mathrm{p}K_a$ et $K_a$ du couple \ce{CH3COOH}/\ce{CH3COO-}.
\item Calculer la concentration $C_A$ de la solution d'acide éthanoïque.
\item Justifier le choix de la phénolphtaléine comme indicateur pour ce dosage.
\item Écrire l'équation de la réaction au point d'équivalence et expliquer pourquoi le pH est basique ($>7$).
\end{enumerate}
\end{exercice}

\begin{corrige}[Exercice : Dosage acide éthanoïque / soude]
\textbf{1.} Au point à mi-équivalence, $n_{\ce{CH3COOH}} = n_{\ce{CH3COO-}}$, donc $[\ce{CH3COOH}] = [\ce{CH3COO-}]$.
$\mathrm{pH} = \mathrm{p}K_a + \log 1 = \mathrm{p}K_a = \textbf{4,8}$.
$K_a = 10^{-\num{4,8}} = \textbf{\num{1,6e-5}}$.

\textbf{2.} À l'équivalence, $n_{\ce{CH3COOH}}^0 = n_{\ce{NaOH}}$ versé.
$C_A \times V_A = C_B \times V_E \Rightarrow C_A = \frac{C_B V_E}{V_A} = \frac{\num{0,100} \times \num{15,0}}{\num{20,0}} = \SI{0,0750}{mol.L^{-1}}$.

\textbf{3.} Zone de virage phénolphtaléine : \num{8,2} à \num{10,0}. pH à l'équivalence $= \num{8,7} \in [\num{8,2}; \num{10,0}]$. L'indicateur vire au voisinage de l'équivalence : \textbf{choix adapté}. (Hélianthine \num{3,1}--\num{4,4} virait bien avant).

\textbf{4.} Équation : $\ce{CH3COOH + HO- -> CH3COO- + H2O}$ (quasi totale, $K = K_a/K_e = 10^{\num{9,2}}$).
À l'équivalence, la solution ne contient que l'ion éthanoate $\ce{CH3COO-}$ (base conjuguée d'un acide faible), qui réagit avec l'eau : $\ce{CH3COO- + H2O <=> CH3COOH + HO-}$. Cette réaction produit des ions $\ce{HO-}$, d'où $\mathrm{pH} > 7$ (basique).
\end{corrige}

\begin{attention}[Les 5 erreurs qui coûtent des points au Bac]
\begin{enumerate}
\item \textbf{Confondre le sens de l'inégalité pH / $-\log C$.} Acide faible $\Rightarrow$ $\mathrm{pH} > -\log C$ ; Base faible $\Rightarrow$ $\mathrm{pH} < 14 + \log C$. Écrire toujours les deux valeurs côte à côte avant de conclure.
\item \textbf{Oublier de diviser par $C$ dans $\alpha$.} $\alpha = \dfrac{[\ce{H3O+}]}{C}$ (acide) ou $\dfrac{[\ce{HO-}]}{C}$ (base). Pour une base, passer par $[\ce{HO-}] = 10^{\mathrm{pH}-14}$ (à \SI{25}{\celsius}).
\item \textbf{Inverser $K_a$ et $\mathrm{p}K_a$ dans les classements.} Grand $K_a$ = petit $\mathrm{p}K_a$ = acide \textbf{fort}. L'acide le plus fort est en \textbf{bas} de l'axe $\mathrm{p}K_a$ croissant.
\item \textbf{Mal écrire Henderson-Hasselbalch.} $\mathrm{pH} = \mathrm{p}K_a + \log\dfrac{[\text{base}]}{[\text{acide}]}$ : \textbf{base au numérateur}. Inversion = domaines de prédominance inversés = mauvais indicateur.
\item \textbf{Se tromper d'indices dans $K = 10^{\mathrm{p}K_{a2}-\mathrm{p}K_{a1}}$.} Indice 1 = couple de l'acide \textbf{réactif} (donneur $\ce{H+}$). Indice 2 = couple de l'acide \textbf{formé} (acide conjugué de la base réactive). Si $K < 1$, tu as écrit la réaction non naturelle : inverse l'équation.
\end{enumerate}
\end{attention}

\newpage

\coursec{Exercices}

\coursec{Leçon 08 : Acides faibles, bases faibles et couples acide/base}

\courssub{1. Acide faible : l'acide éthanoïque}

\begin{exemplebox}[Le vinaigre, un acide faible du quotidien]
Au Cameroun, le vinaigre de palme ou le vinaigre d'alcool (souvent à 8\% massique, soit environ \SI{1,3}{mol.L^{-1}} en acide éthanoïque) sert à la conservation des aliments (pickles, condiments) et au nettoyage. Si l'on mesure le pH d'une solution d'acide éthanoïque \ce{CH3COOH} de concentration \SI{1,0e-1}{mol.L^{-1}}, on trouve \SI{2,9}{\approx}, alors que \SI{-log(1,0e-1)}{= 1,0}. Le pH est \textbf{bien plus élevé} que celui prévu pour un acide fort de même concentration. L'acide éthanoïque ne libère pas tous ses protons : c'est un \cle{acide faible}.
\end{exemplebox}

\begin{definition}[Acide faible]
Un \cle{acide faible} est une espèce chimique qui, en solution aqueuse, ne cède qu'une \textbf{faible fraction} de ses protons \ce{H+} (ions \ce{H3O+}). Sa réaction avec l'eau est \textbf{limitée} et conduit à un \cle{équilibre chimique}.
\end{definition}

\begin{experience}[Mise en évidence de l'ionisation partielle]
\begin{enumerate}
\item Préparer une solution $S$ d'acide éthanoïque de concentration connue $C = \SI{1,0e-2}{mol.L^{-1}}$.
\item Mesurer son pH à \SI{25}{\celsius} : $\mathrm{pH} \approx 3,4$.
\item Calculer $[\ce{H3O+}] = 10^{-\mathrm{pH}} = \SI{4,0e-4}{mol.L^{-1}}$.
\item Comparer : $[\ce{H3O+}] = \SI{4,0e-4}{mol.L^{-1}} \ll C = \SI{1,0e-2}{mol.L^{-1}}$.
\end{enumerate}
\legende{Seule une petite fraction des molécules \ce{CH3COOH} s'est ionisée.}
\end{experience}

\begin{propriete}[Équation de l'équilibre d'ionisation]
L'acide éthanoïque réagit avec l'eau selon l'équilibre :
\[ \ce{CH3COOH_{(aq)} + H2O_{(l)} <=> CH3COO-_{(aq)} + H3O+_{(aq)}} \]
Les deux \cle{couples acide/base} mis en jeu sont :
\begin{itemize}
\item \ce{CH3COOH}/\ce{CH3COO-} (couple de l'acide éthanoïque) ;
\item \ce{H3O+}/\ce{H2O} (couple de l'eau, acide fort).
\end{itemize}
\end{propriete}

\begin{methode}[Étude de l'équilibre par le tableau d'avancement]
Pour un volume $V$ de solution de concentration initiale $C$ :
\begin{enumerate}
\item Écrire l'équation de l'équilibre.
\item Dresser le tableau d'avancement (quantités de matière en mol) :
\begin{center}
\begin{tabularx}{\linewidth}{|c|X|X|X|}\hline
 & \ce{CH3COOH} & \ce{H2O} & \ce{CH3COO-} + \ce{H3O+} \\\hline
État initial & $n_i = C \times V$ & Excès & 0 \\\hline
État final & $n_i - x_f$ & Excès & $x_f$ \\\hline
\end{tabularx}
\end{center}
\item Exprimer les concentrations à l'équilibre : $[\ce{CH3COOH}] = \frac{C V - x_f}{V} = C - \frac{x_f}{V}$, $[\ce{CH3COO-}] = [\ce{H3O+}] = \frac{x_f}{V}$.
\item Le \cle{coefficient d'ionisation} (ou \cle{taux d'avancement final}) est $\alpha = \tau = \frac{x_f}{C V} = \frac{[\ce{H3O+}]}{C}$.
\item La \cle{constante d'acidité} $K_a$ du couple \ce{CH3COOH}/\ce{CH3COO-} s'écrit :
\[ K_a = \frac{[\ce{CH3COO-}][\ce{H3O+}]}{[\ce{CH3COOH}]} = \frac{(C\alpha)^2}{C(1-\alpha)} = \frac{C\alpha^2}{1-\alpha} \]
Si $\alpha \ll 1$ (souvent $\alpha < 5\%$), $K_a \approx C\alpha^2$.
\end{enumerate}
\end{methode}

\begin{exemplebox}[Calcul de $K_a$ et $\alpha$ pour l'acide éthanoïque]
Données : $C = \SI{1,0e-2}{mol.L^{-1}}$, $\mathrm{pH} = 3,4$ ($\Rightarrow [\ce{H3O+}] = \SI{4,0e-4}{mol.L^{-1}}$).
\begin{align*}
\alpha &= \frac{[\ce{H3O+}]}{C} = \frac{4,0\times 10^{-4}}{1,0\times 10^{-2}} = 0,040 = 4,0\% \\
K_a &= \frac{[\ce{H3O+}]^2}{C - [\ce{H3O+}]} = \frac{(4,0\times 10^{-4})^2}{1,0\times 10^{-2} - 4,0\times 10^{-4}} = \frac{1,6\times 10^{-7}}{9,6\times 10^{-3}} = \SI{1,7e-5}{}\\
\mathrm{p}K_a &= -\log K_a = 4,77 \approx 4,8 \quad \text{(valeur tabulée)}
\end{align*}
\end{exemplebox}

\begin{propriete}[Influence de la dilution — Loi d'Ostwald (notion)]
Lorsqu'on dilue une solution d'acide faible (diminution de $C$) :
\begin{itemize}
\item L'avancement final $x_f$ (en mol) \textbf{augmente} (l'équilibre se décale vers les produits, principe de Le Chatelier).
\item Le coefficient d'ionisation $\alpha = x_f/(CV)$ \textbf{augmente}.
\item Le pH \textbf{augmente} (la solution devient moins acide), mais moins vite que pour un acide fort.
\end{itemize}
\emph{Relation approchée (loi d'Ostwald) :} $\alpha \approx \sqrt{K_a/C}$.
\end{propriete}

\begin{attention}[Pièges fréquents]
\begin{itemize}
\item Ne pas confondre \textbf{avancement final} $x_f$ (en mol, dépend du volume) et \textbf{coefficient d'ionisation} $\alpha$ (sans unité, intensif).
\item Ne pas écrire $\ce{CH3COOH -> CH3COO- + H+}$ (flèche simple) : la réaction est \textbf{réversible} (\ce{<=>}).
\item L'approximation $K_a \approx C\alpha^2$ n'est valable que si $\alpha < 5\%$. Sinon, utiliser la formule exacte.
\end{itemize}
\end{attention}

\courssub{2. Base faible : éthanoate de sodium et ammoniac}

\begin{savaistu}[Lessive artisanale et ammoniaque]
Au marché, le savon mou (à base d'huile de palme et de soude) contient de l'éthanoate de sodium résiduel. L'ammoniac \ce{NH3} (gaz dissous dans l'eau, « eau d'ammoniaque ») sert de détachant. Ces deux solutions ont un pH basique (\SI{> 7}{}) mais ne sont pas des bases fortes.
\end{savaistu}

\begin{definition}[Base faible]
Une \cle{base faible} est une espèce qui ne capte qu'une \textbf{faible fraction} des protons de l'eau. Sa réaction avec l'eau est \textbf{limitée} et conduit à un équilibre.
\end{definition}

\paragraph{Éthanoate de sodium \ce{CH3COONa}.} Sel totalement dissocié : \ce{CH3COONa -> Na+ + CH3COO-}. L'anion \ce{CH3COO-} (base conjuguée de l'acide éthanoïque) réagit avec l'eau (hydrolyse) :
\[ \ce{CH3COO-_{(aq)} + H2O_{(l)} <=> CH3COOH_{(aq)} + HO-_{(aq)}} \]
Couples : \ce{CH3COOH}/\ce{CH3COO-} et \ce{H2O}/\ce{HO-}.
La \cle{constante de basicité} $K_b$ s'écrit : $K_b = \frac{[\ce{CH3COOH}][\ce{HO-}]}{[\ce{CH3COO-}]}$.
Relation fondamentale : $K_a(\ce{CH3COOH}/\ce{CH3COO-}) \times K_b(\ce{CH3COOH}/\ce{CH3COO-}) = K_e = 10^{-14}$ (à \SI{25}{\celsius}).

\paragraph{Ammoniac \ce{NH3}.}
\[ \ce{NH3_{(aq)} + H2O_{(l)} <=> NH4+_{(aq)} + HO-_{(aq)}} \]
Couples : \ce{NH4+}/\ce{NH3} et \ce{H2O}/\ce{HO-}.
$K_b = \frac{[\ce{NH4+}][\ce{HO-}]}{[\ce{NH3}]}$. $K_a(\ce{NH4+}/\ce{NH3}) \times K_b = K_e$.
$\mathrm{p}K_a(\ce{NH4+}/\ce{NH3}) = 9,2$ (tabulée) $\Rightarrow \mathrm{p}K_b = 14 - 9,2 = 4,8$.

\begin{aretenir}[Force acide/basique conjuguées]
\begin{itemize}
\item Plus l'acide \ce{AH} est \textbf{fort} (grand $K_a$, petit $\mathrm{p}K_a$), plus sa base conjuguée \ce{A-} est \textbf{faible} (petit $K_b$, grand $\mathrm{p}K_b$).
\item Plus la base \ce{B} est \textbf{forte} (grand $K_b$), plus son acide conjugué \ce{BH+} est \textbf{faible}.
\end{itemize}
\end{aretenir}

\courssub{3. Couples acide/base, couples de l'eau et ampholytes}

\begin{definition}[Couple acide/base]
Un \cle{couple acide/base} est un duo d'espèces chimiques \ce{AH}/\ce{A-} qui se transforment l'une en l'autre par \textbf{perte ou gain d'un proton \ce{H+}} (ion \ce{H3O+} en milieu aqueux).
\emph{Demi-équation :} $\ce{AH <=> A- + H+}$ (ou $\ce{AH + H2O <=> A- + H3O+}$).
\ce{AH} = forme acide (donneur de proton) ; \ce{A-} = forme basique (accepteur de proton).
\end{definition}

\begin{propriete}[Couples de l'eau]
L'eau est un \cle{ampholyte} (elle peut agir comme acide ou comme base) :
\begin{itemize}
\item Comme acide : \ce{H2O <=> HO- + H+} \quad Couple \ce{H2O}/\ce{HO-} \quad ($K_a = K_e = 10^{-14}$, $\mathrm{p}K_a = 14$).
\item Comme base : \ce{H3O+ <=> H2O + H+} \quad Couple \ce{H3O+}/\ce{H2O} \quad ($K_a \approx 1$, $\mathrm{p}K_a \approx 0$).
\end{itemize}
\end{propriete}

\begin{definition}[Ampholyte]
Une espèce \cle{ampholyte} peut \textbf{céder} un proton (agir comme acide) \textbf{et} \textbf{capter} un proton (agir comme base).
\emph{Exemples :} \ce{H2O}, \ce{HCO3-}, \ce{H2PO4-}, \ce{HPO4^2-}, \ce{NH3} (très faiblement acide).
\end{definition}

\begin{exemplebox}[L'ion hydrogénocarbonate \ce{HCO3-}]
\begin{align*}
\text{Comme base :} &\quad \ce{HCO3- + H2O <=> H2CO3 + HO-} \quad (\text{couple } \ce{H2CO3}/\ce{HCO3-}) \\
\text{Comme acide :} &\quad \ce{HCO3- + H2O <=> CO3^2- + H3O+} \quad (\text{couple } \ce{HCO3-}/\ce{CO3^2-})
\end{align*}
\ce{HCO3-} est bien un ampholyte. Dans le sang humain (pH \approx 7,4), c'est le principal tampon (avec \ce{H2CO3}).
\end{exemplebox}

\begin{popfigure}
\begin{tikzpicture}[scale=0.9, every node/.style={font=\small}]
% Axe pH
\draw[->, thick] (0,0) -- (14,0) node[right] {pH};
\foreach \x in {0,2,4,6,8,10,12,14} \draw (\x,0.1) -- (\x,-0.1) node[below] {\x};
% Zones prédominance H2CO3 / HCO3- / CO3^2-
\fill[popblue!15] (0,0.5) rectangle (6.4,1.5) node[midway, align=center] {\ce{H2CO3}\\\textbf{prédomine}};
\fill[popgreen!15] (6.4,0.5) rectangle (10.3,1.5) node[midway, align=center] {\ce{HCO3-}\\\textbf{prédomine}};
\fill[poporange!15] (10.3,0.5) rectangle (14,1.5) node[midway, align=center] {\ce{CO3^2-}\\\textbf{prédomine}};
% Lignes pKa
\draw[dashed, popblue] (6.4,0) -- (6.4,1.5) node[above] {$\mathrm{p}K_{a1}=6,4$};
\draw[dashed, poporange] (10.3,0) -- (10.3,1.5) node[above] {$\mathrm{p}K_{a2}=10,3$};
% Point sang
\draw[popred, thick] (7.4, -0.3) -- (7.4, 1.7) node[above] {Sang (pH 7,4)};
\draw[popred, fill=popred] (7.4, 1.0) circle (2pt);
\end{tikzpicture}
\legende{Diagramme de prédominance du système carbonaté (échelle de pH).}
\end{popfigure}

\courssub{4. Constantes d'acidité $K_a$, de basicité $K_b$ et $\mathrm{p}K_a$}

\begin{propriete}[Constante d'acidité $K_a$]
Pour le couple \ce{AH}/\ce{A-} : $\ce{AH + H2O <=> A- + H3O+}$.
\[ K_a = \frac{[\ce{A-}][\ce{H3O+}]}{[\ce{AH}]} \qquad \mathrm{p}K_a = -\log K_a \]
À \SI{25}{\celsius} : $K_e = [\ce{H3O+}][\ce{HO-}] = 10^{-14}$.
\end{propriete}

\begin{propriete}[Constante de basicité $K_b$]
Pour la réaction de la base \ce{A-} avec l'eau : $\ce{A- + H2O <=> AH + HO-}$.
\[ K_b = \frac{[\ce{AH}][\ce{HO-}]}{[\ce{A-}]} \qquad \mathrm{p}K_b = -\log K_b \]
\textbf{Relation fondamentale :} $K_a \times K_b = K_e = 10^{-14}$ \quad $\Rightarrow$ \quad $\mathrm{p}K_a + \mathrm{p}K_b = 14$ (à \SI{25}{\celsius}).
\end{propriete}

\begin{aretenir}[Classification des couples (à \SI{25}{\celsius})]
\begin{center}
\begin{tabularx}{\linewidth}{|l|X|X|X|}\hline
\rowcolor{popblueL}
\textbf{Couple \ce{AH}/\ce{A-}} & $\mathrm{p}K_a$ & \textbf{Force de l'acide \ce{AH}} & \textbf{Force de la base \ce{A-}} \\\hline
\ce{H3O+}/\ce{H2O} & $\approx 0$ & Très fort & Nulle (base spectatrice) \\\hline
\ce{HCOOH}/\ce{HCOO-} & 3,8 & Fort (pour un acide faible) & Faible \\\hline
\ce{CH3COOH}/\ce{CH3COO-} & 4,8 & Modéré & Modérée \\\hline
\ce{H2CO3}/\ce{HCO3-} & 6,4 & Faible & Modérée \\\hline
\ce{H2O}/\ce{HO-} & 14 & Très faible (négligeable) & Forte \\\hline
\ce{NH4+}/\ce{NH3} & 9,2 & Très faible & Forte \\\hline
\end{tabularx}
\end{center}
\emph{Règle :} $\mathrm{p}K_a$ \textbf{croissant} $\Leftrightarrow$ Force acide \textbf{décroissante} $\Leftrightarrow$ Force basique conjuguée \textbf{croissante}.
\end{aretenir}

\courssub{5. Relation $\mathrm{pH}$-$\mathrm{p}K_a$ et domaines de prédominance}

\begin{methode}[Démonstration de la relation de Henderson-Hasselbalch]
\begin{enumerate}
\item Écrire $K_a = \frac{[\ce{A-}][\ce{H3O+}]}{[\ce{AH}]}$.
\item Isoler $[\ce{H3O+}]$ : $[\ce{H3O+}] = K_a \frac{[\ce{AH}]}{[\ce{A-}]}$.
\item Prendre le $-\log$ : $\mathrm{pH} = \mathrm{p}K_a - \log\frac{[\ce{AH}]}{[\ce{A-}]}$.
\item Utiliser $-\log\frac{a}{b} = \log\frac{b}{a}$ :
\[ \boxed{\mathrm{pH} = \mathrm{p}K_a + \log\frac{[\ce{A-}]}{[\ce{AH}]}} \]
\end{enumerate}
\end{methode}

\begin{propriete}[Domaines de prédominance]
Pour un couple \ce{AH}/\ce{A-} en solution :
\begin{itemize}
\item Si $\mathrm{pH} < \mathrm{p}K_a$ : $\log\frac{[\ce{A-}]}{[\ce{AH}]} < 0 \Rightarrow [\ce{AH}] > [\ce{A-}]$ \quad \textbf{Forme acide prédominante}.
\item Si $\mathrm{pH} = \mathrm{p}K_a$ : $[\ce{AH}] = [\ce{A-}]$ \quad \textbf{Formes équimolaires}.
\item Si $\mathrm{pH} > \mathrm{p}K_a$ : $[\ce{A-}] > [\ce{AH}]$ \quad \textbf{Forme basique prédominante}.
\end{itemize}
\end{propriete}

\begin{exemplebox}[Indicateur coloré : Bleu de bromothymol (BBT)]
Couple \ce{HIn}/\ce{In-} : $\mathrm{p}K_a = 7,0$. \ce{HIn} (jaune), \ce{In-} (bleu).
\begin{itemize}
\item pH < 6,0 : $[\ce{HIn}] > 10[\ce{In-}]$ $\Rightarrow$ \textbf{Jaune}.
\item pH > 8,0 : $[\ce{In-}] > 10[\ce{HIn}]$ $\Rightarrow$ \textbf{Bleu}.
\item 6,0 < pH < 8,0 : \textbf{Zone de virage} (vert, mélange jaune+bleu).
\end{itemize}
\end{exemplebox}

\begin{popfigure}
\begin{tikzpicture}[scale=0.8, every node/.style={font=\small}]
% Axe pH
\draw[->, thick] (0,0) -- (14,0) node[right] {pH};
\foreach \x in {0,2,4,6,8,10,12,14} \draw (\x,0.1) -- (\x,-0.1) node[below] {\x};
% Zones BBT
\fill[yellow!60] (0,0.5) rectangle (6,1.5) node[midway] {\textbf{Jaune} (\ce{HIn})};
\fill[green!40!yellow] (6,0.5) rectangle (8,1.5) node[midway] {\textbf{Vert} (mélange)};
\fill[blue!40] (8,0.5) rectangle (14,1.5) node[midway] {\textbf{Bleu} (\ce{In-})};
% pKa
\draw[dashed, thick, popdark] (7,0) -- (7,1.5) node[above] {$\mathrm{p}K_a = 7,0$};
% Flèches
\draw[<->, thick] (6,1.7) -- (8,1.7) node[midway, above] {Zone de virage};
\end{tikzpicture}
\legende{Zone de virage du bleu de bromothymol (BBT).}
\end{popfigure}

\courssub{6. Prévision des réactions acido-basiques}

\begin{propriete}[Constante de réaction acido-basique]
Réaction entre un acide \ce{AH1} (couple 1, $\mathrm{p}K_{a1}$) et une base \ce{A2-} (couple 2, $\mathrm{p}K_{a2}$) :
\[ \ce{AH1 + A2- <=> A1- + AH2} \]
\[ K = \frac{K_{a1}}{K_{a2}} = 10^{\mathrm{p}K_{a2} - \mathrm{p}K_{a1}} \]
\end{propriete}

\begin{aretenir}[Critère de réaction « totale » (quasi complète)]
\begin{itemize}
\item Si $K > 10^4$ (soit $\mathrm{p}K_{a2} - \mathrm{p}K_{a1} > 4$) : la réaction est \textbf{totale} (avancement final $\approx$ quantité initiale du réactif limitant).
\item Si $K < 10^{-4}$ : la réaction \textbf{ne se produit pas} (ou très peu).
\item Entre les deux : équilibre avec avancement significatif mais non total.
\end{itemize}
\end{aretenir}

\begin{exemplebox}[Réaction acide méthanoïque / ammoniac]
$\mathrm{p}K_{a1}(\ce{HCOOH}/\ce{HCOO-}) = 3,8$ ; $\mathrm{p}K_{a2}(\ce{NH4+}/\ce{NH3}) = 9,2$.
$K = 10^{9,2 - 3,8} = 10^{5,4} \approx \SI{2,5e5}{} \gg 10^4$.
La réaction $\ce{HCOOH + NH3 -> HCOO- + NH4+}$ est \textbf{totale}.
\end{exemplebox}

\begin{savaistu}[Applications au Cameroun]
\begin{itemize}
\item \textbf{Bili-bili / Vin de palme :} Fermentation lactique (acide lactique, $\mathrm{p}K_a \approx 3,8$) et acétique. L'acidité (pH \approx 3,5-4) assure la conservation.
\item \textbf{Savonnerie artisanale :} Hydrolyse basique des triglycérides (huile de palme) par la soude \ce{NaOH} (base forte). L'excès de soude est neutralisé par addition d'acide faible (jus de citron, $\mathrm{p}K_a \approx 3,1$) pour obtenir un savon doux (pH \approx 8-9).
\item \textbf{Eaux minérales (Mont Cameroun) :} Riches en \ce{HCO3-} (ampholyte). Le pH (souvent 6-7) est tamponné par le couple \ce{H2CO3}/\ce{HCO3-}.
\item \textbf{Industrie (SONARA, CIMENCAM) :} Contrôle du pH des effluents par ajout d'acide/base faible pour éviter la corrosion ou la précipitation intempestive.
\end{itemize}
\end{savaistu}

\courssub{Je vérifie mes connaissances}

\begin{exercice}[QCM : réaction acido-basique]
Pour chaque question, une seule réponse est exacte. Entoure-la et justifie brièvement.
\begin{enumerate}
\item Une solution d'acide éthanoïque de concentration $C = \SI{1,0e-2}{mol.L^{-1}}$ a un pH de $3,4$. On peut affirmer que :
\begin{enumerate}
\item l'acide éthanoïque est un acide fort ;
\item l'acide éthanoïque est un acide faible car $\mathrm{pH} > -\log C$ ;
\item la mesure de pH est fausse car un acide a toujours $\mathrm{pH} = -\log C$.
\end{enumerate}
\item Le coefficient d'ionisation $\alpha$ d'un acide faible :
\begin{enumerate}
\item est toujours égal à 1 ;
\item est compris entre 0 et 1 et augmente quand on dilue la solution ;
\item diminue quand on dilue la solution.
\end{enumerate}
\item Le couple \ce{NH4+}/\ce{NH3} a un $\mathrm{p}K_a = 9,2$. Dans une solution de $\mathrm{pH} = 7,0$ :
\begin{enumerate}
\item la forme basique \ce{NH3} prédomine ;
\item la forme acide \ce{NH4+} prédomine ;
\item les deux formes sont en concentrations égales.
\end{enumerate}
\item La constante de la réaction entre un acide \ce{A1H} ($\mathrm{p}K_{a1}$) et une base \ce{A2-} ($\mathrm{p}K_{a2}$) vaut :
\begin{enumerate}
\item $K = 10^{\mathrm{p}K_{a1} - \mathrm{p}K_{a2}}$ ;
\item $K = 10^{\mathrm{p}K_{a2} - \mathrm{p}K_{a1}}$ ;
\item $K = \mathrm{p}K_{a2} - \mathrm{p}K_{a1}$.
\end{enumerate}
\end{enumerate}
\end{exercice}

\begin{exercice}[Vrai ou faux ?]
Réponds par vrai ou faux en justifiant chaque réponse par une phrase ou un calcul court.
\begin{enumerate}
\item Une base faible réagit totalement avec l'eau.
\item Pour un couple acide/base, on a toujours $K_a \times K_b = K_e$ à \SI{25}{\celsius}.
\item Plus le $\mathrm{p}K_a$ d'un couple est grand, plus l'acide du couple est fort.
\item L'ion hydrogénocarbonate \ce{HCO3-} est un ampholyte.
\item Lorsqu'on dilue une solution d'acide faible, son taux d'avancement final diminue.
\item À $\mathrm{pH} = \mathrm{p}K_a$, les concentrations des formes acide et basique d'un couple sont égales.
\end{enumerate}
\end{exercice}

\begin{exercice}[Texte à trous]
Complète le texte suivant avec les mots ou expressions de la liste : \emph{partielle ; équilibre ; $\mathrm{p}K_a$ ; avancement ; ampholyte ; prédomine ; $K_e$ ; totale}.

Un acide faible ne s'ionise que de façon \ldots\ldots\ldots dans l'eau : sa dissolution conduit à un \ldots\ldots\ldots chimique que l'on étudie à l'aide du tableau d'\ldots\ldots\ldots. La force d'un acide se mesure par la constante d'acidité $K_a$ de son couple, ou plus souvent par son \ldots\ldots\ldots. Une espèce comme \ce{HCO3-}, capable de céder ou de capter un proton, est appelée \ldots\ldots\ldots. Pour un couple donné, la forme acide \ldots\ldots\ldots lorsque $\mathrm{pH} < \mathrm{p}K_a$. À \SI{25}{\celsius}, le produit $K_a \times K_b$ d'un couple est égal à \ldots\ldots\ldots, soit $10^{-14}$. Une réaction acido-basique dont la constante $K$ est supérieure à $10^{4}$ est considérée comme \ldots\ldots\ldots.
\end{exercice}

\begin{exercice}[Définitions et couples]
\begin{enumerate}
\item Définis : acide faible, base faible, couple acide/base, constante d'acidité $K_a$.
\item Écris les demi-équations acido-basiques des couples suivants et identifie la forme acide et la forme basique : \ce{HCOOH}/\ce{HCOO-} ; \ce{NH4+}/\ce{NH3} ; \ce{H2O}/\ce{HO-}.
\item Rappelle la relation de Henderson-Hasselbalch reliant le pH d'une solution au $\mathrm{p}K_a$ du couple \ce{AH}/\ce{A-} qu'elle contient.
\end{enumerate}
\end{exercice}

\courssub{Je m'entraîne}

\begin{exercice}[L'ion hydrogénocarbonate, un ampholyte]
On considère les deux équilibres suivants en solution aqueuse :
\[ \ce{HCO3- + H2O <=> H2CO3 + HO-} \qquad \ce{HCO3- + H2O <=> CO3^2- + H3O+} \]
\begin{enumerate}
\item Identifie tous les couples acide/base mis en jeu dans chacun de ces équilibres.
\item Montre, à partir de tes réponses, que l'ion \ce{HCO3-} est un ampholyte.
\item Sachant que $\mathrm{p}K_a(\ce{H2CO3}/\ce{HCO3-}) = 6,4$ et $\mathrm{p}K_a(\ce{HCO3-}/\ce{CO3^2-}) = 10,3$, indique la forme prédominante de l'élément carbone dans une eau minérale de $\mathrm{pH} = 7,2$.
\end{enumerate}
\end{exercice}

\begin{exercice}[Classement de couples et prévision d'une réaction]
On donne les couples suivants : \ce{HCOOH}/\ce{HCOO-} ($\mathrm{p}K_a = 3,8$) ; \ce{CH3COOH}/\ce{CH3COO-} ($\mathrm{p}K_a = 4,8$) ; \ce{NH4+}/\ce{NH3} ($\mathrm{p}K_a = 9,2$).
\begin{enumerate}
\item Classe les acides de ces couples par force acide croissante. Justifie.
\item Classe les bases conjuguées par force basique croissante.
\item Écris l'équation de la réaction entre l'acide méthanoïque \ce{HCOOH} et l'ammoniac \ce{NH3}.
\item Calcule la constante d'équilibre $K$ de cette réaction. Cette réaction est-elle totale ? Conclus.
\end{enumerate}
\end{exercice}

\begin{exercice}[Mélange acide éthanoïque / éthanoate de sodium]
Une solution tampon contient de l'acide éthanoïque et de l'éthanoate de sodium. Son pH mesuré vaut $5,8$. On donne $\mathrm{p}K_a(\ce{CH3COOH}/\ce{CH3COO-}) = 4,8$.
\begin{enumerate}
\item Écris la relation de Henderson-Hasselbalch pour ce couple.
\item Calcule le rapport $\dfrac{[\ce{CH3COO-}]}{[\ce{CH3COOH}]}$ dans cette solution.
\item Quelle est la forme prédominante du couple ? Justifie.
\item Trace le diagramme de prédominance du couple sur un axe de pH et place la solution étudiée.
\end{enumerate}
\end{exercice}

\begin{exercice}[Le bleu de bromothymol (BBT)]
Le bleu de bromothymol est un indicateur coloré dont le couple \ce{HIn}/\ce{In-} a un $\mathrm{p}K_a = 7,0$. La forme acide \ce{HIn} est jaune, la forme basique \ce{In-} est bleue.
\begin{enumerate}
\item Calcule le rapport $\dfrac{[\ce{In-}]}{[\ce{HIn}]}$ dans une solution de $\mathrm{pH} = 5,0$. Quelle couleur prend le BBT dans cette solution ?
\item Même question pour une solution de $\mathrm{pH} = 9,0$.
\item Pourquoi dit-on que la zone de virage du BBT se situe approximativement entre $\mathrm{pH} = 6,0$ et $\mathrm{pH} = 8,0$ ?
\item On verse quelques gouttes de BBT dans du jus de bissap sucré dont le pH vaut $4,5$. Quelle teinte observe-t-on ?
\end{enumerate}
\end{exercice}

\courssub{Je me prépare au Bac}

\begin{exercice}[Type Bac — L'acide méthanoïque est-il un acide faible ?]
L'acide méthanoïque \ce{HCOOH} est l'acide le plus simple de la famille des acides carboxyliques ; il est responsable de la piqûre des fourmis. Au laboratoire, on prépare une solution aqueuse $S$ d'acide méthanoïque de concentration $C = \SI{2,0e-2}{mol.L^{-1}}$. La mesure du pH à \SI{25}{\celsius} donne $\mathrm{pH} = 2,7$.
\begin{enumerate}
\item Calcule la concentration en ions oxonium $[\ce{H3O+}]$ de la solution $S$.
\item Montre, en comparant $[\ce{H3O+}]$ et $C$, que l'acide méthanoïque est un acide faible.
\item Écris l'équation de la réaction de l'acide méthanoïque avec l'eau et précise les deux couples mis en jeu.
\item Dresser le tableau d'avancement de cette réaction pour un volume $V = \SI{1,0}{L}$ de solution, en notant $x_f$ l'avancement final.
\item Déduis de la question 1 la valeur de $x_f$, puis calcule le taux d'avancement final $\tau$ et le coefficient d'ionisation $\alpha$ de l'acide. Commente.
\item Exprime la constante d'acidité $K_a$ du couple \ce{HCOOH}/\ce{HCOO-} en fonction de $[\ce{H3O+}]$ et $C$, puis calcule $K_a$ et le $\mathrm{p}K_a$ du couple.
\item On dilue dix fois la solution $S$. Sans calcul, indique comment évoluent le pH et le coefficient d'ionisation $\alpha$. Quelle loi traduit cette évolution ?
\end{enumerate}
\emph{Données :} $K_e = 10^{-14}$ à \SI{25}{\celsius} ; $10^{0,7} \approx 5,0$ ; $10^{-2,7} \approx 2,0\times 10^{-3}$.
\end{exercice}

\begin{exercice}[Type Bac — L'ammoniac, une base faible de la vie quotidienne]
L'ammoniac \ce{NH3} est utilisé dans certains produits ménagers décapants. On dissout de l'ammoniac dans l'eau distillée pour obtenir une solution $S_b$ de concentration $C_b = \SI{5,0e-3}{mol.L^{-1}}$. À \SI{25}{\celsius}, le pH-mètre indique $\mathrm{pH} = 10,4$.
\begin{enumerate}
\item Calcule les concentrations $[\ce{H3O+}]$ et $[\ce{HO-}]$ de la solution $S_b$.
\item Si l'ammoniac était une base forte, quelle serait la valeur du pH de la solution ? Conclus que l'ammoniac est une base faible.
\item Écris l'équation de la réaction de l'ammoniac avec l'eau et identifie les deux couples acide/base.
\item Dresser le tableau d'avancement pour un litre de solution et exprime l'avancement final $x_f$ en fonction de $[\ce{HO-}]$.
\item Calcule le coefficient d'ionisation $\alpha$ de l'ammoniac dans cette solution.
\item Donne l'expression de la constante de basicité $K_b$ du couple \ce{NH4+}/\ce{NH3} en fonction de $[\ce{HO-}]$ et $C_b$ (on négligera $[\ce{HO-}]$ devant $C_b$ après vérification). Calcule $K_b$.
\item En déduire la constante d'acidité $K_a$ puis le $\mathrm{p}K_a$ du couple \ce{NH4+}/\ce{NH3}. Compare à la valeur tabulée $\mathrm{p}K_a = 9,2$ en calculant l'écart relatif en pourcentage.
\end{enumerate}
\emph{Données :} $K_e = 10^{-14}$ à \SI{25}{\celsius} ; $10^{-10,4} = 4,0\times 10^{-11}$ ; $10^{-3,6} = 2,5\times 10^{-4}$.
\end{exercice}

\begin{exercice}[Type Bac — TP : détermination expérimentale du pKa de l'acide éthanoïque]
Lors d'une séance de travaux pratiques au lycée, un groupe d'élèves dispose de vinaigre commercial et souhaite vérifier la valeur du $\mathrm{p}K_a$ du couple \ce{CH3COOH}/\ce{CH3COO-}. Après dilution, ils préparent une solution $S$ d'acide éthanoïque de concentration $C = \SI{1,0e-2}{mol.L^{-1}}$. La mesure au pH-mètre, à \SI{25}{\celsius}, donne $\mathrm{pH} = 2,9$.
\begin{enumerate}
\item Écris l'équation de l'équilibre d'ionisation de l'acide éthanoïque dans l'eau.
\item Montre que la valeur du pH confirme que l'acide éthanoïque est un acide faible.
\item Calcule $[\ce{H3O+}]$ et montre que $[\ce{CH3COO-}] = [\ce{H3O+}]$ (justifie à l'aide de l'électroneutralité de la solution, en négligeant les ions \ce{HO-}).
\item Exprime puis calcule le coefficient d'ionisation $\alpha$ de l'acide dans la solution $S$.
\item Établis l'expression littérale $K_a = \dfrac{[\ce{H3O+}]^2}{C - [\ce{H3O+}]}$, puis calcule $K_a$ et le $\mathrm{p}K_a$ expérimental du couple.
\item La valeur tabulée est $\mathrm{p}K_a = 4,8$. Calcule l'écart relatif $\left|\dfrac{\mathrm{p}K_{a,\mathrm{exp}} - \mathrm{p}K_{a,\mathrm{tab}}}{\mathrm{p}K_{a,\mathrm{tab}}}\right|$ en pourcentage et cite deux sources d'erreur possibles lors du TP.
\item Les élèves ajoutent à \SI{20,0}{mL} de la solution $S$ une masse $m = \SI{0,164}{g}$ d'éthanoate de sodium \ce{CH3COONa} (sel totalement dissous, sans variation de volume). Calcule le nouveau pH du mélange obtenu. Quel nom donne-t-on à une telle solution ?
\end{enumerate}
\emph{Données :} $M(\ce{CH3COONa}) = \SI{82,0}{g.mol^{-1}}$ ; $10^{-2,9} = 1,26\times 10^{-3}$ ; $\log 2 \approx 0,30$.
\end{exercice}



\newpage

\coursec{Corrigés des exercices}

\begin{corrige}[Exercice 1 — QCM : réaction acido-basique]
\begin{enumerate}
\item \textbf{Réponse b.} Pour un acide fort, on aurait $\mathrm{pH} = -\log C = -\log(1,0\times 10^{-2}) = 2,0$. Or $\mathrm{pH} = 3,4 > 2,0$, donc $[\ce{H3O+}] = 10^{-3,4} \approx \SI{4,0e-4}{mol.L^{-1}} \ll C$. L'ionisation est partielle : l'acide éthanoïque est un \textbf{acide faible}.
\item \textbf{Réponse b.} $\alpha = \dfrac{x_f}{C\,V}$ est un taux, donc $0 < \alpha < 1$ pour un acide faible. D'après la loi d'Ostwald ($\alpha \approx \sqrt{K_a/C}$), quand $C$ diminue (dilution), $\alpha$ \textbf{augmente}.
\item \textbf{Réponse b.} $\mathrm{pH} = 7,0 < \mathrm{p}K_a = 9,2$ : la \textbf{forme acide} \ce{NH4+} prédomine.
\item \textbf{Réponse b.} $K = \dfrac{K_{a1}}{K_{a2}} = \dfrac{10^{-\mathrm{p}K_{a1}}}{10^{-\mathrm{p}K_{a2}}} = 10^{\mathrm{p}K_{a2}-\mathrm{p}K_{a1}}$.
\end{enumerate}
\textbf{Point de vigilance :} ne jamais inverser l'exposant dans l'expression de $K$ ; vérifie mentalement que la réaction la plus favorable est celle où l'écart $\mathrm{p}K_{a2}-\mathrm{p}K_{a1}$ est grand et positif.
\end{corrige}

\begin{corrige}[Exercice 2 — Vrai ou faux ?]
\begin{enumerate}
\item \textbf{Faux.} Une base faible ne capte qu'une faible fraction des protons de l'eau : sa réaction avec l'eau est \textbf{limitée} et conduit à un équilibre (ex. $\alpha \approx 5\%$ pour l'ammoniac à $\SI{5e-3}{mol.L^{-1}}$).
\item \textbf{Vrai.} Pour tout couple \ce{AH}/\ce{A-} à \SI{25}{\celsius} : $K_a \times K_b = K_e = 10^{-14}$, soit $\mathrm{p}K_a + \mathrm{p}K_b = 14$.
\item \textbf{Faux.} $\mathrm{p}K_a = -\log K_a$ : plus $\mathrm{p}K_a$ est \textbf{grand}, plus $K_a$ est \textbf{petit}, donc plus l'acide est \textbf{faible} (et sa base conjuguée forte).
\item \textbf{Vrai.} \ce{HCO3-} peut capter un proton ($\rightarrow$ \ce{H2CO3}, rôle de base) ou en céder un ($\rightarrow$ \ce{CO3^2-}, rôle d'acide) : c'est un \textbf{ampholyte}.
\item \textbf{Faux.} Par dilution, l'équilibre se déplace dans le sens de l'ionisation (Le Chatelier) : le taux d'avancement final \textbf{augmente} ($\tau \approx \sqrt{K_a/C}$). C'est $[\ce{H3O+}]$ qui diminue, donc le pH qui augmente.
\item \textbf{Vrai.} $\mathrm{pH} = \mathrm{p}K_a + \log\dfrac{[\ce{A-}]}{[\ce{AH}]}$ ; si $\mathrm{pH} = \mathrm{p}K_a$, alors $\log\dfrac{[\ce{A-}]}{[\ce{AH}]} = 0$, donc $[\ce{A-}] = [\ce{AH}]$.
\end{enumerate}
\end{corrige}

\begin{corrige}[Exercice 3 — Texte à trous]
Un acide faible ne s'ionise que de façon \textbf{partielle} dans l'eau : sa dissolution conduit à un \textbf{équilibre} chimique que l'on étudie à l'aide du tableau d'\textbf{avancement}. La force d'un acide se mesure par la constante d'acidité $K_a$ de son couple, ou plus souvent par son \textbf{$\mathrm{p}K_a$}. Une espèce comme \ce{HCO3-}, capable de céder ou de capter un proton, est appelée \textbf{ampholyte}. Pour un couple donné, la forme acide \textbf{prédomine} lorsque $\mathrm{pH} < \mathrm{p}K_a$. À \SI{25}{\celsius}, le produit $K_a \times K_b$ d'un couple est égal à \textbf{$K_e$}, soit $10^{-14}$. Une réaction acido-basique dont la constante $K$ est supérieure à $10^{4}$ est considérée comme \textbf{totale}.
\end{corrige}

\begin{corrige}[Exercice 4 — Définitions et couples]
\begin{enumerate}
\item \textbf{Acide faible :} espèce qui ne cède qu'une faible fraction de ses protons à l'eau ; sa réaction avec l'eau est limitée et conduit à un équilibre. \textbf{Base faible :} espèce qui ne capte qu'une faible fraction des protons de l'eau (équilibre limité). \textbf{Couple acide/base :} ensemble de deux espèces \ce{AH}/\ce{A-} qui s'échangent par transfert d'un proton \ce{H+}. \textbf{Constante d'acidité :} constante d'équilibre de la réaction $\ce{AH + H2O <=> A- + H3O+}$ : $K_a = \dfrac{[\ce{A-}][\ce{H3O+}]}{[\ce{AH}]}$.
\item $\ce{HCOOH <=> HCOO- + H+}$ : acide \ce{HCOOH}, base \ce{HCOO-}.\par
$\ce{NH4+ <=> NH3 + H+}$ : acide \ce{NH4+}, base \ce{NH3}.\par
$\ce{H2O <=> HO- + H+}$ : acide \ce{H2O}, base \ce{HO-}.
\item Relation de Henderson-Hasselbalch :
\[ \mathrm{pH} = \mathrm{p}K_a + \log\frac{[\ce{A-}]}{[\ce{AH}]} \]
\end{enumerate}
\end{corrige}

\begin{corrige}[Exercice 5 — L'ion hydrogénocarbonate, un ampholyte]
\begin{enumerate}
\item Premier équilibre : couples \ce{H2CO3}/\ce{HCO3-} et \ce{H2O}/\ce{HO-}. Second équilibre : couples \ce{HCO3-}/\ce{CO3^2-} et \ce{H3O+}/\ce{H2O}.
\item Dans le premier équilibre, \ce{HCO3-} \textbf{capture} un proton (rôle de \textbf{base}, il devient \ce{H2CO3}) ; dans le second, il \textbf{cède} un proton (rôle d'\textbf{acide}, il devient \ce{CO3^2-}). Pouvant jouer les deux rôles, \ce{HCO3-} est un \textbf{ampholyte}.
\item À $\mathrm{pH} = 7,2$ : $6,4 < 7,2 < 10,3$. Pour le couple \ce{H2CO3}/\ce{HCO3-}, $\mathrm{pH} > \mathrm{p}K_{a1}$ donc \ce{HCO3-} prédomine sur \ce{H2CO3} ; pour le couple \ce{HCO3-}/\ce{CO3^2-}, $\mathrm{pH} < \mathrm{p}K_{a2}$ donc \ce{HCO3-} prédomine sur \ce{CO3^2-}. \textbf{Conclusion :} dans cette eau minérale, l'élément carbone se trouve presque exclusivement sous forme \ce{HCO3-} (c'est ce qui explique le pouvoir tampon des eaux du Mont Cameroun).
\end{enumerate}
\end{corrige}

\begin{corrige}[Exercice 6 — Classement de couples et prévision d'une réaction]
\begin{enumerate}
\item Force acide croissante : $\ce{NH4+} < \ce{CH3COOH} < \ce{HCOOH}$. Justification : plus $\mathrm{p}K_a$ est petit, plus $K_a$ est grand, plus l'acide est fort ($9,2 > 4,8 > 3,8$).
\item La base conjuguée est d'autant plus forte que l'acide est faible. Force basique croissante : $\ce{HCOO-} < \ce{CH3COO-} < \ce{NH3}$.
\item \[ \ce{HCOOH + NH3 <=> HCOO- + NH4+} \]
\item $K = 10^{\mathrm{p}K_{a2}-\mathrm{p}K_{a1}} = 10^{9,2-3,8} = 10^{5,4} \approx \SI{2,5e5}{}$. Comme $K \gg 10^{4}$, la réaction est \textbf{totale} : elle peut servir de réaction de dosage.
\end{enumerate}
\end{corrige}

\begin{corrige}[Exercice 7 — Mélange acide éthanoïque / éthanoate de sodium]
\begin{enumerate}
\item $\mathrm{pH} = \mathrm{p}K_a + \log\dfrac{[\ce{CH3COO-}]}{[\ce{CH3COOH}]}$.
\item $\log\dfrac{[\ce{CH3COO-}]}{[\ce{CH3COOH}]} = \mathrm{pH} - \mathrm{p}K_a = 5,8 - 4,8 = 1,0$, donc $\dfrac{[\ce{CH3COO-}]}{[\ce{CH3COOH}]} = 10^{1} = 10$.
\item $[\ce{CH3COO-}] = 10\,[\ce{CH3COOH}]$ : la \textbf{forme basique} \ce{CH3COO-} prédomine (cohérent avec $\mathrm{pH} > \mathrm{p}K_a$).
\item Diagramme de prédominance :
\begin{popfigure}
\begin{tikzpicture}[scale=0.85, every node/.style={font=\small}]
\draw[->, thick] (0,0) -- (12,0) node[right] {pH};
\foreach \x in {0,2,4,6,8,10} \draw (\x,0.1) -- (\x,-0.1) node[below] {\x};
\fill[popblue!20] (0,0.5) rectangle (4.8,1.5) node[midway, align=center] {\ce{CH3COOH}\\prédomine};
\fill[poporange!20] (4.8,0.5) rectangle (12,1.5) node[midway, align=center] {\ce{CH3COO-}\\prédomine};
\draw[dashed, thick, popdark] (4.8,0) -- (4.8,1.7) node[above] {$\mathrm{p}K_a = 4,8$};
\draw[popink, thick] (5.8,-0.35) -- (5.8,1.7);
\fill[popink] (5.8,1.0) circle (2pt) node[right, popink] {solution étudiée (pH 5,8)};
\end{tikzpicture}
\legende{Diagramme de prédominance du couple \ce{CH3COOH}/\ce{CH3COO-}.}
\end{popfigure}
La solution (pH $= 5,8$) se situe dans le domaine de \ce{CH3COO-}.
\end{enumerate}
\end{corrige}

\begin{corrige}[Exercice 8 — Le bleu de bromothymol (BBT)]
\begin{enumerate}
\item $\log\dfrac{[\ce{In-}]}{[\ce{HIn}]} = \mathrm{pH} - \mathrm{p}K_a = 5,0 - 7,0 = -2,0 \Rightarrow \dfrac{[\ce{In-}]}{[\ce{HIn}]} = 10^{-2} = 0,01$. La forme acide \ce{HIn} est cent fois majoritaire : teinte \textbf{jaune}.
\item $\log\dfrac{[\ce{In-}]}{[\ce{HIn}]} = 9,0 - 7,0 = 2,0 \Rightarrow \dfrac{[\ce{In-}]}{[\ce{HIn}]} = 10^{2} = 100$. La forme basique domine : teinte \textbf{bleue}.
\item L'œil perçoit la teinte d'une forme lorsqu'elle est au moins dix fois plus concentrée que l'autre : $\dfrac{[\ce{In-}]}{[\ce{HIn}]} \in [0,1\,;\,10] \Leftrightarrow \mathrm{pH} \in [\mathrm{p}K_a - 1\,;\,\mathrm{p}K_a + 1] = [6,0\,;\,8,0]$. C'est la \textbf{zone de virage} (teinte verte, mélange).
\item $\mathrm{pH} = 4,5 < 6,0$ : $\dfrac{[\ce{In-}]}{[\ce{HIn}]} = 10^{4,5-7,0} = 10^{-2,5} \approx 3\times 10^{-3}$. \ce{HIn} ultra-majoritaire : le jus de bissap additionné de BBT paraît \textbf{jaune}.
\end{enumerate}
\end{corrige}

\begin{corrige}[Exercice 9 — Type Bac : l'acide méthanoïque est-il un acide faible ?]
\begin{enumerate}
\item $[\ce{H3O+}] = 10^{-\mathrm{pH}} = 10^{-2,7} \approx \SI{2,0e-3}{mol.L^{-1}}$.
\item $[\ce{H3O+}] = \SI{2,0e-3}{mol.L^{-1}} < C = \SI{2,0e-2}{mol.L^{-1}}$ (dix fois moins). L'acide n'est pas totalement ionisé : c'est un \textbf{acide faible}.
\item $\ce{HCOOH_{(aq)} + H2O_{(l)} <=> HCOO-_{(aq)} + H3O+_{(aq)}}$ ; couples : \ce{HCOOH}/\ce{HCOO-} et \ce{H3O+}/\ce{H2O}.
\item Tableau d'avancement pour $V = \SI{1,0}{L}$ ($n_i = C\,V = \SI{2,0e-2}{mol}$) :
\begin{center}
\begin{tabularx}{\linewidth}{|l|X|X|X|X|}\hline
 & \ce{HCOOH} & \ce{H2O} & \ce{HCOO-} & \ce{H3O+} \\\hline
État initial & $n_i = \SI{2,0e-2}{mol}$ & excès & $0$ & $0$ \\\hline
État final & $n_i - x_f$ & excès & $x_f$ & $x_f$ \\\hline
\end{tabularx}
\end{center}
\item $x_f = [\ce{H3O+}]\times V = \SI{2,0e-3}{mol}$. $\tau = \alpha = \dfrac{x_f}{n_i} = \dfrac{2,0\times 10^{-3}}{2,0\times 10^{-2}} = 0,10 = \textbf{10\%}$. Seule une molécule sur dix est ionisée : l'ionisation est bien partielle.
\item $K_a = \dfrac{[\ce{HCOO-}][\ce{H3O+}]}{[\ce{HCOOH}]} = \dfrac{[\ce{H3O+}]^2}{C - [\ce{H3O+}]} = \dfrac{(2,0\times 10^{-3})^2}{2,0\times 10^{-2} - 2,0\times 10^{-3}} = \dfrac{4,0\times 10^{-6}}{1,8\times 10^{-2}} \approx \SI{2,2e-4}{}$. $\mathrm{p}K_a = -\log(2,2\times 10^{-4}) = 4 - \log 2,2 \approx 4 - 0,34 = \textbf{3,66}$, proche de la valeur tabulée $3,8$.
\item Dilution dix fois : $C$ diminue ; d'après la \textbf{loi d'Ostwald} ($\alpha \approx \sqrt{K_a/C}$), $\alpha$ \textbf{augmente} ; $[\ce{H3O+}] \approx \sqrt{K_a C}$ diminue, donc le \textbf{pH augmente} (mais de moins d'une unité).
\end{enumerate}
\textbf{Barème indicatif (sur 10) :} Q1 : 1 ; Q2 : 1 ; Q3 : 1,5 ; Q4 : 1,5 ; Q5 : 1,5 ; Q6 : 2 ; Q7 : 1,5.\par
\textbf{Points de vigilance :} flèche \ce{<=>} obligatoire (équilibre) ; ne pas oublier de soustraire $[\ce{H3O+}]$ au dénominateur de $K_a$ ici, car $\alpha = 10\% > 5\%$ (l'approximation $K_a \approx C\alpha^2$ n'est pas licite) ; distinguer $x_f$ (mol) et $\tau$ (sans unité).
\end{corrige}

\begin{corrige}[Exercice 10 — Type Bac : l'ammoniac, une base faible de la vie quotidienne]
\begin{enumerate}
\item $[\ce{H3O+}] = 10^{-10,4} = \SI{4,0e-11}{mol.L^{-1}}$ ; $[\ce{HO-}] = \dfrac{K_e}{[\ce{H3O+}]} = \dfrac{10^{-14}}{4,0\times 10^{-11}} = \SI{2,5e-4}{mol.L^{-1}}$.
\item Si \ce{NH3} était une base forte : $[\ce{HO-}] = C_b = \SI{5,0e-3}{mol.L^{-1}}$, $\mathrm{pOH} = 2,3$, $\mathrm{pH} = 11,7$. Or $\mathrm{pH} = 10,4 < 11,7$ : la protonation est partielle, \ce{NH3} est une \textbf{base faible}.
\item $\ce{NH3_{(aq)} + H2O_{(l)} <=> NH4+_{(aq)} + HO-_{(aq)}}$ ; couples : \ce{NH4+}/\ce{NH3} et \ce{H2O}/\ce{HO-}.
\item Pour $V = \SI{1,0}{L}$ :
\begin{center}
\begin{tabularx}{\linewidth}{|l|X|X|X|X|}\hline
 & \ce{NH3} & \ce{H2O} & \ce{NH4+} & \ce{HO-} \\\hline
État initial & $n_i = \SI{5,0e-3}{mol}$ & excès & $0$ & $0$ \\\hline
État final & $n_i - x_f$ & excès & $x_f$ & $x_f$ \\\hline
\end{tabularx}
\end{center}
avec $x_f = [\ce{HO-}]\times V = \SI{2,5e-4}{mol}$.
\item $\alpha = \dfrac{x_f}{n_i} = \dfrac{2,5\times 10^{-4}}{5,0\times 10^{-3}} = 0,050 = \textbf{5,0\%}$.
\item $[\ce{HO-}] = 0,05\,C_b \ll C_b$ : approximation justifiée (limite de validité usuelle $\alpha < 5\%$). $K_b = \dfrac{[\ce{NH4+}][\ce{HO-}]}{[\ce{NH3}]} \approx \dfrac{[\ce{HO-}]^2}{C_b} = \dfrac{(2,5\times 10^{-4})^2}{5,0\times 10^{-3}} = \dfrac{6,25\times 10^{-8}}{5,0\times 10^{-3}} = \SI{1,25e-5}{}$.
\item $K_a = \dfrac{K_e}{K_b} = \dfrac{10^{-14}}{1,25\times 10^{-5}} = \SI{8,0e-10}{}$ ; $\mathrm{p}K_a = 10 - \log 8,0 = 10 - 0,90 = \textbf{9,10}$. Écart relatif : $\dfrac{|9,10 - 9,2|}{9,2}\times 100 \approx \textbf{1,1\%}$ : excellent accord (écart dû aux arrondis et à la température).
\end{enumerate}
\textbf{Barème indicatif (sur 10) :} Q1 : 1,5 ; Q2 : 1,5 ; Q3 : 1,5 ; Q4 : 1,5 ; Q5 : 1 ; Q6 : 1,5 ; Q7 : 1,5.\par
\textbf{Points de vigilance :} pour une base, raisonner sur $[\ce{HO-}]$ et non $[\ce{H3O+}]$ ; toujours vérifier l'approximation avant de négliger $[\ce{HO-}]$ devant $C_b$ ; bien utiliser $K_a \times K_b = K_e$ pour le \emph{même} couple.
\end{corrige}

\begin{corrige}[Exercice 11 — Type Bac : TP, détermination expérimentale du pKa de l'acide éthanoïque]
\begin{enumerate}
\item $\ce{CH3COOH_{(aq)} + H2O_{(l)} <=> CH3COO-_{(aq)} + H3O+_{(aq)}}$.
\item $-\log C = 2,0$ et $\mathrm{pH} = 2,9 > 2,0$, donc $[\ce{H3O+}] < C$ : ionisation partielle, l'acide éthanoïque est \textbf{faible}.
\item $[\ce{H3O+}] = 10^{-2,9} = \SI{1,26e-3}{mol.L^{-1}}$. Électroneutralité : $[\ce{H3O+}] = [\ce{CH3COO-}] + [\ce{HO-}]$ avec $[\ce{HO-}] = \dfrac{10^{-14}}{1,26\times 10^{-3}} \approx \SI{7,9e-12}{mol.L^{-1}}$, négligeable. Donc $[\ce{CH3COO-}] \approx [\ce{H3O+}] = \SI{1,26e-3}{mol.L^{-1}}$.
\item $\alpha = \dfrac{[\ce{CH3COO-}]}{C} = \dfrac{1,26\times 10^{-3}}{1,0\times 10^{-2}} = 0,126 = \textbf{12,6\%}$.
\item $K_a = \dfrac{[\ce{CH3COO-}][\ce{H3O+}]}{[\ce{CH3COOH}]} = \dfrac{[\ce{H3O+}]^2}{C - [\ce{H3O+}]} = \dfrac{(1,26\times 10^{-3})^2}{1,0\times 10^{-2} - 1,26\times 10^{-3}} = \dfrac{1,59\times 10^{-6}}{8,74\times 10^{-3}} \approx \SI{1,8e-5}{}$. $\mathrm{p}K_{a,\mathrm{exp}} = 5 - \log 1,8 \approx \textbf{4,74}$.
\item Écart relatif : $\dfrac{|4,74 - 4,8|}{4,8}\times 100 \approx \textbf{1,3\%}$. Sources d'erreur : température différente de \SI{25}{\celsius} ($K_a$ dépend de $T$) ; pH-mètre mal étalonné (solutions tampons périmées) ; présence d'autres acides dans le vinaigre.
\item $n(\ce{CH3COONa}) = \dfrac{m}{M} = \dfrac{0,164}{82,0} = \SI{2,0e-3}{mol}$, soit $[\ce{CH3COO-}]_{\text{apporté}} = \dfrac{2,0\times 10^{-3}}{2,0\times 10^{-2}} = \SI{0,10}{mol.L^{-1}}$. L'acide apporte $[\ce{CH3COOH}] \approx C = \SI{1,0e-2}{mol.L^{-1}}$ (l'ion commun \ce{CH3COO-} refoule l'ionisation, déjà faible). Alors :
\[ \mathrm{pH} = \mathrm{p}K_a + \log\frac{[\ce{CH3COO-}]}{[\ce{CH3COOH}]} = 4,8 + \log\frac{0,10}{0,010} = 4,8 + 1,0 = \textbf{5,8} \]
Ce mélange acide faible / base conjuguée en quantités comparables est une \textbf{solution tampon} : son pH varie très peu par ajout modéré d'acide, de base ou par dilution.
\end{enumerate}
\textbf{Barème indicatif (sur 12) :} Q1 : 1 ; Q2 : 1,5 ; Q3 : 2 ; Q4 : 1,5 ; Q5 : 2 ; Q6 : 2 ; Q7 : 2.\par
\textbf{Points de vigilance :} justifier $[\ce{CH3COO-}] = [\ce{H3O+}]$ par l'électroneutralité (et non l'affirmer sans preuve) ; dans la question 7, convertir la masse en quantité de matière \emph{avant} de calculer la concentration, et vérifier que le volume ne change pas ; conclure en nommant explicitement la solution tampon.
\end{corrige}

\newpage

\coursec{Fiche bilan}

\coursec{Acides faibles, bases faibles et couples acide/base}

\begin{popfigure}
\begin{tikzpicture}[
  mindmap,
  grow cyclic,
  every node/.style={concept, font=\small\bfseries},
  root concept/.append style={concept color=popblue, font=\large\bfseries, text=white, minimum size=3.4cm},
  level 1 concept/.append style={sibling angle=60, level distance=4.8cm, minimum size=2.6cm, font=\footnotesize\bfseries, text=white},
  level 2 concept/.append style={sibling angle=45, level distance=3.0cm, minimum size=2.2cm, font=\scriptsize, text=popdark},
]
\node[root concept]{Acides faibles, bases faibles et couples acide/base}
  child[concept color=poporange]{ node {Acide faible : \ce{CH3COOH}}
    child[concept color=poporangeL]{ node {Ionisation partielle : $\alpha < 1$}}
    child[concept color=poporangeL]{ node {$\mathrm{pH} > -\log C$}}
  }
  child[concept color=popgreen]{ node {Base faible : \ce{CH3COO-}, \ce{NH3}}
    child[concept color=popgreenL]{ node {$[\ce{HO-}] < C$}}
    child[concept color=popgreenL]{ node {$\mathrm{pH} < 14 + \log C$}}
  }
  child[concept color=poppurple]{ node {Couples acide/base}
    child[concept color=poppurpleL]{ node {\ce{AH}/\ce{A-} ; couples de l'eau}}
    child[concept color=poppurpleL]{ node {Ampholytes : \ce{H2O}, \ce{HCO3-}}}
  }
  child[concept color=popgold]{ node {$K_a$, $K_b$, $\mathrm{p}K_a$}
    child[concept color=popgoldL]{ node {$K_a \cdot K_b = K_e$}}
    child[concept color=popgoldL]{ node {$K = 10^{\mathrm{p}K_{a2}-\mathrm{p}K_{a1}}$}}
  }
  child[concept color=poppink]{ node {Relation pH--$\mathrm{p}K_a$}
    child[concept color=poppinkL]{ node {$\mathrm{pH} = \mathrm{p}K_a + \log\dfrac{[\ce{A-}]}{[\ce{AH}]}$}}
    child[concept color=poppinkL]{ node {Domaines de prédominance}}
  }
  child[concept color=popblue!70]{ node {Force des acides et bases}
    child[concept color=popblueL]{ node {$K_a$ grand $\Rightarrow$ acide fort}}
    child[concept color=popblueL]{ node {Indicateurs colorés}}
  };
\end{tikzpicture}
\legende{Carte mentale de la leçon : les six idées-forces à maîtriser pour le Bac.}
\end{popfigure}

\begin{definition}[Définitions fondamentales]
\begin{itemize}
  \item \cle{Acide faible} : acide dont la réaction avec l'eau est \textbf{partielle} et réversible (équilibre chimique). Pour une solution de concentration molaire $C$ : la concentration en ions \ce{H3O+} est \textbf{strictement inférieure} à $C$, d'où $\mathrm{pH} > -\log C$.
  \item \cle{Base faible} : base dont la réaction avec l'eau est partielle et réversible. Pour une concentration $C$ : $[\ce{HO-}] < C$, d'où $\mathrm{pH} < 14 + \log C$ (à \SI{25}{\celsius}).
  \item \cle{Couple acide/base} \ce{AH}/\ce{A-} : ensemble de deux espèces chimiques qui s'interconvertissent par transfert d'un proton \ce{H+} : \ce{AH <=> A- + H+}. \ce{AH} est l'acide, \ce{A-} sa base conjuguée.
  \item \cle{Ampholyte} (espèce amphotère) : espèce capable d'agir comme un acide (donneur de \ce{H+}) vis-à-vis d'une base plus forte, et comme une base (accepteur de \ce{H+}) vis-à-vis d'un acide plus fort. Exemples : \ce{H2O}, \ce{HCO3-}, \ce{HSO4-}, \ce{H2PO4-}.
  \item \cle{Coefficient d'ionisation} $\alpha$ : quotient de la quantité d'entités ionisées sur la quantité d'entités introduites. Pour un acide faible \ce{AH} : $\alpha = \dfrac{[\ce{H3O+}]_{\text{eq}}}{C} = \dfrac{10^{-\mathrm{pH}}}{C}$. Il est sans unité, $0 < \alpha < 1$.
\end{itemize}
\end{definition}

\begin{propriete}[Équations-types et formulaires à connaître par cœur]
\textbf{Équations d'ionisation dans l'eau (double flèche \ce{<=>} obligatoire)}
\begin{itemize}
  \item Acide éthanoïque : \ce{CH3COOH + H2O <=> CH3COO- + H3O+}
  \item Ion éthanoate (base conjuguée) : \ce{CH3COO- + H2O <=> CH3COOH + HO-}
  \item Ammoniac : \ce{NH3 + H2O <=> NH4+ + HO-}
  \item Couples de l'eau (convention Bac) : \ce{H3O+}/\ce{H2O} ($\mathrm{p}K_a = 0$) et \ce{H2O}/\ce{HO-} ($\mathrm{p}K_a = 14$)
  \item Autoprotolyse de l'eau : \ce{2 H2O <=> H3O+ + HO-} \quad ($K_e = 10^{-14}$ à \SI{25}{\celsius})
  \item Réaction acido-basique générale : \ce{A1H + B2 <=> A1- + B2H}
\end{itemize}

\textbf{Constantes d'équilibre}
\begin{center}
\begin{tabularx}{\linewidth}{|l|X|}
\hline
\rowcolor{popblueL}
\textbf{Grandeur} & \textbf{Expression et relation} \\
\hline
\rowcolor{popcream}
Constante d'acidité & $K_a = \dfrac{[\ce{A-}]\,[\ce{H3O+}]}{[\ce{AH}]}$ \quad ; \quad $\mathrm{p}K_a = -\log K_a$ \\
\hline
Constante de basicité & $K_b = \dfrac{[\ce{AH}]\,[\ce{HO-}]}{[\ce{A-}]}$ \quad ; \quad $\mathrm{p}K_b = -\log K_b$ \\
\hline
\rowcolor{popcream}
Relation fondamentale & $K_a \times K_b = K_e = 10^{-14}$ \quad $\Rightarrow$ \quad $\mathrm{p}K_a + \mathrm{p}K_b = 14$ (à \SI{25}{\celsius}) \\
\hline
Relation pH--$\mathrm{p}K_a$ (Henderson--Hasselbalch) & $\mathrm{pH} = \mathrm{p}K_a + \log\dfrac{[\ce{A-}]}{[\ce{AH}]}$ \\
\hline
\rowcolor{popcream}
Constante de réaction & \ce{A1H + B2 <=> A1- + B2H} : $K = \dfrac{K_{a1}}{K_{a2}} = 10^{\mathrm{p}K_{a2}-\mathrm{p}K_{a1}}$ \\
\hline
Coefficient d'ionisation & $\alpha = \dfrac{[\ce{H3O+}]}{C} = \dfrac{10^{-\mathrm{pH}}}{C}$ \quad (acide faible monoprotique) \\
\hline
\rowcolor{popcream}
$K_a$ en fonction de $\alpha$ & $K_a = \dfrac{C\,\alpha^{2}}{1-\alpha}$ \quad ; \quad si $\alpha \ll 1$ (souvent $\alpha < 5\%$) : $K_a \approx C\alpha^{2}$ \\
\hline
\end{tabularx}
\end{center}
\end{propriete}

\begin{methode}[Démarches types pour le Bac]
\textbf{1. Prouver qu'un acide (ou une base) est faible}
\begin{enumerate}
  \item Calculer la concentration théorique en \ce{H3O+} si l'ionisation était totale : $[\ce{H3O+}]_{\text{th}} = C$.
  \item Calculer le pH théorique : $\mathrm{pH}_{\text{th}} = -\log C$.
  \item Comparer au pH mesuré (ou donné) : si $\mathrm{pH}_{\text{mes}} > \mathrm{pH}_{\text{th}}$ (ou $\alpha < 1$), l'ionisation est partielle $\Rightarrow$ \textbf{acide faible}.
  \item Pour une base : si $\mathrm{pH}_{\text{mes}} < 14 + \log C$ (ou $[\ce{HO-}] < C$), c'est une \textbf{base faible}.
\end{enumerate}

\textbf{2. Déterminer $K_a$ (ou $K_b$) expérimentalement}
\begin{enumerate}
  \item Mesurer le pH d'une solution de concentration $C$ connue.
  \item Calculer $[\ce{H3O+}] = 10^{-\mathrm{pH}}$.
  \item À l'équilibre : $[\ce{H3O+}] = [\ce{A-}] = 10^{-\mathrm{pH}}$ et $[\ce{AH}] = C - 10^{-\mathrm{pH}}$.
  \item Appliquer $K_a = \dfrac{(10^{-\mathrm{pH}})^2}{C - 10^{-\mathrm{pH}}}$ (ou l'approximation $K_a \approx \dfrac{10^{-2\mathrm{pH}}}{C}$ si $\alpha < 5\%$).
\end{enumerate}

\textbf{3. Classer des couples acide/base et comparer les forces}
\begin{enumerate}
  \item Ranger les couples par $\mathrm{p}K_a$ croissant (ou $K_a$ décroissant).
  \item \textbf{Règle} : Plus $\mathrm{p}K_a$ est \textbf{petit} (plus $K_a$ est \textbf{grand}), plus l'acide \ce{AH} est \textbf{fort} et plus sa base conjuguée \ce{A-} est \textbf{faible}.
  \item L'acide le plus fort réagit avec la base la plus forte (règle du $\gamma$).
\end{enumerate}

\textbf{4. Prévoir le sens et l'avancement d'une réaction acido-basique}
\begin{enumerate}
  \item Identifier les deux couples : \ce{A1H}/\ce{A1-} et \ce{B2H}/\ce{B2} (avec \ce{B2H} acide, \ce{B2} base).
  \item Calculer $K = 10^{\mathrm{p}K_{a}(\ce{B2H}) - \mathrm{p}K_{a}(\ce{A1H})}$.
  \item Interpréter :
  \begin{itemize}
    \item $K > 10^4$ : réaction \textbf{quasi totale} (flèche \ce{->}).
    \item $10^{-4} < K < 10^4$ : réaction \textbf{limitée} (équilibre, flèche \ce{<=>}).
    \item $K < 10^{-4}$ : réaction \textbf{quasi nulle} dans le sens écrit.
  \end{itemize}
\end{enumerate}

\textbf{5. Choisir un indicateur coloré pour un dosage}
\begin{enumerate}
  \item Connaître le pH à l'équivalence ($\mathrm{pH}_{\text{eq}}$).
  \item Choisir un indicateur dont la zone de virage ($\mathrm{p}K_i \pm 1$) \textbf{englobe} $\mathrm{pH}_{\text{eq}}$.
\end{enumerate}
\end{methode}

\begin{savaistu}[Chimie et vie quotidienne au Cameroun]
\begin{itemize}
  \item \textbf{Vin de palme et bili-bili} : La fermentation alcoolique produit de l'éthanol, puis les bactéries acétiques (acétobacters) oxydent l'éthanol en \textbf{acide éthanoïque} (vinaigre). C'est un acide faible ($\mathrm{p}K_a = 4,8$) : le vinaigre de palme a un pH $\approx 2,5$--$3,0$ pour $C \approx 0,5$--$1\,\text{mol/L}$.
  \item \textbf{Savonneries artisanales (savon noir, savon de Marseille local)} : La saponification de l'huile de palme (triglycérides) par la soude \ce{NaOH} (base forte) libère du glycérol et des \textbf{carboxylates} (bases faibles, conjuguées des acides gras). Le savon en solution est basique (pH 9--10) car l'ion carboxylate hydrolyse l'eau : \ce{RCOO- + H2O <=> RCOOH + HO-}.
  \item \textbf{SONARA (Raffinerie de Limbé)} : Traitement des eaux acides (contenant \ce{H2S}, \ce{CO2}) par neutralisation avec de la chaux \ce{Ca(OH)2} ou de l'ammoniac \ce{NH3} (base faible, $\mathrm{p}K_a(\ce{NH4+}) = 9,2$).
  \item \textbf{Pharmacopée et jus de bissap} : L'hibiscus (bissap) contient des anthocyanes (indicateurs naturels rouge/bleu selon pH) et des acides organiques faibles (citrique, malique). Le pH $\approx 2,5$--$3,0$ explique l'acidité gustative et la conservation.
  \item \textbf{Mont Cameroun} : Émanations volcaniques de \ce{CO2}, \ce{SO2}, \ce{H2S}. \ce{CO2} dissous forme l'acide carbonique \ce{H2CO3} (couple \ce{H2CO3}/\ce{HCO3-}, $\mathrm{p}K_{a1} = 6,3$), acidifiant les eaux de ruissellement.
\end{itemize}
\end{savaistu}

\begin{aretenir}[Checklist avant le Bac -- Module 2 : Acides et bases]
\begin{itemize}
  \item[$\square$] Je sais définir un acide faible et une base faible et citer des exemples classiques (\ce{CH3COOH}, \ce{NH3}, \ce{CH3COO-}, \ce{HCO3-}).
  \item[$\square$] Je sais montrer, à partir du pH mesuré et de la concentration $C$, qu'un acide (ou une base) n'est pas totalement ionisé(e) ($\mathrm{pH} > -\log C$ ou $\alpha < 1$).
  \item[$\square$] Je sais écrire l'équation d'ionisation d'un acide ou d'une base faible dans l'eau avec la double flèche d'équilibre \ce{<=>} et l'eau comme réactif.
  \item[$\square$] Je sais identifier les deux couples acide/base mis en jeu dans un équilibre acido-basique (ex: \ce{CH3COOH}/\ce{CH3COO-} et \ce{H3O+}/\ce{H2O}).
  \item[$\square$] Je connais les couples de l'eau (\ce{H3O+}/\ce{H2O} $\mathrm{p}K_a=0$ ; \ce{H2O}/\ce{HO-} $\mathrm{p}K_a=14$) et je sais reconnaître un ampholyte (\ce{H2O}, \ce{HCO3-}, \ce{HSO4-}, \ce{H2PO4-}).
  \item[$\square$] Je sais écrire l'expression de $K_a$ et de $K_b$ (produits sur réactifs, espèces pures exclues) et utiliser la relation $K_a \cdot K_b = K_e = 10^{-14}$ (à \SI{25}{\celsius}).
  \item[$\square$] Je sais utiliser la relation $\mathrm{pH} = \mathrm{p}K_a + \log\dfrac{[\ce{A-}]}{[\ce{AH}]}$ pour calculer un pH, un rapport de concentrations ou un $\mathrm{p}K_a$.
  \item[$\square$] Je sais tracer et interpréter les domaines de prédominance d'un couple \ce{AH}/\ce{A-} : \ce{AH} prédomine si $\mathrm{pH} < \mathrm{p}K_a - 1$, \ce{A-} prédomine si $\mathrm{pH} > \mathrm{p}K_a + 1$, $[\ce{AH}] = [\ce{A-}]$ si $\mathrm{pH} = \mathrm{p}K_a$.
  \item[$\square$] Je sais calculer le coefficient d'ionisation $\alpha = 10^{-\mathrm{pH}}/C$ et déterminer $K_a$ à partir d'une mesure de pH (avec ou sans approximation $\alpha \ll 1$).
  \item[$\square$] Je sais que la dilution augmente le coefficient d'ionisation $\alpha$ (loi d'Ostwald) mais diminue la concentration en \ce{H3O+} (le pH augmente).
  \item[$\square$] Je sais classer des couples selon $K_a$ ou $\mathrm{p}K_a$ : plus $K_a$ grand / $\mathrm{p}K_a$ petit $\Rightarrow$ acide plus fort, base conjuguée plus faible.
  \item[$\square$] Je sais calculer la constante de réaction $K = 10^{\mathrm{p}K_{a2}-\mathrm{p}K_{a1}}$ pour \ce{A1H + B2 <=> A1- + B2H} et conclure sur le sens et l'avancement (quasi total si $K > 10^4$).
  \item[$\square$] Je sais choisir un indicateur coloré adapté à un dosage : zone de virage ($\mathrm{p}K_i \pm 1$) contenant le pH à l'équivalence.
\end{itemize}
\end{aretenir}

\findecours

\end{document}
